% 附录: 习题训练 (Aufgabentraining)
% 对应源文档 appendix_src.txt (第 1–6822 行).
% 注: 第 1 章解答为完整翻译; 引用第 2–7 章定理处所用标签 (def:2.1.1, thm:2.1.10,
% def:5.2.10, thm:5.2.12 等) 在其对应章节译出前会显示为 "??", 属预期现象.

\section*{总论}

书稿中各个章节的习题提供了将相应呈现的理论运用于具体例子的机会, 从而得以更深入
地加工与理解.有时需要的只是纯粹的例行验证, 但有时也需要稍具创造性的思考.本附录
旨在说明, 如何对习题的解答过程进行有意义的组织安排, 从而使解题更加有效.

在每一章的各个小节中, 我都选取了一些习题, 其解答将在下文以范例的方式加以讨论.
这些习题都带有形如 (AT xyz) 的标注, 表示相应解答的讨论见于习题训练中的第 xyz 页.
自然, 解答一道习题应以最终给出一个尽可能简单的,形式流畅的解答为目标.但至少同样
重要的是事先必须进行的那些思考, 它们导向真正的解题路径.我将这些思考概括在
"策略"这一关键词之下.一般而言, 策略性思考在完成的解答中已不再显现, 但由于其
核心地位, 在随后的讨论中它们往往占据比解答本身更大的篇幅.事实上, 若要以系统的
方式独立地构思一个解答, 或者仅仅是想复现或更好地理解一个现成的解答, 策略性思考
都是不可或缺的.下面将更精确地描述解答一道习题时所需要经历的各个步骤.

\noindent\textbf{起点:} 在一开始, 应从习题文本中提炼出初始情形以及可资利用的前提
条件.此时, 有必要重新回顾所用到的概念构设, 连同在所呈现理论的框架内处理过的相应
性质与结论.

\noindent\textbf{目标:} 在第二步, 应当将所要证明的结论(如有必要)清晰地加以梳理.
尤其应当明确, 具体有哪些论断需要证明.与第一步一样, 在此也有必要再次把所用到的
概念构设以及相应的性质与结论展现在眼前.

\noindent\textbf{策略:} 此时开始规划从起点到目标的实际路线.为了找到一条可行的
路径, 可以首先更仔细地审视初始情形, 并考虑是否存在与所处理的理论相关联的,指向
目标方向的等价或类似的起始条件.另一方面, 也可以进一步分析目标, 提出如下类型的
问题: 目标是否存在等价的描述? 是否有明显的结论可从目标推出? 换言之, 目标是否
存在可能更容易达成的等价版本或弱化形式? 或者是否存在一个适用于完全归纳法的约简
步骤? 简言之, 任何有助于更精确地勘测起点与目标之间"地形"的思考都是有价值的.
在此还尤其应当估量: 目标是否适合直接证明, 抑或更适合通过反证法间接地加以处理.

此外, 考察初始情形中的具体例子通常也极为有用, 无论是几何性质的例子, 还是通过
进一步限制起始前提得到的例子.在此过程中往往能获得有价值的线索, 得知在放宽条件下
如何达成目标, 并由此或许学会在一般情形下如何行事.遗憾的是, 并不存在一个普遍
有效的方法, 能保证将我们从起点带到目标.诚然, 基于上述思考, 通常可以把两点之间
的未知区域大为缩小, 但常常仍剩下一个看似不可逾越的障碍.此时就需要一些创造力,
以发现某个隐藏的突破口.一旦可以说"看见"了一条通向目标的可行路径, 便可进入
下一步, 尝试表述一个完整而数学上精确的解答.

\noindent\textbf{解题路径:} 在这一步, 应当把通过策略性思考所规划的通向目标的路径
以精确的数学形式加以实现.所有那些仅仅有助于发现解题路径,但对实际的求解过程而言
无关紧要的细节在此均被略去.另一方面, 在策略考察的框架内所采用的种种简化, 在这一
步不再被允许.当然, 撰写一份流畅的数学解答需要一定的经验, 而这种经验会随着时间
的推移而逐渐获得.这种经验尤其涉及选择清晰而合用的记号, 以及精心设计论证与所用
结论的先后顺序.这一切都应当在尽可能大的清晰性与简单性的观点下进行.

不过, 倘若在实现预期的解题路径时出乎意料地遭遇不可逾越的困难, 那就别无他法, 只能
回到上一步, 并带着新获得的经验重新进行策略考察.

\noindent\textbf{补充:} 在此要点下, 我们给出关于相应习题问题的若干附加的可选信息,
倘若这对一般读者具有意义的话.就此而言, 这一要点并不属于解答一道习题时所应完成的
常规步骤.

\section*{解法讨论}

下文所讨论的习题解答, 连同相应的策略性思考, 只能被视为范例性的.一般说来, 一个
习题问题存在多种不同,各有道理且地位对等的处理方法.对于导向真正解答的策略性思考,
这一点尤为适用.因此, 我们绝不声称后续的解答讨论是必须遵循的,不可用其他思考替代
的解题方式.

\section*{对应 1.1-1.6: 向量空间}

% 源文档第 118–1392 行.

\subsection*{习题 1.1/1}

\noindent\textbf{起点:} 给定集合 $X$ 的子集 $A, B, C$.

\noindent\textbf{目标:} 需证明下列恒等式:
\[
\text{(i)}\ A \cap (B \cup C) = (A \cap B) \cup (A \cap C),
\]
\[
\text{(ii)}\ A \cup (B \cap C) = (A \cup B) \cap (A \cup C),
\]
\[
\text{(iii)}\ A - (B \cup C) = (A - B) \cap (A - C),
\]
\[
\text{(iv)}\ A - (B \cap C) = (A - B) \cup (A - C)\text{.}
\]

\noindent\textbf{策略:} 我们尚无可用来描述集合的并,交与差之间相互作用的结果.
因此别无他法, 只能依据这些运算的定义, 逐元素地比较各等式左右两侧.为了由此断定
两侧相等, 最清晰的做法是证明: 左侧的每个元素也都包含于右侧, 反之亦然.

\noindent\textbf{解题路径:} 我们从 (i) 开始, 先证包含关系 $\subset$.设
$x \in A \cap (B \cup C)$.则得 $x \in A$ 且 $x \in B \cup C$, 后者意味着
$x \in B$ 或 $x \in C$.对于 $x \in B$, 得 $x \in A \cap B$, 从而
$x \in (A \cap B) \cup (A \cap C)$.相应地, $x \in C$ 同样导致
$x \in (A \cap B) \cup (A \cap C)$.

(i) 的包含关系 $\supset$.设 $x \in (A \cap B) \cup (A \cap C)$.则得
$x \in A \cap B$ 或 $x \in A \cap C$, 但在任何情形下都有 $x \in A$.若
$x \in A \cap B$, 则另外得到 $x \in B$, 从而 $x \in A \cap (B \cup C)$.
相应地, $x \in A \cap C$ 同样给出 $x \in A \cap (B \cup C)$.

(ii) 的包含关系 $\subset$.设 $x \in A \cup (B \cap C)$, 即 $x \in A$ 或
$x \in B \cap C$.对于 $x \in A$, 得 $x \in A \cup B$ 且 $x \in A \cup C$,
从而 $x \in (A \cup B) \cap (A \cup C)$.另一方面, $x \in B \cap C$ 自然给出
$x \in B \subset (A \cup B)$ 且 $x \in C \subset (A \cup C)$, 因此同样有
$x \in (A \cup B) \cap (A \cup C)$.

(ii) 的包含关系 $\supset$.设 $x \in (A \cup B) \cap (A \cup C)$.则特别地
$x \in A \cup B$; 对 $x \in A$, 得 $x \in A \cup (B \cap C)$.但若 $x \notin A$,
则又有 $x \in B$, 相应地 $x \in C$, 从而 $x \in B \cap C \subset A \cup (B \cap C)$.

(iii) 的包含关系 $\subset$.设 $x \in A - (B \cup C)$.则 $x \in A$, 但
$x \notin B \cup C$, 即 $x$ 既不属于 $B$ 也不属于 $C$.于是得
$x \in (A - B) \cap (A - C)$.

(iii) 的包含关系 $\supset$.设 $x \in (A - B) \cap (A - C)$.则 $x \in A - B$
且 $x \in A - C$, 故无论如何 $x \in A$, 并且 $x$ 既不属于 $B$ 也不属于 $C$.
但这意味着 $x \in A - (B \cup C)$.

(iv) 的包含关系 $\subset$.设 $x \in A - (B \cap C)$.则 $x \in A$, 但
$x \notin B \cap C$, 即 $x \notin B$ 或 $x \notin C$.在第一种情形得
$x \in A - B$, 在第二种得 $x \in A - C$; 无论何种情形都有
$x \in (A - B) \cup (A - C)$.

(iv) 的包含关系 $\supset$.设 $x \in (A - B) \cup (A - C)$, 即 $x \in A - B$
或 $x \in A - C$.则无论如何 $x \in A$.在第一种情形有 $x \notin B$, 从而也有
$x \notin B \cap C$.在第二种情形相应地得 $x \notin C$, 从而同样
$x \notin B \cap C$.总而言之得 $x \in A - (B \cap C)$.

\noindent\textbf{补充:} 等式 (iii) 与 (iv) 包含了命题的否定, 例如"$x$ 是 $Y$ 的
元素"这种类型.如人们容易确信的那样, 这对命题 $P$ 与 $Q$ 是按如下原则进行的:
\[
\overline{P \text{ 且 } Q} = \overline{P} \text{ 或 } \overline{Q}, \qquad
\overline{P \text{ 或 } Q} = \overline{P} \text{ 且 } \overline{Q},
\]
其中上横线表示对相应命题的否定.

\subsection*{习题 1.1/5}

\noindent\textbf{起点:} $X$ 是一个集合, $\mathfrak{P}(X)$ 是其幂集, 即 $X$ 的一切
子集的集合.此外给定一个映射 $f \colon X \longrightarrow \mathfrak{P}(X)$.

\noindent\textbf{目标:} 需证明 $f$ 不可能是满射.

\noindent\textbf{策略:} 由于映射 $f$ 并未以任何具体方式给出, 我们猜想所证论断是
建立在某个一般原则之上的.为探究这一原则, 我们考察简单集合 $X$ 的幂集
$\mathfrak{P}(X)$.空集 $X = \emptyset$ 不含任何元素, 而其幂集 $\mathfrak{P}(X)$
恰好含一个元素, 即 $X$.因此, 在这种情况下, 由于 $X$ 与 $\mathfrak{P}(X)$ 的
元素个数之故, 不可能存在满射 $f \colon X \longrightarrow \mathfrak{P}(X)$.
同一个论证也适用于由 $n$ 个不同元素组成的有限集合 $X = \{ x_1, \ldots, x_n \}$.
因为此时 $\mathfrak{P}(X)$ 含有 $n$ 个元素 $\{ x_1 \}, \ldots, \{ x_n \}$, 但此外
还含有其他元素, 例如空集, 故总共至少有 $n + 1$ 个元素.一个计数论证再次表明,
不可能存在满射 $f \colon X \longrightarrow \mathfrak{P}(X)$.

所考察的这些特殊情形提示我们, 在一般情形下不应试图显式地给出一个不属于 $f$ 之像
的 $\mathfrak{P}(X)$ 元素.相反, 应当采用间接方式, 即假设存在一个满射
$f \colon X \longrightarrow \mathfrak{P}(X)$, 并由此导出一个矛盾.刚才考察的有限
集合 $X$ 的例子正是这样运作的.\emph{罗素悖论}(\textup{Paradoxon von Russel})\index{luosubeilun@罗素悖论(Paradoxon von Russel)}为一般做法提供了一些线索.即考虑子集
\[
U = \{ x \in X ; x \notin f(x) \} \subset X.
\]
若 $f$ 是满射, 则存在 $y \in X$ 使得 $U = f(y)$, 而正如我们将看到的, 由此确实
可以导出一个矛盾.

\noindent\textbf{解题路径:} 我们于是假设 $f \colon X \longrightarrow \mathfrak{P}(X)$
是满射, 并考虑子集 $U = \{ x \in X ; x \notin f(x) \} \subset X$.则存在一个元素
$y \in X$ 使得 $U = f(y)$.若 $y \in U$, 则由 $U$ 的定义可知必有
$y \notin f(y) = U$, 这与 $y \in U$ 矛盾.另一方面, 由 $y \notin f(y) = U$ 又可得
$y \in U$, 也构成矛盾.总之, 满射 $f \colon X \longrightarrow \mathfrak{P}(X)$ 的
假设导致矛盾, 于是可知这样的映射不可能存在.

\noindent\textbf{补充:} 顺便提一下, 可以证明: 对含有 $n$ 个不同元素的有限集合 $X$,
其幂集 $\mathfrak{P}(X)$ 恰好有 $2^n$ 个元素.

\subsection*{习题 1.2/4}

\noindent\textbf{起点:} $G$ 是一个群, 带有两个子群 $H_1, H_2 \subset G$; 关于子群
的定义参见定义~\mtref{def:1.2.3}.

\noindent\textbf{目标:} 问的是在何种情形下 $H_1 \cup H_2$ 仍是 $G$ 的子群.对于
$H_1 \subset H_2$, 有 $H_1 \cup H_2 = H_2$; 相应地对于 $H_2 \subset H_1$, 有
$H_1 \cup H_2 = H_1$.因此在这些情形下, $H_1 \cup H_2$ 是 $G$ 的子群.要证明的是,
在所有其他情形下, $H_1 \cup H_2$ 不可能是 $G$ 的子群.

\noindent\textbf{策略:} 若取 $H_1 \cup H_2$ 中两个元素 $a, b$ 的乘积, 那么只要
$a, b$ 同属于 $H_1$ 或同属于 $H_2$, 自然就有 $ab \in H_1 \cup H_2$, 因为此时乘积
$ab$ 又属于 $H_1$ 或 $H_2$.但若取 $H_1 \cup H_2$ 中两个并非同属 $H_1$ 或 $H_2$ 的
元素 $a, b$ 的乘积, 情形就变得关键了, 例如 $a \in H_1 - H_2$ 且 $b \in H_2 - H_1$.
这种情形需要更仔细地考察.

\noindent\textbf{解题路径:} 我们已经看到, 当 $H_1 \subset H_2$ 或 $H_2 \subset H_1$
时, $H_1 \cup H_2$ 出于平凡的理由是 $G$ 的子群.反过来要证明的是, 只要
$H_1 \cup H_2$ 是 $G$ 的子群, 就有 $H_1 \subset H_2$ 或 $H_2 \subset H_1$ 这两个
包含关系之一成立.对这一步我们采用间接方法, 即把 $H_1 \cup H_2$ 设为 $G$ 的子群,
并假设既无 $H_1 \subset H_2$ 也无 $H_2 \subset H_1$, 目标是导出矛盾.于是, 若
$H_1 \not\subset H_2$ 且同时 $H_2 \not\subset H_1$, 则存在元素 $a \in H_1 - H_2$ 与
$b \in H_2 - H_1$.由于我们已假设 $H_1 \cup H_2$ 是 $G$ 的子群, 得
$ab \in H_1 \cup H_2$, 不妨设 $ab \in H_1$.则可从左侧乘以 $a^{-1} \in H_1$, 得到
$b \in H_1$, 与 $b$ 的选取矛盾.但若 $ab \in H_2$, 则从右侧乘以 $b^{-1} \in H_2$
相应地给出 $a \in H_2$, 也与 $a$ 的选取矛盾.因此, 若 $H_1 \cup H_2$ 是 $G$ 的子群,
则假设 $H_1 \not\subset H_2$ 且同时 $H_2 \not\subset H_1$ 是不能成立的.于是如所愿地
得到 $H_1 \subset H_2$ 或 $H_2 \subset H_1$.

\subsection*{习题 1.2/7}

\noindent\textbf{起点:} $G$ 是一个有限交换群; 设 $G = \{ g_1, \ldots, g_n \}$,
其中 $n$ 是 $G$ 的元素个数.

\noindent\textbf{目标:} 需证明
$\displaystyle\prod_{i=1}^{n} g_i^2 = 1$.

\noindent\textbf{策略:} 用展开的记号, 要计算乘积 $g_1 g_1 \cdots g_n g_n$, 并证明它
等于单位元 $1 \in G$.有哪些可能性可以简化这个乘积, 以证明其值为 $1$? 首先我们知道
单位元 $1$ 出现在元素 $g_1, \ldots, g_n$ 之中; 设 $g_1 = 1$.则仍须证明
$g_2 g_2 \cdots g_n g_n = 1$, 但这几乎算不上什么简化.另一方面, 由于 $G$ 是交换的,
我们在计算乘积时可以任意重新安排因子.此外我们记得, 对每个 $g_i \in G$, 在 $G$ 中
存在逆元 $g_j = g_i^{-1}$, 且此时 $g_i g_j = 1$.因此, 在乘积
$g_1^2 \cdots g_n^2$ 内部, 我们可以把元素 $g_i$ 与 $g_j$ 组合起来, 然后把 $g_i g_j$
替换为 $1$.这就是计算乘积 $\prod_{i=1}^{n} g_i^2$ 的关键.

\noindent\textbf{解题路径:} 映射 $\tau \colon G \longrightarrow G,\ g \longmapsto g^{-1}$
是双射, 因为由 $(g^{-1})^{-1} = g$($g \in G$)得 $\tau^2 = \operatorname{id}$.
由于 $G$ 是交换的, 我们可以如下计算:
\[
\prod_{g \in G} g^2 = \prod_{g \in G} g \cdot \prod_{g \in G} \tau(g)
= \prod_{g \in G} g \cdot \prod_{g \in G} g^{-1}
= \prod_{g \in G} (g g^{-1}) = \prod_{g \in G} 1 = 1\text{.}
\]

\subsection*{习题 1.3/2}

\noindent\textbf{起点:} 设 $K$ 是一个由有限多个元素组成的域.对 $n \in \mathbb{N}$
与 $a \in K$, 记 $n \cdot a$ 为 $a$ 与自身的 $n$ 倍和.

\noindent\textbf{目标:} 存在一个自然数 $n > 0$, 使得对所有 $a \in K$ 有
$n \cdot a = 0_K$, 其中 $0_K \in K$ 是零元.若 $n$ 是具有此性质的最小者, 则 $n$ 是
一个素数.

\noindent\textbf{策略:} 对一个元素 $a \in K$, 我们考察其所有倍数
$n \cdot a \in K$, $n \in \mathbb{N}$.这些倍数不可能全都不同, 因为
$\mathbb{N}$ 含有无穷多个元素, 而 $K$ 只含有限多个.因此存在元素
$n', n'' \in \mathbb{N}$, $n' < n''$, 使得 $n' \cdot a = n'' \cdot a$.由此得
\[
n' \cdot a = n'' \cdot a = n' \cdot a + (n'' - n') \cdot a,
\]
从而对 $n = n'' - n' > 0$ 有 $n \cdot a = 0_K$.于是对每个元素 $a \in K$, 都存在
一个自然数 $n > 0$ 使得 $n \cdot a = 0_K$.一般来说, $n$ 会依赖于 $a$ 的选取, 例如
对所有 $n \in \mathbb{N}$ 都有 $n \cdot 0_K = 0_K$, 但 $n \cdot 1_K \neq 0_K$ 当
$n = 1$ 且 $1_K \in K$ 是单位元时.然而, 通过一个技巧, 我们可以使自然数 $n$ 与 $a$
无关.例如, 若 $a_1, a_2 \in K$, 且 $n_1, n_2 \in \mathbb{N} - \{ 0 \}$ 满足
$n_i a_i = 0_K$($i = 1, 2$), 则得
\[
(n_1 n_2) \cdot a_1 = n_2 \cdot (n_1 \cdot a_1) = n_2 \cdot 0_K = 0_K,
\]
\[
(n_1 n_2) \cdot a_2 = n_1 \cdot (n_2 \cdot a_2) = n_1 \cdot 0_K = 0_K\text{.}
\]
如此这般, 我们终于可以找到作为自然数($> 0$)之乘积的 $n$, 使得对 $K$ 的每个元素
$a \in K$(共有限多个)都有 $n \cdot a = 0_K$.

可以想见, 用这种方法却不可能看出这样的 $n$ 在取最小者时是素数.因此我们必须考虑:
寻找满足对所有 $a \in K$ 都有 $n \cdot a = 0_K$ 的自然数 $n > 0$ 这一问题, 是否不能
归结到少数几个元素, 或者最好归结到单个元素 $a \in K$.为此我们注意到, 由
$n \cdot a = 0_K$ 自动推出 $n \cdot (a \cdot b) = (n \cdot a) \cdot b = 0_K$ 对进一步
的元素 $b \in K$ 成立.因此只须考察满足 $n \cdot 1_K = 0_K$ 的自然数 $n > 0$.
这个简化是解决该问题的关键.

\noindent\textbf{解题路径:} 我们上面已经证明, 存在自然数 $n > 0$ 使得
$n \cdot 1_K = 0_K$, 且此时对所有 $a \in K$ 得 $n \cdot a = (n \cdot 1_K) \cdot a = 0_K$.
现在选取具有此性质的最小 $n > 0$.自然有 $n > 1$.若 $n$ 不是素数, 则有分解
$n = n_1 \cdot n_2$, 其中 $n_1, n_2 < n$ 是自然数因子.由此得
\[
\begin{aligned}
0_K &= n \cdot 1_K = (n_1 \cdot n_2) \cdot 1_K = n_1 \cdot (n_2 \cdot 1_K) \\
&= n_1 \cdot \bigl( 1_K \cdot (n_2 \cdot 1_K) \bigr) = (n_1 \cdot 1_K) \cdot (n_2 \cdot 1_K),
\end{aligned}
\]
从而 $n_1 \cdot 1_K = 0_K$ 或 $n_2 \cdot 1_K = 0_K$, 这与 $n$ 的最小性矛盾.故 $n$
必为素数.

\noindent\textbf{补充:} 我们使用了一些诸如对 $n \in \mathbb{N}$,$a \in K$ 形如
$n \cdot a$ 的乘积的结合性之类的计算法则, 这些法则本应加以证明.对
$m, n \in \mathbb{Z}$, $a, b \in K$ 成立下列基本法则:
\[
\begin{array}{rcl}
(m \cdot n) \cdot a = m \cdot (n \cdot a), & & m \cdot (a \cdot b) = (m \cdot a) \cdot b, \\
(m + n) \cdot a = m \cdot a + n \cdot a, & & m \cdot (a + b) = m \cdot a + m \cdot b\text{.}
\end{array}
\]
关于这些法则的证明或多或少是"显然的", 但在更严格的意义上也可以用完全归纳法进行.

这一习题的主题日后还会借助环论的方法再次加以处理; 参见 \ref{sec:5.2} 节末尾关于域的特征的
论述以及那里的习题 9.

\subsection*{习题 1.3/5}

\noindent\textbf{起点:} 要考察的域是 $\mathbb{Q}(\sqrt{2})$, 它由形如
$a + b\sqrt{2}$($a, b \in \mathbb{Q}$)的一切实数组成.根据引理~\mtref{lem:1.3.4}
我们知道 $\sqrt{2} \notin \mathbb{Q}$.特别地, $\mathbb{Q}$ 是
$\mathbb{Q}(\sqrt{2})$ 的真子域.

\noindent\textbf{目标:} 需证明 $\sqrt{3} \notin \mathbb{Q}(\sqrt{2})$.

\noindent\textbf{策略:} 由于我们已经很好地了解了域 $\mathbb{Q}(\sqrt{2})$, 却对其中
不包含的实数知之甚少, 因此宜采用间接论证.故应假设
$\sqrt{3} \in \mathbb{Q}(\sqrt{2})$, 并尝试由此导出矛盾.

\noindent\textbf{解题路径:} 我们假设 $\sqrt{3} \in \mathbb{Q}(\sqrt{2})$.则存在有理数
$a, b \in \mathbb{Q}$ 使得 $a + b\sqrt{2} = \sqrt{3}$, 于是得到
\[
(a + b\sqrt{2})^2 = 3 \quad \text{亦即} \quad a^2 + 2b^2 - 3 = -2ab\sqrt{2}\text{.}
\]
对于 $ab \neq 0$, 我们可以把第二个方程对 $\sqrt{2}$ 求解, 将得到
$\sqrt{2} \in \mathbb{Q}$, 但这根据引理~\mtref{lem:1.3.4} 是被排除的.因此必有
$ab = 0$, 即 $a = 0$ 或 $b = 0$, 且 $a^2 + 2b^2 - 3 = 0$.对 $a = 0$ 得
$2b^2 = 3$.把 $b$ 取为约分后的分数 $\frac{p}{q}$, 其中 $p, q$ 是整数, 则得
$2p^2 = 3q^2$.与引理~\mtref{lem:1.3.4} 的证明一样可推出, $2$ 既是 $q$ 的因子也是
$p$ 的因子, 即分数 $\frac{p}{q}$ 未约分, 与我们的假设矛盾.

剩下要考察 $b = 0$ 的情形, 它归结为方程 $a^2 = 3$.仍把 $a$ 取为约分后的分数
$\frac{p}{q}$, 其中 $p, q$ 是整数.由此得 $p^2 = 3q^2$, 与
引理~\mtref{lem:1.3.4} 的证明一样可推出 $3$ 既是 $p$ 的因子也是 $q$ 的因子, 与我们的
假设矛盾.于是 $\sqrt{3} \in \mathbb{Q}(\sqrt{2})$ 总导致矛盾, 所以 $\sqrt{3}$ 不
可能是 $\mathbb{Q}(\sqrt{2})$ 的元素.

\noindent\textbf{补充:} 严格说来, 我们的论证以简化的形式使用了整数环 $\mathbb{Z}$
中的整除理论, 这将在 \ref{sec:5.2} 节再作详尽的处理.我们确实用到了这样的性质: 像 $2$ 或 $3$
这样的素数是素数元; 为此特别参见定理~\mtref{thm:5.2.12} 之前的论述.$\mathbb{Z}$ 中的
一个素数元 $p$, 参见定义~\mtref{def:5.2.10}, 其特征是: $p$ 整除 $\mathbb{Z}$ 中一个
乘积 $xy$ 当且仅当 $p$ 整除因子 $x$ 与 $y$ 之一.

\subsection*{习题 1.3/11}

\noindent\textbf{起点:} 对 $n, k \in \mathbb{N}$, $n \geq 1$, 考虑集合
\[
X(n, k) = \bigl\{ (a_1, \ldots, a_n) \in \mathbb{N}^n ;\ a_1 + \cdots + a_n = k \bigr\};
\]
记 $\#X(n, k)$ 为元素个数.

\noindent\textbf{目标:} 对 $n, k \in \mathbb{N}$, $n \geq 1$, 需证明下列论断:
\[
A(n, k) \colon \quad \#X(n, k) = \binom{k + n - 1}{n - 1}\text{.}
\]

\noindent\textbf{策略:} 由于我们不能通过明显的计数方式确定 $\#X(n, k)$, 我们尝试用
归纳法证明所求公式.首先显然有
\[
\#X(1, k) = 1 = \binom{k + 1 - 1}{1 - 1}, \qquad
\#X(n, 0) = 1 = \binom{0 + n - 1}{n - 1}
\]
对 $n \geq 1$, $k \geq 0$ 成立, 这些关系可供归纳基础使用.进一步, 对 $X(n+1, k)$ 中
的 $(n+1)$ 元组 $(a_1, \ldots, a_{n+1})$, 我们注意到
$\sum_{i=1}^{n} a_i = k - a_{n+1}$, 从而 $(a_1, \ldots, a_n) \in X(n, k - a_{n+1})$.
由于 $a_{n+1}$ 可以取 $0, \ldots, k$ 各值, 通过对 $X(n+1, k)$ 作相应的分割, 得
关系
\[
\#X(n+1, k) = \#X(n, k) + \#X(n, k-1) + \cdots + \#X(n, 0)\text{.}
\]
这个方程把 $X(n+1, k)$ 中 $(n+1)$ 元组所求的个数归结为若干 $n$ 元组的个数, 因此会
鼓励人们尝试对固定的 $k$ 关于 $n$ 作归纳.对于归纳步骤, 必须把所列的各个个数分别
换成相应的二项式系数, 并且必须验证如下关系:
\[
\binom{k + n}{n} = \binom{k + n - 1}{n - 1} + \cdots + \binom{0 + n - 1}{n - 1}\text{.}
\]
但由于右侧有多个加项, 这并不容易做到.由此我们学到: 应该把 $X(n+1, k)$ 以尽可能
简单的方式加以分割, 例如分为两部分
\[
\begin{array}{l}
X' = \bigl\{ (a_1, \ldots, a_{n+1}) \in X(n+1, k) ;\ a_{n+1} = 0 \bigr\}, \\
X'' = \bigl\{ (a_1, \ldots, a_{n+1}) \in X(n+1, k) ;\ a_{n+1} \geq 1 \bigr\},
\end{array}
\]
其中我们假设 $k \geq 1$.则 $X'$ 与 $X''$ 分别恰好有 $\#X(n, k)$ 与
$\#X(n+1, k-1)$ 个元素, 因为映射
\[
\begin{array}{l}
X' \longrightarrow X(n, k), \quad (a_1, \ldots, a_{n+1}) \longmapsto (a_1, \ldots, a_n), \\
X'' \longrightarrow X(n+1, k-1), \quad (a_1, \ldots, a_{n+1}) \longmapsto (a_1, \ldots, a_{n+1} - 1),
\end{array}
\]
是双射.从而成立等式
\[
\#X(n+1, k) = \#X(n, k) + \#X(n+1, k-1),
\]
我们虽不能把它用于对 $n$ 或 $k$ 的通常归纳步骤, 但或许可用于二重归纳.把满足
$n \geq 1$,$k \geq 0$ 的配对 $(n, k) \in \mathbb{N}^2$ 想象为实平面中的格点, 那么我们
知道所有形如 $(1, k)$ 与 $(n, 0)$ 的点上的所求个数.借助上述公式, 我们可以在点
$(n+1, k)$($k \geq 1$)计算个数, 倘若已知它在点 $(n, k)$ 与 $(n+1, k-1)$ 的值.
几何上可以相信, 通过反复应用这样的步骤, 从形如 $(1, k)$ 与 $(n, 0)$ 的点作为
"归纳基础"出发, 最终可以到达一切有关的格点.为了证明个数 $\#X(n, k)$ 由相应的
二项式系数给出, 我们还需在上述求和分解式中把各个个数替换为相应的二项式系数, 并检验
所得到的方程
\[
\binom{k + n}{n} = \binom{k + n - 1}{n - 1} + \binom{k + n - 1}{n}
\]
对 $n \geq 1$, $k \geq 1$ 成立.但这以乔装的形式正是关于二项式系数的著名公式
\[
\binom{n+1}{i} = \binom{n}{i-1} + \binom{n}{i},
\]
我们在证明二项式公式时已对 $1 \leq i \leq n$ 推导过它.所有这些都鼓励我们把上述思考
扩充为一个完整的解题方案.

\noindent\textbf{解题路径:} 我们用对 $n \geq 1$ 与 $k \geq 0$ 的二重归纳来证明论断
$A(n, k)$.为此需要下列步骤:

(I) $A(n, k)$ 对所有形如 $(1, k)$ 与 $(n, 0)$ 的配对 $(n, k)$ 成立, 这一点我们上面
已经验证过.

(II) 对 $n \geq 1$, $k \geq 1$ 成立如下蕴含式:
\[
\bigl[ A(n, k) \ \text{且}\ A(n+1, k-1) \bigr] \Longrightarrow A(n+1, k)\text{.}
\]
为证明这一蕴含式, 我们使用上面推导出的等式
\[
\#X(n+1, k) = \#X(n, k) + \#X(n+1, k-1)\text{.}
\]
结合 $A(n, k)$ 与 $A(n+1, k-1)$ 以及所提到的二项式系数公式, 得
\[
\#X(n+1, k) = \binom{k + n - 1}{n - 1} + \binom{k + n - 1}{n} = \binom{k + n}{n},
\]
即论断 $A(n+1, k)$.

我们现在还必须考虑, 步骤 (I) 与 (II) 是否足以得到 $A(n, k)$ 对所有 $n \geq 1$ 与
$k \geq 0$ 的成立.为此我们对 $n$ 作归纳, 证明 $A(n, k)$ 对所有 $n \geq 1$ 与任意
$k \geq 0$ 成立.由 (I) 我们知道 $A(1, k)$ 对所有 $k \geq 0$ 成立, 故归纳基础
$n = 1$ 有保证.现设 $n \geq 1$ 任意.我们假设 $A(n, k)$ 对所有 $k \geq 0$ 成立, 并想
由此推出 $A(n+1, k)$ 对所有 $k \geq 0$ 成立.为达此目的, 我们临时用一次对 $k$ 的
归纳.归纳基础 $k = 0$ 由 (I) 是显然的.现设 $k \geq 1$ 且 $A(n+1, k-1)$ 已知成立.
则结合 $A(n, k)$, 由蕴含式 (II) 得到 $A(n+1, k)$ 成立.对 $k$ 的归纳至此完成, 于是
$A(n+1, k)$ 对所有 $k \geq 0$ 成立, 这也完成了对 $n$ 的归纳.这样就证明了 $A(n, k)$
对所有 $n \geq 1$ 与 $k \geq 0$ 成立.

\noindent\textbf{补充:} 我们开始时考察了求和分解式
\[
\#X(n+1, k) = \#X(n, k) + \#X(n, k-1) + \cdots + \#X(n, 0),
\]
但不能用它计算 $\#X(n+1, k)$.由于我们现在知道各个个数都是二项式系数, 由此也就
证明了公式
\[
\binom{k + n}{n} = \binom{k + n - 1}{n - 1} + \cdots + \binom{0 + n - 1}{n - 1}
\]
对 $n \geq 1$, $k \geq 0$ 成立.

\subsection*{习题 1.4/1}

\noindent\textbf{起点:} $U$ 是 $K$-向量空间 $V$ 的一个线性子空间, 即根据
定义~\mtref{def:1.4.2}, $U$ 是 $V$ 中一个非空的,对加法与数乘封闭的子集.特别地,
有 $0 \in U$.

\noindent\textbf{目标:} 我们要刻画所有这样的元素 $a \in V$: 使得连同 $U$,
$a + U = \{ a + u ;\ u \in U \}$ 也是 $V$ 的线性子空间.

\noindent\textbf{策略:} 作为例子, 我们考虑域 $K = \mathbb{R}$ 上的向量空间
$V = \mathbb{R}^3$, 以及 $\mathbb{R}^3$ 中包含原点 $0 \in \mathbb{R}^3$ 的一个平面
$E \subset \mathbb{R}^3$.则 $E$ 是 $\mathbb{R}^3$ 中的线性子空间, 可以把向量
$a \in \mathbb{R}^3$ 的 $a + E$ 想象为把平面 $E$ 沿向量 $a$ 平行移动.对 $a \in E$,
$E$ 被双射地移到自身; 而对 $a \notin E$, 平移沿某个方向发生, 会离开 $E$.此时特别地,
$a + E$ 与 $E$ 不相交, 且 $a + E$ 不再可能是 $\mathbb{R}^3$ 的线性子空间, 因为零向量
$0$ 属于 $E$, 却不属于 $a + E$.

受这一几何例子的启发, 我们也在一般情形下试图证明: 对 $a \in U$, $a + U$ 与 $U$
重合; 而对 $a \notin U$, $a + U$ 不含零向量 $0 \in V$.因此 $a + U$ 是 $V$ 的线性
子空间当且仅当 $a \in U$.

\noindent\textbf{解题路径:} 先设 $a \in U$.由于 $U$ 作为 $V$ 的线性子空间对加法封闭,
得 $a + U \subset U$.由于 $U$ 与 $a$ 一起也包含 $a$, 相应地得 $(-a) + U \subset U$,
两侧同时加上 $a$ 得 $U \subset a + U$.故对 $a \in U$ 有 $a + U = U$, 且 $a + U$ 是
$V$ 的线性子空间.

接下来我们要弄清楚, 对哪些向量 $a \in V$, 集合 $a + U$ 会含有零向量.设
$0 \in a + U$.则存在向量 $b \in U$ 使得 $0 = a + b$.这意味着
$a = -b = (-1) \cdot b \in U$, 因为 $U$ 作为 $V$ 的线性子空间对数乘封闭.于是, 对
$a \notin U$, $a + U$ 不含零向量, 因此不可能是 $V$ 的线性子空间.总而言之我们断定:
$a + U$ 是 $V$ 的线性子空间当且仅当 $a \in U$.

\noindent\textbf{补充:} 另一种做法是, 尝试对 $a + U$ 验证定义~\mtref{def:1.4.2} 中所列
的线性子空间的诸定义性质, 以弄清对哪些 $a \in V$ 这可以做到.原则上这种途径是可行
的, 但最终同样归结为如下本质性认识: 对 $a \in U$ 有 $a + U = U$, 而对 $a \notin U$
有 $0 \notin a + U$.

\subsection*{习题 1.4/5}

\noindent\textbf{起点:} 给定两点 $x, y \in \mathbb{R}^2 - \{ 0 \}$, 它们不共同位于
经过 $0 \in \mathbb{R}^2$ 的同一条直线上; 这意味着 $\mathbb{R}x \neq \mathbb{R}y$,
特别地, 对所有 $\alpha, \beta \in \mathbb{R}^*$ 有 $\alpha x \neq \beta y$.

\noindent\textbf{目标:} $x, y$ 构成 $\mathbb{R}^2$ 的一个生成系, 即
$\mathbb{R}^2 = \langle x, y \rangle$, 其中根据定义--定理~\mtref{std:1.4.5} 有
$\langle x, y \rangle = \{ \alpha x + \beta y ;\ \alpha, \beta \in \mathbb{R} \}$.
此外还需检验, 若把 $\mathbb{R}$ 换成任意域 $K$, 相应的论断是否仍然成立.

\noindent\textbf{策略:} 几何直观已说明应当如何进行.点 $0, x, y$ 与 $x + y$ 张成
一个真正的平行四边形, 因为它们不位于经过 $0$ 的直线上.若 $z \in \mathbb{R}^2$ 非
平凡, 则可过 $z$ 作平行于边 $\overline{0x}$ 与 $\overline{0y}$ 的平行线, 并确定它们
与直线 $\mathbb{R}y$ 和 $\mathbb{R}x$ 的交点.由此得到一个由点 $0, z$ 以及直线
$\mathbb{R}x$ 上一点 $\alpha x$ 与直线 $\mathbb{R}y$ 上一点 $\beta y$ 张成的平行
四边形.于是得 $z = \alpha x + \beta y$, 如所愿.

原则上, 所需的直线及其交点都可以用公式描述或计算, 从而可以这种方式显式地算出常数
$\alpha, \beta \in \mathbb{R}$.但根据点 $x, y, z$ 的位置所需的分情形讨论, 使这一
途径并不吸引人.因此应当考虑如何减少计算量.

\noindent\textbf{解题路径:} 有 $\mathbb{R}^2 = \langle e_1, e_2 \rangle$, 其中
$e_1 = (1, 0)$, $e_2 = (0, 1)$ 是单位向量.若我们能证明 $e_1, e_2 \in \langle x, y
\rangle$, 则由关于生成线性子空间的法则立得
$\mathbb{R}^2 = \langle e_1, e_2 \rangle \subset \langle x, y \rangle \subset
\mathbb{R}^2$, 从而 $\langle x, y \rangle = \mathbb{R}^2$.为证
$e_1 \in \langle x, y \rangle$, 只需证明 $\langle x, y \rangle$ 含有类型
$(\varepsilon, 0)$ 且带有实常数 $\varepsilon \neq 0$ 的一个向量, 因为由此得
$e_1 = \varepsilon^{-1}(\varepsilon, 0) \in \langle x, y \rangle$.若
$x = (x_1, x_2)$ 且 $y = (y_1, y_2)$, 那么为了证明 $e_1 \in \langle x, y \rangle$,
我们最终必须找到常数 $\alpha, \beta \in \mathbb{R}$, 使得对某个
$\varepsilon \in \mathbb{R}^*$ 有 $\alpha x + \beta y = (\varepsilon, 0)$, 即
\[
\alpha x_1 + \beta y_1 \neq 0, \qquad \alpha x_2 + \beta y_2 = 0\text{.}
\]
自然存在不全为 $0$ 的常数 $\alpha, \beta \in \mathbb{R}$ 满足 $\alpha x_2 + \beta y_2
= 0$.此时必有 $\alpha x_1 + \beta y_1 \neq 0$, 因为否则将得 $\alpha x + \beta y = 0$,
即 $\alpha x = -\beta y$, 而这与前提条件--$x$ 和 $y$ 不得共同位于经过 $0$ 的直线
上--相矛盾.于是得 $e_1 \in \langle x, y \rangle$.类似地可以证明
$e_2 \in \langle x, y \rangle$, 然后断定 $\langle x, y \rangle = \mathbb{R}^2$.

最后, 若把 $\mathbb{R}$ 换成任意域 $K$, 具体的几何关联就消失了, 但所进行的计算论证
毫无限制地保持有效, 因此在那种情形下也有 $\langle x, y \rangle = K^2$.

\noindent\textbf{补充:} 解决此题需要运用 \ref{sec:1.4} 节的手段.另一方面, 也可以直接从将在
\ref{sec:1.5} 节处理的关于向量线性无关与向量空间维数的理论得出解答.此外, \ref{sec:4.4} 节中求解线性
方程组的克拉默法则也有帮助, 参见特别是推论~\mtref{cor:4.4.6}, 如果人们想要对给定的
$z \in \mathbb{R}^2$ 确定满足 $\alpha x + \beta y = z$ 的常数 $\alpha, \beta \in
\mathbb{R}$ 的话.

\subsection*{习题 1.5/2}

\noindent\textbf{起点:} $V$ 是带有线性子空间 $U, U' \subset V$ 的 $K$-向量空间, 其中
$U \cap U' = 0$.此外给定 $U$ 中线性无关的向量组 $x_1, \ldots, x_r$ 与 $U'$ 中线性
无关的向量组 $y_1, \ldots, y_s$.

\noindent\textbf{目标:} 向量 $x_1, \ldots, x_r, y_1, \ldots, y_s$ 构成 $V$ 中一个线性
无关的向量组.

\noindent\textbf{策略:} 我们试图按照定义~\mtref{def:1.5.1} 证明向量组
$x_1, \ldots, x_r, y_1, \ldots, y_s$ 的线性无关性, 并考察一个表示零向量的线性组合,
即
\[
\sum_{i=1}^{r} \alpha_i x_i + \sum_{j=1}^{s} \beta_j y_j = 0
\]
其中系数 $\alpha_i, \beta_j \in K$.由 $x_1, \ldots, x_r$ 的线性无关性, 只要我们知道
线性组合 $u = \sum_{i=1}^{r} \alpha_i x_i \in U$ 表示零向量, 就可以推出
$\alpha_1 = \cdots = \alpha_r = 0$.相应地, 若 $v = \sum_{j=1}^{s} \beta_j y_j \in U'$
表示零向量, 则 $\beta_1 = \cdots = \beta_s = 0$.实现等式 $u = 0$ 与 $v = 0$ 的唯一
机会, 在于利用前提 $U \cap U' = 0$.

\noindent\textbf{解题路径:} 设 $\alpha_1, \ldots, \alpha_r, \beta_1, \ldots, \beta_s
\in K$ 满足
\[
\sum_{i=1}^{r} \alpha_i x_i + \sum_{j=1}^{s} \beta_j y_j = 0\text{.}
\]
则考虑向量 $u = \sum_{i=1}^{r} \alpha_i x_i \in U$ 与 $v = \sum_{j=1}^{s} \beta_j y_j
\in U'$.由 $u + v = 0$, 即 $u = -v$ 与 $v = -u$, 得 $u, v \in U \cap U' = 0$,
从而 $u = v = 0$.向量组 $x_1, \ldots, x_r$ 与 $y_1, \ldots, y_s$ 的线性无关性于是
给出: 对所有 $i$ 有 $\alpha_i = 0$, 对所有 $j$ 有 $\beta_j = 0$, 由
定义~\mtref{def:1.5.1} 可知向量组 $x_1, \ldots, x_r, y_1, \ldots, y_s$ 在 $V$ 中线性
无关.

\noindent\textbf{补充:} 这个证明没有前提 $U \cap U' = 0$ 是不会成立的.若确实在
$U \cap U'$ 中存在一个向量 $u \neq 0$, 则它既可表示为 $x_i$ 的非平凡线性组合, 也可
表示为 $y_j$ 的非平凡线性组合, 例如 $u = \sum_{i=1}^{r} \alpha_i x_i$ 与
$v = \sum_{j=1}^{s} \beta_j y_j$, 则
\[
\sum_{i=1}^{r} \alpha_i x_i + \sum_{j=1}^{s} (-\beta_j) y_j = 0
\]
是 $0$ 的一个非平凡线性组合.即向量组 $x_1, \ldots, x_r, y_1, \ldots, y_s$ 将线性
相关.

\subsection*{习题 1.5/3}

\noindent\textbf{起点:} 要考察的是可变 $n \in \mathbb{N}$ 的 $\mathbb{R}$-向量空间
$\mathbb{R}^n$.

\noindent\textbf{目标:} 问的是哪些指数 $n \in \mathbb{N}$ 具有如下性质: 存在一个向量
序列 $a_1, a_2, \ldots \in \mathbb{R}^n$, 使得该序列任意两个向量都构成一个线性无关
的向量组.

\noindent\textbf{策略:} 我们已经看到 $\dim \mathbb{R}^n = n$.因此根据
推论~\mtref{cor:1.5.12}, 在 $\mathbb{R}^n$ 中只对 $n \geq 2$ 才可能存在长度为 $2$ 的
线性无关向量组.另一方面, 一个向量空间中两个向量线性相关, 当且仅当其中一个是另一个
的倍数.从而几何上可以相信: 两点 $x, y \in \mathbb{R}^n$ 恰好当 $x$ 与 $y$ 不位于
经过 $0$ 的直线上时, 才给出 $\mathbb{R}^n$ 中一个线性无关的向量组.由于在
$\mathbb{R}^2$ 中可以给出无穷多条不同的直线, 在 $\mathbb{R}^2$ 中应当也存在一个具有
所要求性质的无穷向量序列.把 $\mathbb{R}^2$ 视为 $\mathbb{R}^n$($n \geq 2$)的线性
子空间, 我们读得: 我们的问题对所有 $n \geq 2$ 可解, 而对 $n < 2$ 不可解, 正如我们
开头所见.

\noindent\textbf{解题路径:} 我们已经论证过, 只有在 $n \geq 2$ 时, $\mathbb{R}^n$ 中才
至多可能存在一个向量序列, 使得其中任意两个向量都构成线性无关的向量组.为了给出这样
的序列 $a_1, a_2, \ldots \in \mathbb{R}^2$, 我们取 $a_i = (i, 1)$, 并断言对 $i \neq j$,
向量 $a_i, a_j$ 线性无关.事实上, 设 $\alpha, \beta \in \mathbb{R}$ 满足
$\alpha a_i + \beta a_j = 0$, 即
\[
\alpha i + \beta j = 0, \qquad \alpha + \beta = 0\text{.}
\]
则得 $\alpha = -\beta$, 从而 $\beta \cdot (j - i) = 0$.由 $j - i \neq 0$ 得
$\beta = 0$, 进而 $\alpha = 0$.故向量 $a_i, a_j$ 在 $\mathbb{R}^2$ 中线性无关.
同样的论证也适用于 $n \geq 2$ 时的 $\mathbb{R}^n$, 只要把 $\mathbb{R}^2$ 等同于线性
子空间 $\{ (\alpha_1, \ldots, \alpha_n) \in \mathbb{R}^n ;\ \alpha_3 = \cdots = \alpha_n
= 0 \} \subset \mathbb{R}^n$.

\noindent\textbf{补充:} 这个例子非常漂亮地表明: 一个向量组的某些(真)部分组的线性
无关性, 一般说来不能说明整个向量组的线性无关性.但请注意, 一个向量组线性无关当且仅
当它的每个有限部分组都线性无关.我们在 \ref{sec:1.5} 节中对无穷向量组认识到了这一刻画, 但它
对有限向量组也平凡地成立.

\subsection*{习题 1.5/8}

\noindent\textbf{起点:} 对域 $K$ 与 $n \in \mathbb{N}$, 考虑 $K$-向量空间 $K^n$, 以及
如下形式的子集
\[
U = \left\{ (\alpha_1, \ldots, \alpha_n) \in K^n ;\ \sum_{i=1}^{n} \alpha_i \gamma_i = 0 \right\} \subset K^n
\]
其中 $\gamma_1, \ldots, \gamma_n \in K$ 是某些常数.

\noindent\textbf{目标:} $U$ 是 $K^n$ 的线性子空间, 需计算其相应的维数.

\noindent\textbf{策略:} 依据定义~\mtref{def:1.4.2} 验证 $U$ 的线性子空间诸性质没有任何
困难.但我们必须考虑如何确定 $U$ 的维数, 并从特殊情形 $n = 0$ 开始.此时
$U = K^0 = 0$, 从而 $\dim_K U = 0$.其次, 对 $n = 1$, $U$ 由
$U = \{ \alpha_1 \in K^1 ;\ \alpha_1 \gamma_1 = 0 \}$ 给出.而 $\alpha_1 \gamma_1 = 0$
当且仅当 $\alpha_1 = 0$ 或 $\gamma_1 = 0$.因此对 $\gamma_1 = 0$ 有 $U = K^1$, 对
$\gamma_1 \neq 0$ 有 $U = 0$, 相应地 $\dim_K U = \dim_K K^1 = 1$ 或
$\dim_K U = 0 = \dim_K K^1 - 1$.我们还应当考察特殊情形 $n = 2$, 即
\[
U = \left\{ (\alpha_1, \alpha_2) \in K^2 ;\ \alpha_1 \gamma_1 + \alpha_2 \gamma_2 = 0 \right\}\text{.}
\]
对于 $(\gamma_1, \gamma_2) = (0, 0)$, 同样可以看出 $U = K^2$, 从而
$\dim_K U = \dim_K K^2 = 2$.但若假设 $(\gamma_1, \gamma_2) \neq (0, 0)$, 例如
$\gamma_2 \neq 0$, 则我们还可以如下描述 $U$:
\[
U = \left\{ (\alpha_1, \alpha_2) \in K^2 ;\ \alpha_2 = -\gamma_2^{-1} \gamma_1 \cdot \alpha_1 \right\} = K \cdot (1, -\gamma_2^{-1} \gamma_1)\text{.}
\]
向量 $(1, -\gamma_2^{-1} \gamma_1)$ 于是构成 $U$ 的一个基, 我们得到
$\dim_K U = 1 = \dim_K K^2 - 1$.基于这些预考察, 可以设想一般情形的如下结果: 若
$\gamma_1 = \cdots = \gamma_n = 0$, 则 $\dim_K U = n$; 否则 $\dim_K U = n - 1$.

\noindent\textbf{解题路径:} 为证明 $U$ 是 $K^n$ 的线性子空间, 我们验证
定义~\mtref{def:1.4.2} 的条件.显然 $0 \in U$, 从而 $U \neq \emptyset$.再设
$a = (\alpha_1, \ldots, \alpha_n)$ 与 $b = (\beta_1, \ldots, \beta_n)$ 是 $U$ 的元素,
即 $\sum_{i=1}^{n} \alpha_i \gamma_i = 0$ 且 $\sum_{i=1}^{n} \beta_i \gamma_i = 0$, 则得
$\sum_{i=1}^{n} (\alpha_i + \beta_i) \gamma_i = 0$, 即 $a + b \in U$.此外, 对
$\alpha \in K$ 也有 $\sum_{i=1}^{n} \alpha \alpha_i \gamma_i = 0$, 从而
$\alpha \cdot a \in U$.于是 $U$ 被确认为 $K^n$ 的线性子空间.

为确定 $\dim_K U$, 我们首先注意到: 在平凡情形 $\gamma_1 = \cdots = \gamma_n = 0$ 中,
线性子空间 $U$ 与 $K^n$ 重合, 得到 $\dim_K U = \dim_K K^n = n$.现设 $\gamma_i$ 中至少
有一个不为 $0$, 例如 $\gamma_n \neq 0$, 这特别要求 $n \geq 1$.则
\[
U = \left\{ (\alpha_1, \ldots, \alpha_n) \in K^n ;\ \alpha_n = -\gamma_n^{-1} \sum_{i=1}^{n-1} \alpha_i \gamma_i \right\},
\]
我们看到, 对 $U$ 的元素 $(\alpha_1, \ldots, \alpha_n)$, 分量 $\alpha_1, \ldots,
\alpha_{n-1} \in K$ 不受任何条件约束, 而 $\alpha_n$ 由这些前几个分量唯一确定.设
$e'_1, \ldots, e'_{n-1} \in K^{n-1}$ 是标准单位向量, 则因此存在唯一确定的向量
$a_1, \ldots, a_{n-1} \in U$, 其前 $n-1$ 个分量分别与 $e'_1, \ldots, e'_{n-1}$ 的分量
一致.我们断言 $a_1, \ldots, a_{n-1}$ 是 $U$ 的一个基, 从而
$\dim_K U = n - 1$.事实上, 考察一个方程 $\sum_{i=1}^{n-1} \lambda_i a_i = 0$, 系数
$\lambda_i \in K$, 则删去下标为 $n$ 的分量后得方程 $\sum_{i=1}^{n-1} \lambda_i e'_i = 0$.
由于 $e'_1, \ldots, e'_{n-1}$ 构成 $K^{n-1}$ 的一个基, 所有 $\lambda_i$ 都平凡, 特别地,
向量 $a_1, \ldots, a_{n-1}$ 构成一个线性无关的向量组.以类似的方式可以证明这些向量
生成 $U$, 从而构成 $U$ 的一个基.为此我们考察任意向量 $a \in U$.删去最后一个分量后,
由它得到一个向量 $a' \in K^{n-1}$, 我们可以找到常数 $\lambda_1, \ldots, \lambda_{n-1}
\in K$, 使得 $a' = \sum_{i=1}^{n-1} \lambda_i e'_i$.则 $a$ 与
$\sum_{i=1}^{n-1} \lambda_i a_i$ 是 $U$ 中两个前 $n-1$ 个分量一致的向量.由于下标为 $n$
的分量由前 $n-1$ 个分量唯一确定, 得 $a = \sum_{i=1}^{n-1} \lambda_i a_i$.因此
$a_1, \ldots, a_{n-1}$ 构成 $U$ 的一个生成系, 从而构成一个基, 我们断定
$\dim_K U = n - 1$.

\noindent\textbf{补充:} 对 $\gamma_1, \ldots, \gamma_n \in K$ 可以考察映射
\[
f \colon K^n \longrightarrow K^1, \quad (\alpha_1, \ldots, \alpha_n) \longmapsto \sum_{i=1}^{n} \alpha_i \gamma_i\text{.}
\]
这是一个线性映射, 参见定义~\mtref{def:2.1.1}, 这意味着 $f$ 保持加法与数乘, 即对所有
$a, b \in K^n$ 与 $\alpha \in K$ 有
\[
f(a + b) = f(a) + f(b), \qquad f(\alpha \cdot a) = \alpha \cdot f(a)\text{.}
\]
此外, $U$ 与 $f^{-1}(0)$ 一致, 即所谓 $f$ 的核, 它由于 $f$ 的线性性始终是 $K^n$ 的
线性子空间.同样容易由 $f$ 的线性性看出, 像 $f(K^n)$ 是 $K^1$ 的线性子空间.对线性
映射现在成立所谓维数公式, 在我们的情形即
$\dim_K K^n = \dim_K f^{-1}(0) + \dim_K f(K^n)$, 我们将在定理~\mtref{thm:2.1.10} 中作
一般证明.由于
\[
f(K^n) = \begin{cases} 0 & \text{若 } \gamma_1 = \cdots = \gamma_n = 0, \\ K^1 & \text{否则,} \end{cases}
\]
我们可以轻易读出: 若 $\gamma_i$ 都平凡, 则 $\dim_K U = n$; 若 $\gamma_i$ 至少有一个
非零, 则 $\dim_K U = n - 1$.

\subsection*{习题 1.6/4}

\noindent\textbf{起点:} 给定 $K$-向量空间 $V$ 的线性子空间 $U_1, U_2, U_3$.

\noindent\textbf{目标:} 需证明如下维数公式:
\[
\dim_K U_1 + \dim_K U_2 + \dim_K U_3 = \dim_K (U_1 + U_2 + U_3) + \dim_K \bigl( (U_1 + U_2) \cap U_3 \bigr) + \dim_K (U_1 \cap U_2)\text{.}
\]

\noindent\textbf{策略:} 在定理~\mtref{thm:1.6.5} 中, 对两个线性子空间 $U, U' \subset V$
证明了维数公式
\[
\dim_K U + \dim_K U' = \dim_K (U + U') + \dim_K (U \cap U')\text{.}
\]
自然要把此公式应用于 $U = U_1 + U_2$ 与 $U' = U_3$, 何况在要证的公式右侧出现了项
$\dim_K ((U_1 + U_2) \cap U_3)$.

\noindent\textbf{解题路径:} 我们把维数公式定理~\mtref{thm:1.6.5} 应用于 $V$ 的线性子
空间 $U_1 + U_2$ 与 $U_3$, 得到
\[
\dim_K (U_1 + U_2) + \dim_K U_3 = \dim_K (U_1 + U_2 + U_3) + \dim_K \bigl( (U_1 + U_2) \cap U_3 \bigr)\text{.}
\]
再把维数公式应用于 $U_1$ 与 $U_2$, 得
\[
\dim_K U_1 + \dim_K U_2 = \dim_K (U_1 + U_2) + \dim_K (U_1 \cap U_2),
\]
于是总共
\[
\dim_K U_1 + \dim_K U_2 + \dim_K U_3 = \dim_K (U_1 + U_2 + U_3) + \dim_K \bigl( (U_1 + U_2) \cap U_3 \bigr) + \dim_K (U_1 \cap U_2),
\]
这就是所要证明的.

\noindent\textbf{补充:} 用完全归纳法可以类似地推导出对有限多个线性子空间
$U_1, \ldots, U_n \subset V$ 的维数公式.

\subsection*{习题 1.6/8}

\noindent\textbf{起点:} 给定一个有限维 $K$-向量空间 $V$ 及其线性子空间 $U \subset V$.
设 $n = \dim_K V$, $s = \dim_K U$.

\noindent\textbf{目标:} 需考虑: 对哪些 $r \in \mathbb{N}$, 以及在 $U$ 与 $V$ 的维数
满足何种条件之下, 存在 $U$ 在 $V$ 中的补空间 $U_1, \ldots, U_r$, 使其和
$\sum_{i=1}^{r} U_i$ 是直和.最优的是一个必要且充分的条件.

\noindent\textbf{策略:} 对 $r = 0$ 无须证明, 因为此时没有涉及 $U$ 的补空间, 且空和
总是取值 $0$.因此值得关注的只有个数 $r > 0$.其中 $r = 1$ 的情形是平凡的, 但
$r = 2$ 的情形已不能一眼看穿.我们由此学到: 参数 $r$ 不应处于中心位置, 而应服从
其余参数 $n$ 与 $s$.

为作进一步的定向, 我们看一个例子, 即 $\mathbb{R}$-向量空间 $V = \mathbb{R}^3$ 与
线性子空间(平面)$U = \{ (\alpha_1, \alpha_2, \alpha_3) \in \mathbb{R}^3 ;\ \alpha_3 = 0 \}$.
由于 $\dim U = 2$, 每个不包含于 $U$ 中的向量 $a = (\alpha_1, \alpha_2, \alpha_3) \in
\mathbb{R}^3$, 即满足 $\alpha_3 \neq 0$ 的向量, 都张成 $U$ 在 $V$ 中的一个补空间.
这样的向量 $a$ 有无限多个, 事实上甚至有不计其数的,不位于 $U$ 中的类型
$\mathbb{R} \cdot a \subset \mathbb{R}^3$ 的直线, 它们定义 $U$ 的补空间.若对
$i = 1, \ldots, r$, $\mathbb{R} \cdot a_i$ 都是这样的补空间, 且和
$\sum_{i=1}^{r} \mathbb{R} \cdot a_i$ 是直和, 则向量 $a_1, \ldots, a_r$ 构成一个线性
无关的向量组, 其中由 $\dim \mathbb{R}^3 = 3$ 必然推出 $r \leq 3$.另一方面, 也确实
存在一个不在 $U$ 中,从而张成 $U$ 在 $\mathbb{R}^3$ 中的补空间的线性无关向量组
$a_1, a_2, a_3 \in \mathbb{R}^3$, 即向量
\[
(0, 0, 1), \qquad (1, 0, 1), \qquad (1, 1, 1)\text{.}
\]
上述考察在一般情形下也把我们引向一个至少是必要的条件: 若存在 $U$ 在 $V$ 中的补空间
$U_1, \ldots, U_r$, 其和 $\sum_{k=1}^{r} U_k \subset V$ 是直和, 则由
定理~\mtref{thm:1.6.4} 对有限多个被加项的推广得
\[
\sum_{k=1}^{r} \dim_K U_k \leq \dim_K V = n\text{.}
\]
又由定理~\mtref{thm:1.6.4} 得, 对所有 $k$ 有 $\dim_K U_k = \dim_K V - \dim_K U = n - s$,
于是给出必要条件 $r \cdot (n - s) \leq n$.在我们上面的例子中, $n = 3$, $s = 2$, 条件
于是取形式 $r \leq 3$, 并且如我们所看到的, 它尤其也是充分的.

为在一般情形下也找到我们问题的充分条件, 我们首先来考虑, 可以用什么方式构造 $U$ 在
$V$ 中的不同补空间.我们首先选取 $U$ 在 $V$ 中的任意一个补空间 $U'$.与
定理~\mtref{thm:1.6.4} 的证明一样, 它可以如下构造: 选取 $U$ 的一个基 $x_1, \ldots, x_s$,
并用元素 $y_1, \ldots, y_t$($t = n - s$)把它扩充为 $V$ 的一个基.则
$U' = \bigoplus_{j=1}^{t} K y_j$ 是 $U$ 在 $V$ 中的一个补空间.若我们现在想改变 $U'$,
那么有意义的做法是改变 $U'$ 的生成元 $y_j$, 而且--为了有把握离开 $U'$--用 $U$ 中
的元素来改变.因此我们选取元素 $a_1, \ldots, a_t \in U$, 并检验向量
$y_1 + a_1, \ldots, y_t + a_t$ 是否也张成一个补空间.为此只需检验这些元素连同
$x_1, \ldots, x_s$ 构成一个线性无关的向量组, 从而构成 $V$ 的一个基.于是设
$\lambda_1, \ldots, \lambda_s, \lambda'_1, \ldots, \lambda'_t \in K$ 满足
\[
\sum_{i=1}^{s} \lambda_i x_i + \sum_{j=1}^{t} \lambda'_j (y_j + a_j) = \left( \sum_{i=1}^{s} \lambda_i x_i + \sum_{j=1}^{t} \lambda'_j a_j \right) + \sum_{j=1}^{t} \lambda'_j y_j = 0\text{.}
\]
则 $\sum_{i=1}^{s} \lambda_i x_i + \sum_{j=1}^{t} \lambda'_j a_j \in U$ 且
$\sum_{j=1}^{t} \lambda'_j y_j \in U'$.由于和 $U + U'$ 是直和, 这两个元素必都为零.
由此得 $\lambda'_1, \ldots, \lambda'_t = 0$, 进而 $\lambda_1, \ldots, \lambda_s = 0$, 即
$U'' = \sum_{j=1}^{t} K \cdot (y_j + a_j)$ 确实是 $U$ 在 $V$ 中的又一个补空间, 而且当
向量 $a_1, \ldots, a_t$ 中至少有一个非平凡时, 它不同于 $U'$.我们想检验: 和
$U' + U''$ 能否是直和, 即向量 $y_1, \ldots, y_t, y_1 + a_1, \ldots, y_t + a_t$ 是否
构成一个线性无关的向量组.因此设 $\lambda'_1, \ldots, \lambda'_t, \lambda''_1, \ldots,
\lambda''_t \in K$, 使得
\[
\sum_{j=1}^{t} \lambda'_j y_j + \sum_{j=1}^{t} \lambda''_j (y_j + a_j) = \sum_{j=1}^{t} \lambda''_j a_j + \sum_{j=1}^{t} (\lambda'_j + \lambda''_j) y_j = 0\text{.}
\]
则 $\sum_{j=1}^{t} \lambda''_j a_j \in U$ 且 $\sum_{j=1}^{t} (\lambda'_j + \lambda''_j) y_j
\in U'$, 而由于和 $U + U'$ 是直和, 两个元素又必为零.于是可以对所有 $j$ 推出
$\lambda'_j + \lambda''_j = 0$, 但只有当向量 $a_1, \ldots, a_t$ 在 $U$ 中构成一个线性
无关的向量组时, 才能对所有 $j$ 推出 $\lambda''_j = 0$.后者例如在 $t \leq s$ 且对
$j = 1, \ldots, t$ 取 $a_j = x_j$ 时成立.这是关键性观察, 它将向我们表明: 条件
$r \cdot (n - s) \leq n$ 对我们的问题的求解确实也是充分的.

\noindent\textbf{解题路径:} 我们要证明: 存在 $U$ 在 $V$ 中的补空间 $U_1, \ldots, U_r$,
其和是直和, 当且仅当 $r(n - s) \leq n$.正如我们所见, 这一条件是必要的.若
$U_1, \ldots, U_r$ 确是如此的补空间, 则由定理~\mtref{thm:1.6.4} 对所有 $k = 1, \ldots, r$
有 $\dim_K U_k = n - s$, 于是由推论~\mtref{cor:1.5.14} 与定理~\mtref{thm:1.6.4} 对有限多
个被加项的推广得
\[
r(n - s) = \sum_{k=1}^{r} \dim_K U_k = \dim_K \bigoplus_{k=1}^{r} U_k \leq \dim_K V = n\text{.}
\]
反之, 现设 $r \in \mathbb{N}$ 满足 $r(n - s) \leq n$.则得 $(r - 1)(n - s) \leq s$, 即
$(r - 1)t \leq s$, 其中 $t = n - s$.此外我们可以假定 $r \geq 2$, 因为要证的论断在其
余情形下是平凡的.现在我们在 $U$ 中选取一个由 $(r - 1)t$ 个向量组成的线性无关向量组,
写成形式
\[
x_{jk}, \qquad j = 1, \ldots, t, \qquad k = 2, \ldots, r;
\]
这是可能的, 因为 $(r - 1)t \leq s = \dim_K U$.又设 $U_1$ 是 $U$ 在 $V$ 中的任意一个
补空间, $y_1, \ldots, y_t$ 是 $U_1$ 的一个基, 其中我们利用 $t = n - s = \dim_K U_1$.
然后令
\[
U_k = \sum_{j=1}^{t} K \cdot (y_j + x_{jk}) \quad \text{对} \ k = 2, \ldots, r\text{.}
\]
由于显然对所有 $j = 1, \ldots, t$ 与 $k = 2, \ldots, r$ 有 $y_j \in U + U_k$, 得
$U + U_k = V$ 对所有 $k$ 成立, 并且借助推论~\mtref{cor:1.6.6} 从维数理由得知和
$U + U_k$ 是直和.于是除 $U_1$ 之外, $U_2, \ldots, U_r$ 也是 $U$ 在 $V$ 中的补空间.
为证明和 $\sum_{k=1}^{r} U_k$ 是直和, 我们考察一个线性组合
\[
\sum_{j=1}^{t} \lambda_j y_j + \sum_{k=2}^{r} \sum_{j=1}^{t} \lambda_{jk} (y_j + x_{jk}) = 0,
\]
系数 $\lambda_j, \lambda_{jk} \in K$.这个方程也可以写成形式
\[
\left( \sum_{j=1}^{t} \lambda_j y_j + \sum_{k=2}^{r} \sum_{j=1}^{t} \lambda_{jk} y_j \right) + \left( \sum_{k=2}^{r} \sum_{j=1}^{t} \lambda_{jk} x_{jk} \right) = 0,
\]
其中左边括号中的表达式属于 $U_1$, 右边的属于 $U$.由于和 $V = U + U_1$ 是直和, 两个
表达式必都为零.由 $x_{jk}$ 的线性无关性推出所有系数 $\lambda_{jk}$ 为零, 由 $y_j$ 的
线性无关性推出系数 $\lambda_j$ 也为零.于是 $y_j$($j = 1, \ldots, t$)连同向量
$y_j + x_{jk}$($j = 1, \ldots, t$, $k = 2, \ldots, r$)构成 $V$ 中一个线性无关的向量
组, 由此得和 $\sum_{k=1}^{r} U_k$ 在 $V$ 中是直和.这样就证明了条件 $r(n - s) \leq n$
对解决我们的问题不仅必要而且充分.

\noindent\textbf{补充:} 若 $x_1, \ldots, x_s$ 是 $U$ 的一个基, 则在上面的证明中还可以
更具体地选取向量 $x_{jk}$, 即令
\[
x_{jk} = x_{j + (k-2) t}, \qquad j = 1, \ldots, t, \qquad k = 2, \ldots, r\text{.}
\]

\section*{对应 2.1-2.3: 线性映射}

\subsection*{习题 2.1/2}

\noindent\textbf{起点:} 在不同的情形 (i)--(iv) 中, 要考虑的是若干显式给出的向量 $a_i \in \mathbb{R}^4$ 以及 $b_i \in \mathbb{R}^3$($i = 1, 2, 3, 4$ 或 $i = 1, 2, 3$).

\noindent\textbf{目标:} 问是否存在 $\mathbb{R}$-线性映射 $\mathbb{R}^4 \longrightarrow \mathbb{R}^3$, 使它对所有 $i$ 都把 $a_i$ 映到 $b_i$.

\noindent\textbf{策略:} 我们在定理~\mtref{thm:2.1.7} 及其相应的证明中学到, 可以通过线性扩张来定义一个 $\mathbb{R}$-线性映射 $f \colon \mathbb{R}^4 \longrightarrow \mathbb{R}^3$: 考虑 $\mathbb{R}^4$ 的一个基 $a_1, \ldots, a_4$, 并任意指定它们的像 $b_1, \ldots, b_4 \in \mathbb{R}^3$.所得的 $\mathbb{R}$-线性映射由 $f(a_i) = b_i$, $i = 1, 2, 3, 4$, 唯一确定.因此在相应的情形中, 如果向量 $a_1, \ldots, a_4$ 线性无关, 那么它们构成 $\mathbb{R}^4$ 的一个基, 并且总存在唯一确定的 $\mathbb{R}$-线性映射 $f \colon \mathbb{R}^4 \longrightarrow \mathbb{R}^3$, 把 $a_i$ 映到相应的向量 $b_i \in \mathbb{R}^3$.如果只给出向量 $a_1, a_2, a_3 \in \mathbb{R}^4$ 且它们构成一个线性无关组, 那么可以类似地处理.借助定理~\mtref{thm:1.5.8} 用一个向量 $a$ 把这个组扩充为 $\mathbb{R}^4$ 的一个基, 就可以找到 $\mathbb{R}$-线性映射 $f \colon \mathbb{R}^4 \longrightarrow \mathbb{R}^3$, 把 $a_1, a_2, a_3$ 映到 $b_1, b_2, b_3$, 并把 $a$ 映到另一个任意选取的向量 $b \in \mathbb{R}^3$(例如 $b = 0$).这时映射 $f$ 自然依赖于 $b$ 的选取, 因而不再由 $a_1, a_2, a_3$ 到 $b_1, b_2, b_3$ 的像唯一确定.

如果向量 $a_1, \ldots, a_4 \in \mathbb{R}^4$ 构成一个线性相关组, 那该怎么办? 这时可以考虑由 $a_1, \ldots, a_4$ 生成的线性子空间 $U = \langle a_1, \ldots, a_4 \rangle \subset \mathbb{R}^4$, 并把组 $a_1, \ldots, a_4$ 缩减为 $U$ 的一个基.例如, 如果 $a_1, a_2, a_3$ 构成 $U$ 的一个基, 那么正如我们所看到的, 总存在 $\mathbb{R}$-线性映射 $f \colon \mathbb{R}^4 \longrightarrow \mathbb{R}^3$, 把 $a_1, a_2, a_3$ 映到 $b_1, b_2, b_3$.按定理~\mtref{thm:2.1.7} 这个映射在 $U$ 上唯一确定.然后通过把 $a_4$ 表示成 $a_1, a_2, a_3$ 的线性组合, 来检验是否也有 $f(a_4) = b_4$.在第一种情形存在所希望类型的 $\mathbb{R}$-线性映射 $\mathbb{R}^4 \longrightarrow \mathbb{R}^3$, 在第二种情形则不存在.

\noindent\textbf{解题路径:} (i) 我们断言, 向量组
\[
a_1 = (1, 1, 0, 0), \quad a_2 = (1, 1, 1, 0), \quad a_3 = (0, 1, 1, 1), \quad a_4 = (0, 0, 1, 1)
\]
线性无关, 从而构成 $\mathbb{R}^4$ 的一个基.为验证这一点, 我们考虑系数 $\lambda_i \in \mathbb{R}$ 的方程 $\sum_{i = 1}^{4} \lambda_i a_i = 0$, 即按分量的写法考虑下列方程组:
\[
\begin{array}{l}
\lambda_1 \cdot 1 + \lambda_2 \cdot 1 + \lambda_3 \cdot 0 + \lambda_4 \cdot 0 = 0\\
\lambda_1 \cdot 1 + \lambda_2 \cdot 1 + \lambda_3 \cdot 1 + \lambda_4 \cdot 0 = 0\\
\lambda_1 \cdot 0 + \lambda_2 \cdot 1 + \lambda_3 \cdot 1 + \lambda_4 \cdot 1 = 0\\
\lambda_1 \cdot 0 + \lambda_2 \cdot 0 + \lambda_3 \cdot 1 + \lambda_4 \cdot 1 = 0
\end{array}
\]
前两个方程给出 $\lambda_3 = 0$, 后两个给出 $\lambda_2 = 0$.于是又得 $\lambda_1 = 0$ 和 $\lambda_4 = 0$, 我们看到向量 $a_1, a_2, a_3, a_4$ 构成一个线性无关组, 从而是 $\mathbb{R}^4$ 的一个基.于是由定理~\mtref{thm:2.1.7} 读出, 存在 $\mathbb{R}$-线性映射 $\mathbb{R}^4 \longrightarrow \mathbb{R}^3$, 把 $a_1, a_2, a_3, a_4$ 映到预先给定的向量 $b_1, b_2, b_3, b_4 \in \mathbb{R}^3$.

(ii) 我们断言, 这里的向量组
\[
a_1 = (0, 1, 1, 1), \quad a_2 = (1, 0, 1, 1), \quad a_3 = (1, 1, 0, 1)
\]
也线性无关.同样考虑系数 $\lambda_i \in \mathbb{R}$ 的方程 $\sum_{i = 1}^{3} \lambda_i a_i = 0$, 即下列方程组:
\[
\begin{array}{l}
\lambda_1 \cdot 0 + \lambda_2 \cdot 1 + \lambda_3 \cdot 1 = 0\\
\lambda_1 \cdot 1 + \lambda_2 \cdot 0 + \lambda_3 \cdot 1 = 0\\
\lambda_1 \cdot 1 + \lambda_2 \cdot 1 + \lambda_3 \cdot 0 = 0\\
\lambda_1 \cdot 1 + \lambda_2 \cdot 1 + \lambda_3 \cdot 1 = 0
\end{array}
\]
第一个和第四个方程给出 $\lambda_1 = 0$, 第二个和第四个给出 $\lambda_2 = 0$.于是也有 $\lambda_3 = 0$, 我们看到 $a_1, a_2, a_3$ 构成一个线性无关组.把这个组扩充为 $\mathbb{R}^4$ 的一个基, 再次由定理~\mtref{thm:2.1.7} 看到, 存在 $\mathbb{R}$-线性映射 $f \colon \mathbb{R}^4 \longrightarrow \mathbb{R}^3$, 把 $a_1, a_2, a_3$ 映到预先给定的向量 $b_1, b_2, b_3 \in \mathbb{R}^3$.再次应用定理~\mtref{thm:2.1.7} 得到, 只有 $f$ 到线性子空间 $\langle a_1, a_2, a_3 \rangle \subset V$ 上的限制是唯一确定的.

(iii) 我们再次检验向量组
\[
a_1 = (0, 1, 1, 1), \qquad a_2 = (1, 0, 1, 1), \qquad a_3 = (1, 1, 0, 1), \qquad a_4 = (-1, 1, 0, 0)
\]
是否线性无关, 并设系数 $\lambda_i \in \mathbb{R}$ 的方程 $\sum_{i = 1}^{4} \lambda_i a_i = 0$.因此要解下列方程组:
\[
\begin{array}{l}
\lambda_1 \cdot 0 + \lambda_2 \cdot 1 + \lambda_3 \cdot 1 - \lambda_4 \cdot 1 = 0\\
\lambda_1 \cdot 1 + \lambda_2 \cdot 0 + \lambda_3 \cdot 1 + \lambda_4 \cdot 1 = 0\\
\lambda_1 \cdot 1 + \lambda_2 \cdot 1 + \lambda_3 \cdot 0 + \lambda_4 \cdot 0 = 0\\
\lambda_1 \cdot 1 + \lambda_2 \cdot 1 + \lambda_3 \cdot 1 + \lambda_4 \cdot 0 = 0
\end{array}
\]
第一个和第四个方程给出 $\lambda_1 + \lambda_4 = 0$, 第二个和第四个给出 $\lambda_2 - \lambda_4 = 0$, 第三个和第四个给出 $\lambda_3 = 0$.特别地, 基于这些关系, 值
\[
\lambda_1 = 1, \qquad \lambda_2 = -1, \qquad \lambda_3 = 0, \qquad \lambda_4 = -1
\]
是可能的, 而代入上述方程表明, 这的确是方程组的一个解.于是 $a_1 - a_2 - a_4 = 0$, 即 $a_4 = a_1 - a_2$, 向量组 $a_1, a_2, a_3, a_4$ 线性相关.

如我们在 (ii) 中所见, 向量 $a_1, a_2, a_3$ 构成一个线性无关组, 且存在 $\mathbb{R}$-线性映射 $f \colon \mathbb{R}^4 \longrightarrow \mathbb{R}^3$ 满足 $f(a_i) = b_i$, $i = 1, 2, 3$, 其中 $f$ 到线性子空间 $\langle a_1, a_2, a_3 \rangle \subset V$ 上的限制唯一确定.此外有
\[
f(a_4) = f(a_1 - a_2) = f(a_1) - f(a_2) = b_1 - b_2 = (-1, -1, 2) \neq b_4\text{.}
\]
在这种情形下, 因此不存在满足 $f(a_i) = b_i$, $i = 1, 2, 3, 4$, 的 $\mathbb{R}$-线性映射 $f \colon \mathbb{R}^4 \longrightarrow \mathbb{R}^3$.

(iv) 在向量
\[
a_1 = (0, 1, 1, 1), \qquad a_2 = (1, 0, 1, 1), \qquad a_3 = (1, 1, 0, 1), \qquad a_4 = (0, 2, 0, 1)
\]
的情形, 我们同样考虑系数 $\lambda_i \in \mathbb{R}$ 的方程 $\sum_{i = 1}^{4} \lambda_i a_i = 0$, 它把我们引到下列方程组:
\[
\begin{array}{l}
\lambda_1 \cdot 0 + \lambda_2 \cdot 1 + \lambda_3 \cdot 1 + \lambda_4 \cdot 0 = 0\\
\lambda_1 \cdot 1 + \lambda_2 \cdot 0 + \lambda_3 \cdot 1 + \lambda_4 \cdot 2 = 0\\
\lambda_1 \cdot 1 + \lambda_2 \cdot 1 + \lambda_3 \cdot 0 + \lambda_4 \cdot 0 = 0\\
\lambda_1 \cdot 1 + \lambda_2 \cdot 1 + \lambda_3 \cdot 1 + \lambda_4 \cdot 1 = 0
\end{array}
\]
第一个和第四个方程给出 $\lambda_1 + \lambda_4 = 0$, 第二个和第四个给出 $\lambda_2 - \lambda_4 = 0$, 第三个和第四个给出 $\lambda_3 + \lambda_4 = 0$.基于这些关系, 值
\[
\lambda_1 = 1, \qquad \lambda_2 = -1, \qquad \lambda_3 = 1, \qquad \lambda_4 = -1
\]
是可能的, 而代入上述方程表明, 这的确是一个解.于是 $a_4 = a_1 - a_2 + a_3$.与 (ii) 一样, 向量 $a_1, a_2, a_3$ 线性无关, 与 (iii) 一样, 存在 $\mathbb{R}$-线性映射 $f \colon \mathbb{R}^4 \longrightarrow \mathbb{R}^3$ 满足 $f(a_i) = b_i$, $i = 1, 2, 3$.现在有
\[
f(a_4) = f(a_1 - a_2 + a_3) = f(a_1) - f(a_2) + f(a_3) = b_1 - b_2 + b_3 = (2, 0, 4) = b_4\text{.}
\]
因此 $f \colon \mathbb{R}^4 \longrightarrow \mathbb{R}^3$ 正如所希望的那样, 是满足 $f(a_i) = b_i$, $i = 1, 2, 3, 4$, 的 $\mathbb{R}$-线性映射.

\noindent\textbf{补充:} 在确定系数 $\lambda_i$ 时所考虑的方程组是所谓的线性方程组.这些方程组及其相关的实用求解方法的理论, 将在后面的 \ref{sec:3.5} 节中详细讨论.

\subsection*{习题 2.1/5}

\noindent\textbf{起点:} 给定 $K$-向量空间 $V$ 的一个自同态 $f \colon V \longrightarrow V$, 满足 $f^2 = f$.

\noindent\textbf{目标:} 要证明 $V = \ker f \oplus \operatorname{im} f$, 即线性子空间的和 $\ker f + \operatorname{im} f \subset V$ 是直和并等于 $V$.

\noindent\textbf{策略:} 对每个自同态 $f \colon V \longrightarrow V$ 都可以把和 $\ker f + \operatorname{im} f$ 作为 $V$ 的线性子空间来构造.这个和一般不是直的, 例如由定理~\mtref{thm:2.1.7} (ii) 定义的,带典范单位向量 $e_1, e_2 \in K^2$ 的自同态
\[
f \colon K^2 \longrightarrow K^2, \qquad e_1 \longmapsto 0, \qquad e_2 \longmapsto e_1,
\]
就表明了这一点; 这里 $f^2 = 0 \neq f$, 且 $\ker f = \operatorname{im} f = K e_1$.我们由此学到, 条件 $f^2 = f$ 对和 $\ker f + \operatorname{im} f$ 的直和性将是关键的.此外, 如果已认识到这个和是直的, 那么由定理~\mtref{thm:1.6.4} 与维数公式(定理~\mtref{thm:2.1.10})得
\[
\dim_K(\ker f \oplus \operatorname{im} f) = \dim_K(\ker f) + \dim_K(\operatorname{im} f) = \dim_K V\text{.}
\]
因此如果 $V$ 有限维, 我们就能借助推论~\mtref{cor:1.5.14} (ii) 直接读出 $V = \ker f \oplus \operatorname{im} f$.

于是要说明的是, 和 $\ker f + \operatorname{im} f \subset V$ 是直的, 即 $\ker f \cap \operatorname{im} f = 0$.这可以例行地验证.设 $x \in \ker f \cap \operatorname{im} f$.则一方面 $f(x) = 0$, 另一方面存在向量 $x' \in V$ 使 $f(x') = x$.于是 $f^2(x') = 0$, 又因 $f^2 = f$ 还有 $x = f(x') = f^2(x') = 0$, 即 $\ker f \cap \operatorname{im} f = 0$.

还须证明 $V = \ker f + \operatorname{im} f$, 而不假定 $V$ 有限维.为此考虑向量 $x \in V$, 尝试找到向量 $x'' \in \operatorname{im} f$, 使得 $x - x'' \in \ker f$.这里一个显然的候选是 $x'' = f(x)$.

\noindent\textbf{解题路径:} 我们已经证明和 $\ker f + \operatorname{im} f \subset V$ 是直的.为看出 $\ker f + \operatorname{im} f = V$, 选取 $x \in V$.则得
\[
f(x - f(x)) = f(x) - f^2(x) = f(x) - f(x) = 0,
\]
即 $x - f(x) \in \ker f$, 从而 $x \in \ker f + \operatorname{im} f$, 这正是要证明的.

\noindent\textbf{补充:} 如我们所认识到的, 条件 $f^2 = f$ 对关系式 $V = \ker f \oplus \operatorname{im} f$ 是充分的.但它不是必要的, 因为对向量空间 $V \neq 0$, 总存在满足 $f^2 \neq f$ 的自同构 $f \colon V \longrightarrow V$.若 $f$ 是这样的自同构, 则 $\ker f = 0$ 且 $\operatorname{im} f = V$, 从而平凡地有 $V = \ker f \oplus \operatorname{im} f$.

\subsection*{习题 2.1/8}

\noindent\textbf{起点:} 给定有限维 $K$-向量空间之间的 $K$-线性映射 $V_1 \xrightarrow{f} V_2 \xrightarrow{g} V_3$.

\noindent\textbf{目标:} 要证明 $\operatorname{rg} f + \operatorname{rg} g \leq \operatorname{rg}(g \circ f) + \dim_K V_2$, 其中\emph{线性映射的秩}(\textup{Rang einer linearen Abbildung})\index{xianxingyingshedezhi@线性映射的秩(Rang einer linearen Abbildung)}定义为相应像的维数.

\noindent\textbf{策略:} 由于在待证不等式中出现 $K$-线性映射 $f$, $g$ 以及 $g \circ f$ 的秩, 我们尝试使用维数公式(定理~\mtref{thm:2.1.10})的相应版本, 即
\[
\begin{array}{l}
\dim_K V_1 = \dim_K(\ker f) + \operatorname{rg} f,\\
\dim_K V_2 = \dim_K(\ker g) + \operatorname{rg} g,\\
\dim_K V_1 = \dim_K(\ker(g \circ f)) + \operatorname{rg}(g \circ f).
\end{array}
\]
于是得
\[
\operatorname{rg} f + \operatorname{rg} g = \dim_K V_1 + \dim_K V_2 - \dim_K(\ker f) - \dim_K(\ker g),
\]
再利用第三个公式得
\[
\operatorname{rg} f + \operatorname{rg} g = \dim_K(\ker(g \circ f)) + \operatorname{rg}(g \circ f) + \dim_K V_2 - \dim_K(\ker f) - \dim_K(\ker g).
\]
因此我们可以通过证明
\[
\dim_K(\ker(g \circ f)) \leq \dim_K(\ker f) + \dim_K(\ker g)
\]
来达到目标.为获得这个不等式的思路, 我们先弄清所涉及的向量空间.$\ker f$ 和 $\ker(g \circ f)$ 是 $V_1$ 的线性子空间, 且 $\ker f \subset \ker(g \circ f) \subset V_1$, 这容易验证.此外 $\ker g$ 是 $V_2$ 的线性子空间, 确切地说有 $f(\ker(g \circ f)) \subset \ker g \subset V_2$.这鼓励我们考察由 $f$ 诱导的线性映射
\[
\tilde{f} \colon \ker(g \circ f) \longrightarrow \ker g,
\]
其核与 $\ker f$ 一致, 并对它应用维数公式(定理~\mtref{thm:2.1.10}).这的确给出
\[
\dim_K(\ker(g \circ f)) = \dim_K(\ker \tilde{f}) + \dim_K(\operatorname{im} \tilde{f}) \leq \dim_K(\ker f) + \dim_K(\ker g),
\]
即所希望的,为达到目标还须证明的不等式.

上述途径看起来有点复杂, 我们自问是否也能用稍简单些的论证达到目标.我们论证的一个要点是, 把线性子空间 $\operatorname{im} \tilde{f} \subset V_2$ 换成 $\ker g$, 从而使用不等式 $\dim_K(\operatorname{im} \tilde{f}) \leq \dim_K(\ker g)$.其中更确切地有
\[
\operatorname{im} \tilde{f} = \ker g \cap f(V_1),
\]
这容易看出.这使我们想到, 不用给定的映射 $f$ 和 $g$, 而用 $K$-线性映射
\[
V_1 \xrightarrow{f'} f(V_1) \xrightarrow{g'} V_3,
\]
其中 $f'$ 和 $g'$ 分别由 $f$ 和 $g$ 通过把 $V_2$ 换成 $f(V_1)$ 的限制得到.这一做法事实上导致一个更简单的解, 我们马上就会看到.

\noindent\textbf{解题路径:} 我们于是把 $V_2$ 限制到 $f(V_1)$, 考虑所得的 $K$-线性映射 $V_1 \xrightarrow{f'} f(V_1) \xrightarrow{g'} V_3$.然后把维数公式(定理~\mtref{thm:2.1.10})应用于 $g'$, 得
\[
\begin{array}{r l}
& \operatorname{rg} f = \dim_K f(V_1) = \dim_K(\ker g') + \dim_K(\operatorname{im} g')\\
& \qquad = \dim_K\bigl(\ker g \cap f(V_1)\bigr) + \operatorname{rg}(g \circ f)\\
& \qquad \leq \dim_K(\ker g) + \operatorname{rg}(g \circ f).
\end{array}
\]
再利用按 $g$ 的维数公式给出的 $\dim_K(\ker g) = \dim_K V_2 - \operatorname{rg} g$, 即得所希望的不等式 $\operatorname{rg} f + \operatorname{rg} g \leq \operatorname{rg}(g \circ f) + \dim_K V_2$.

\subsection*{习题 2.2/4}

\noindent\textbf{起点:} 给定一个带子集 $A \subset V$ 的 $K$-向量空间.

\noindent\textbf{目标:} 要考虑满足 $A \subset B$ 的子集 $B \subset V$, 使得由
\[
x \sim_B y : \Longleftrightarrow x - y \in B \quad \text{对 } x, y \in V
\]
给出的关系 $\sim_B$ 是一个等价关系.设 $\mathfrak{B}$ 是所有这些子集 $B \subset V$ 的集合.要证明 $\mathfrak{B}$ 含有一个最小元素, 记为 $A'$, 即对每个另外的 $B \in \mathfrak{B}$ 都有 $A' \subset B$.还须在 $A'$ 对数乘封闭的情形, 即对 $\alpha \in K$, $a \in A$ 恒有 $\alpha a \in A$ 的情形, 更精确地确定 $A'$.

\noindent\textbf{策略:} 在定义--定理~\mtref{std:1.4.5} 中由子集 $A \subset V$ 生成的线性子空间 $A' = \langle A \rangle \subset V$ 的例子, 给出了一些基本做法的提示.一种可能的做法是, 把 $A'$ 取作对所有集合 $B \in \mathfrak{B}$ 的交.如我们下面将要看到的, 这一方法毫无问题地可行.但它不能给出所希望的集合 $A'$ 的具体描述.因此自然也要研究通过具体构造步骤把 $A$ 增大到所希望集合 $A'$ 的方法.为此, 必须先把等价关系 $\sim_B$ 的性质转化为相应子集 $B \subset V$ 的性质.

因此我们考虑任意子集 $B \subset V$ 以及相应的关系 $\sim_B$.下面这个联系是自然的:
\[
\left( \begin{array}{c}
0 \in B\\
a \in B \Longrightarrow -a \in B\\
a,\, b \in B \Longrightarrow a + b \in B
\end{array} \right)
\Longleftrightarrow
\binom{\text{关系 ``$\sim_B$''}}{\text{是一个等价关系}}
\]
事实上, 蕴含 ``$\Longrightarrow$'' 是平凡的, 因为列出的对 $B$ 的条件依次保证了关系 $\sim_B$ 的自反性,对称性和传递性.为看出余下的蕴含 ``$\Longleftarrow$'', 我们假定 $\sim_B$ 是一个等价关系.则对所有 $x \in V$ 有 $x \sim_B x$, 从而 $0 = x - x \in B$.对 $a \in B$ 又有 $a - 0 = a \in B$, 即 $a \sim_B 0$.由对称性得 $0 \sim_B a$, 从而 $-a = 0 - a \in B$.现在设 $a, b \in B$.则因 $a - 0 = a \in B$ 有 $a \sim_B 0$, 又因 $0 - (-b) = b \in B$ 有 $0 \sim_B (-b)$.传递性给出 $a \sim_B (-b)$, 即 $a + b = a - (-b) \in B$, 正如所断言.

注意到左边的条件按定义~\mtref{def:1.2.3} 恰刻画 $V$ 的加法群的一个子群, 我们于是认识到:

对子集 $B \subset V$, 关系 $\sim_B$ 在 $V$ 上定义等价关系, 当且仅当 $B$ 是 $V$ 的一个加法子群.

我们由此为我们的问题学到, 必须从 $A$ 过渡到包含 $A$ 的最小加法子群 $A' \subset V$.为此既可以把 $A'$ 定义为 $V$ 的所有把 $A$ 作为子集包含的加法子群之交, 也可以具体给出 $A'$, 类似我们在定义--定理~\mtref{std:1.4.5} 中解释由 $A$ 生成的线性子空间 $\langle A \rangle \subset V$ 那样.

\noindent\textbf{解题路径:} 如开头所说, 设 $\mathfrak{B}$ 是所有满足 $A \subset B$,使 $\sim_B$ 在 $V$ 上定义等价关系的子集 $B \subset V$ 的集合.显然 $V \in \mathfrak{B}$, 从而 $\mathfrak{B}$ 非空.于是我们设
\[
A' = \bigcap_{B \in \mathfrak{B}} B,
\]
并容易看出 $A' \in \mathfrak{B}$.事实上, 对向量 $x, y \in V$, 因 $A' = \bigcap_{B \in \mathfrak{B}} B$, 有 $x \sim_{A'} y$ 当且仅当对所有 $B \in \mathfrak{B}$ 有 $x \sim_B y$.特别地, 显然 $A'$ 是 $A$ 的使 $\sim_{A'}$ 在 $V$ 上给出等价关系的最小上集.

我们上面已经考虑过, 对子集 $B \subset V$, 关系 $\sim_B$ 恰在 $B$ 是 $V$ 的加法子群时于 $V$ 上定义等价关系.因此 $\mathfrak{B}$ 由 $V$ 的所有包含 $A$ 的加法子群组成, 而 $A' = \bigcap_{B \in \mathfrak{B}} B$ 是唯一确定的,包含 $A$ 的 $V$ 的最小子群, 换言之, 是 $A$ 在 $V$ 中生成的加法子群.后者还可以写成
\[
A' = \left\{\sum_{i = 1}^{r} \varepsilon_i a_i;\; r \in \mathbb{N},\ \varepsilon_i \in \{1, -1\},\ a_i \in A \right\},
\]
因为每个满足 $A \subset B$ 的加法子群 $B \subset V$ 必然包含右端, 并且不难看出我们把右端恰好定义为一个包含 $A$ 的 $V$ 的加法子群.现在若对 $\alpha \in K$ 和 $a \in A$ 恒有 $\alpha a \in A$, 则还可以写成
\[
A' = \left\{\sum_{i = 1}^{r} \alpha_i a_i;\; r \in \mathbb{N},\ \alpha_i \in K,\ a_i \in A \right\},
\]
并得到 $A' = \langle A \rangle$, 即在这种情形 $A'$ 恰是 $A$ 在 $V$ 中生成的线性子空间; 参见定义--定理~\mtref{std:1.4.5}.

\noindent\textbf{补充:} 在我们的解末使用生成群的显式描述看起来有些绕弯, 用一个更形式化的论证可以避免.事实上, 对数量 $\alpha \in K$, 考虑乘法
\[
V \longrightarrow V, \qquad x \longmapsto \alpha x,
\]
它保持子集的包含关系, 并把 $V$ 的加法子群映为加法子群.现在设 $A \subset V$ 是一个子集, $A'$ 是 $A$ 在 $V$ 中生成的加法子群, 即 $V$ 中所有包含 $A$ 的加法子群之交.则 $\alpha A'$ 是由 $\alpha A$ 生成的子群, 而在 $\alpha A = A$ 的情形也得到 $\alpha A' = A'$, 因为由 $A$ 生成的 $V$ 的加法子群是唯一确定的.因此若 $A$ 对数乘封闭, 则 $A'$ 亦然, 由定义~\mtref{def:1.4.2} 我们看到 $A'$ 是 $V$ 的线性子空间.于是 $A'$ 也是包含 $A$ 的 $V$ 的最小线性子空间, 即由 $A$ 生成的线性子空间 $\langle A \rangle \subset V$.

\subsection*{习题 2.2/8}

\noindent\textbf{起点:} 对 $K$-向量空间 $V$ 的线性子空间 $U, U'$, 考虑 $K$-线性映射
\[
f \colon U \longrightarrow U + U' \longrightarrow (U + U')/U',
\]
它由包含 $U \subset U + U'$ 与典范满同态 $U + U' \longrightarrow (U + U')/U'$ 复合而成.

\noindent\textbf{目标:} 要证明 $\ker f = U \cap U'$, 并进一步证明 $f$ 诱导同构
\[
\overline{f} \colon U/(U \cap U') \longrightarrow (U + U')/U'\text{.}
\]

\noindent\textbf{策略:} 要考虑 $K$-线性映射 $f \colon V \longrightarrow V'$, 其中 $V = U$ 且 $V' = (U + U')/U'$.问的是形如 $\overline{f} \colon V/\ker f \longrightarrow V'$ 的诱导同构, 这提示我们应用同态定理~\mtref{thm:2.2.8}.

\noindent\textbf{解题路径:} 我们知道, $U'$ 是典范满同态 $\pi \colon U + U' \longrightarrow (U + U')/U'$ 的核.此外 $f$ 是 $\pi$ 到 $U$ 上的限制.于是 $\ker f = U \cap \ker \pi = U \cap U'$.按同态定理(定理~\mtref{thm:2.2.8}), $f$ 于是诱导 $K$-线性映射 $\overline{f} \colon U/(U \cap U') \longrightarrow (U + U')/U'$, 且因 $\ker f = U \cap U'$ 它是单射.为看出 $\overline{f}$ 也满射, 从而是同构, 我们证明 $f$ 满射.设 $y \in (U + U')/U'$ 是一个剩余类, $x + x' \in U + U'$($x \in U$, $x' \in U'$)是 $y$ 的一个代表元, 即 $y = (x + x') + U'$.则也有 $y = x + U'$, 从而 $f(x) = y$.因此 $f$ 满射, 总体按推论~\mtref{cor:2.2.9} $\overline{f}$ 是同构.

\subsection*{习题 2.2/11}

\noindent\textbf{起点:} 没有描述具体的起点情形.

\noindent\textbf{目标:} 要构造一个带线性子空间 $U \neq 0$ 的 $K$-向量空间 $V$, 使得存在同构 $V/U \longrightarrow V$.此外还问, 在 $\dim_K V < \infty$ 的情形能否存在这样的例子.

\noindent\textbf{策略:} 我们力求给出尽可能简单的例子.因此先考察 $\dim_K V < \infty$ 的情形.若 $U \neq 0$ 是 $V$ 的线性子空间, 则 $\dim_K U > 0$, 从而按定理~\mtref{thm:2.2.7}
\[
\dim_K(V/U) = \dim_K V - \dim_K U < \dim_K V.
\]
我们由此读出, 在这种情形不存在同构 $V/U \longrightarrow V$, 因为同构保持向量空间的维数.

因此我们必须求助于无穷维的 $K$-向量空间 $V$.在最简单的情形, 这样的向量空间具有一个可数基.设 $V$ 是带可数基 $\{x_i;\; i \in \mathbb{N}\}$ 的 $K$-向量空间, 例如 $K$-向量空间
\[
V = \left\{(\alpha_j)_{j \in \mathbb{N}} \in K^{\mathbb{N}};\; \alpha_j = 0 \text{ 对几乎所有 } j \in \mathbb{N} \right\},
\]
带逐分量加法与数量乘法.这里单位向量 $x_i = (\delta_{ij})_{j \in \mathbb{N}}$, $i \in \mathbb{N}$, 构成一个基.

问的是对某个合适的线性子空间 $U \subset V$ 的 $K$-线性映射 $\overline{f} \colon V/U \longrightarrow V$.借助同态定理(定理~\mtref{thm:2.2.8}), $\overline{f}$ 可以由合适的 $K$-线性映射 $f \colon V \longrightarrow V$ 诱导.为定义这样的映射 $f$, 又可以应用定理~\mtref{thm:2.1.7}, 即应用于未必有限的生成系和基的版本.若再利用存在从 $\mathbb{N}$ 到 $\mathbb{N}$ 的满射但非单射的映射, 我们就备齐了构造所希望类型例子的所有要素.

\noindent\textbf{解题路径:} 设 $V$ 是带可数基 $\{x_i;\; i \in \mathbb{N}\}$ 的 $K$-向量空间.然后考虑由
\[
x_0 \longmapsto 0, \quad x_i \longmapsto x_{i - 1} \quad \text{对 } i > 0,
\]
即由
\[
\sum_{i = 0}^{\infty} \alpha_i x_i \longmapsto \sum_{i = 0}^{\infty} \alpha_{i + 1} x_i,
\]
以及线性扩张描述的 $K$-线性映射 $f \colon V \longrightarrow V$, 其中我们假定系数 $\alpha_i \in K$ 对几乎所有 $i \in \mathbb{N}$ 为零.显然 $f$ 满射, 且 $\ker f = K \cdot x_0$.于是设 $U = K \cdot x_0$, 则 $U \neq 0$, 且按同态定理(推论~\mtref{cor:2.2.9}) $f$ 如所希望地诱导同构 $\overline{f} \colon V/U \longrightarrow V$.在 $\dim_K V < \infty$ 的情形不可能存在这样的同构, 我们上面已经考虑过.

\subsection*{习题 2.3/2}

\noindent\textbf{起点:} 对 $n \in \mathbb{N}$, 把 $K^n$ 视为带投影
\[
p_i \colon K^n \longrightarrow K, \quad (\alpha_1, \ldots, \alpha_n) \longmapsto \alpha_i, \quad i = 1, \ldots, n,
\]
的 $K$-向量空间, 其中 $p_i$ 是 $K^n$ 上的线性形式.

\noindent\textbf{目标:} $p_1, \ldots, p_n$ 构成对偶空间 $(K^n)^*$ 的一个基, 即 $K^n$ 的,由单位向量 $e_j = (\delta_{1j}, \ldots, \delta_{nj})$, $j = 1, \ldots, n$, 组成的典范基的对偶基.此外, 对 $n = 4$ 和 $K = \mathbb{R}$, 还要确定 $(\mathbb{R}^4)^*$ 中基
\[
x_1 = (1, 0, 0, 0), \quad x_2 = (1, 1, 0, 0), \quad x_3 = (1, 1, 1, 0), \quad x_4 = (1, 1, 1, 1)
\]
的对偶基, 方法是把其元素表示为 $p_i$ 的线性组合.

\noindent\textbf{策略:} 对 $i, j = 1, \ldots, n$ 有 $p_i(e_j) = \delta_{ij}$, 因此按定理~\mtref{thm:2.3.5}, 线性形式 $p_1, \ldots, p_n \in (K^n)^*$ 构成 $e_1, \ldots, e_n \in K^n$ 的对偶基.在 $n = 4$, $K = \mathbb{R}$ 的情形容易验证, 向量 $x_1, \ldots, x_4 \in \mathbb{R}^4$ 构成一个线性无关组, 从而是 $\mathbb{R}^4$ 的一个基.按定理~\mtref{thm:2.3.5}, 于是在 $(\mathbb{R}^4)^*$ 中存在 $x_1, \ldots, x_4$ 的对偶基, 记为 $\varphi_1, \ldots, \varphi_4$.由于 $p_1, \ldots, p_4$ 也是 $(\mathbb{R}^4)^*$ 的一个基, 存在常数 $\lambda_{ij} \in \mathbb{R}$ 使
\[
\varphi_j = \sum_{i = 1}^{4} \lambda_{ij} p_i, \quad j = 1, \ldots, 4.
\]
然后要确定常数 $\lambda_{ij}$, 其中对偶基的刻画性质 $\varphi_j(x_k) = \delta_{jk}$ 给出一些方程, 我们可以尝试用它们来确定常数 $\lambda_{ij}$.

\noindent\textbf{解题路径:} 我们已经说明, $p_1, \ldots, p_n \in (\mathbb{R}^n)^*$ 是典范基 $e_1, \ldots, e_n \in \mathbb{R}^n$ 的对偶基, 并且对 $n = 4$, $K = \mathbb{R}$, 向量 $x_1, \ldots, x_4$ 也构成 $\mathbb{R}^4$ 的一个基, 因而可以考虑其相应的对偶基 $\varphi_1, \ldots, \varphi_4 \in (\mathbb{R}^4)^*$.于是有表示
\[
\varphi_j = \sum_{i = 1}^{4} \lambda_{ij} p_i, \quad j = 1, \ldots, 4,
\]
其系数 $\lambda_{ij} \in \mathbb{R}$ 待定.定理~\mtref{thm:2.3.5} 中对偶基的刻画性质 $\varphi_j(x_k) = \delta_{jk}$ 于是给出方程
\[
\sum_{i = 1}^{4} \lambda_{ij} p_i(x_k) = \delta_{jk}, \quad j, k = 1, \ldots, 4,
\]
即按 $j = 1, \ldots, 4$ 排列的下列方程组:
\[
\begin{array}{rcl}
\lambda_{11} \cdot 1 + \lambda_{21} \cdot 0 + \lambda_{31} \cdot 0 + \lambda_{41} \cdot 0 &=& 1\\
\lambda_{11} \cdot 1 + \lambda_{21} \cdot 1 + \lambda_{31} \cdot 0 + \lambda_{41} \cdot 0 &=& 0\\
\lambda_{11} \cdot 1 + \lambda_{21} \cdot 1 + \lambda_{31} \cdot 1 + \lambda_{41} \cdot 0 &=& 0\\
\lambda_{11} \cdot 1 + \lambda_{21} \cdot 1 + \lambda_{31} \cdot 1 + \lambda_{41} \cdot 1 &=& 0
\end{array}
\]
\[
\begin{array}{rcl}
\lambda_{12} \cdot 1 + \lambda_{22} \cdot 0 + \lambda_{32} \cdot 0 + \lambda_{42} \cdot 0 &=& 0\\
\lambda_{12} \cdot 1 + \lambda_{22} \cdot 1 + \lambda_{32} \cdot 0 + \lambda_{42} \cdot 0 &=& 1\\
\lambda_{12} \cdot 1 + \lambda_{22} \cdot 1 + \lambda_{32} \cdot 1 + \lambda_{42} \cdot 0 &=& 0\\
\lambda_{12} \cdot 1 + \lambda_{22} \cdot 1 + \lambda_{32} \cdot 1 + \lambda_{42} \cdot 1 &=& 0
\end{array}
\]
\[
\begin{array}{rcl}
\lambda_{13} \cdot 1 + \lambda_{23} \cdot 0 + \lambda_{33} \cdot 0 + \lambda_{43} \cdot 0 &=& 0\\
\lambda_{13} \cdot 1 + \lambda_{23} \cdot 1 + \lambda_{33} \cdot 0 + \lambda_{43} \cdot 0 &=& 0\\
\lambda_{13} \cdot 1 + \lambda_{23} \cdot 1 + \lambda_{33} \cdot 1 + \lambda_{43} \cdot 0 &=& 1\\
\lambda_{13} \cdot 1 + \lambda_{23} \cdot 1 + \lambda_{33} \cdot 1 + \lambda_{43} \cdot 1 &=& 0
\end{array}
\]
\[
\begin{array}{rcl}
\lambda_{14} \cdot 1 + \lambda_{24} \cdot 0 + \lambda_{34} \cdot 0 + \lambda_{44} \cdot 0 &=& 0\\
\lambda_{14} \cdot 1 + \lambda_{24} \cdot 1 + \lambda_{34} \cdot 0 + \lambda_{44} \cdot 0 &=& 0\\
\lambda_{14} \cdot 1 + \lambda_{24} \cdot 1 + \lambda_{34} \cdot 1 + \lambda_{44} \cdot 0 &=& 0\\
\lambda_{14} \cdot 1 + \lambda_{24} \cdot 1 + \lambda_{34} \cdot 1 + \lambda_{44} \cdot 1 &=& 1
\end{array}
\]
由此得
\[
\begin{array}{l l l l}
\lambda_{11} = 1, & \lambda_{12} = 0, & \lambda_{13} = 0, & \lambda_{14} = 0\\
\lambda_{21} = -1, & \lambda_{22} = 1, & \lambda_{23} = 0, & \lambda_{24} = 0\\
\lambda_{31} = 0, & \lambda_{32} = -1, & \lambda_{33} = 1, & \lambda_{34} = 0\\
\lambda_{41} = 0, & \lambda_{42} = 0, & \lambda_{43} = -1, & \lambda_{44} = 1
\end{array}
\]
从而
\[
\varphi_1 = p_1 - p_2, \qquad \varphi_2 = p_2 - p_3, \qquad \varphi_3 = p_3 - p_4, \qquad \varphi_4 = p_4.
\]
这些就是所希望的 $p_1, \ldots, p_4$ 的线性组合, 它们描述 $(\mathbb{R}^4)^*$ 中 $x_1, \ldots, x_4$ 的对偶基.

\subsection*{习题 2.3/4}

\noindent\textbf{起点:} 对 $K$-向量空间 $V$ 的一个元素 $a$, 考虑由 $\varphi \longmapsto \varphi(a)$ 给出的线性映射 $a^* \colon V^* \longrightarrow K$.则 $a^*$ 是对偶空间 $V^*$ 上的一个线性形式, 即\emph{双重对偶空间}(\textup{doppelt dualer Vektorraum})\index{shuangchongduioukongjian@双重对偶空间(doppelt dualer Vektorraum)} $V^{**}$ 的一个元素.此外, 对 $a \in V$ 还可以考虑映射 $[a] \colon K \longrightarrow V$, $\alpha \longmapsto \alpha a$, 可以说是``乘以 $a$''.若把 $K$ 视为向量空间 $K^1$, 即视为它自身的向量空间, 则它是 $K$-线性的.

\noindent\textbf{目标:} 首先要证明映射 $\iota_V \colon V \longrightarrow \operatorname{Hom}_K(K, V)$, $a \longmapsto [a]$, 是 $K$-向量空间的同构.然后选取元素 $a \in V$, 考虑相应的 $K$-线性映射 $\iota_V(a)$, 即 $[a] \colon K \longrightarrow V$, 以及它的对偶映射 $[a]^* \colon V^* \longrightarrow K^*$, 其中 $K^* = \operatorname{Hom}_K(K, K)$ 指 $K$ 作为它自身向量空间的对偶空间.要证明: 当把 $K$ 作为 $K$-向量空间通过 $\iota_K \colon K \longrightarrow \operatorname{Hom}_K(K, K)$ 与它的对偶空间 $\operatorname{Hom}_K(K, K)$ 等同起来时, $a^*$ 恰是 $[a]$ 的对偶映射 $[a]^*$.

\noindent\textbf{策略:} $K$ 作为它自身上的向量空间具有一个典范基, 即由单位元 $1 \in K$ 组成.因此按定理~\mtref{thm:2.1.7}, $[a] \colon K \longrightarrow V$ 由对应 $1 \longmapsto a$ 唯一确定, 并且显然 $\iota_V \colon V \longrightarrow \operatorname{Hom}_K(K, V)$ 是一个同构.此外, 与 $[a]$ 对偶的线性映射由
\[
[a]^* \colon V^* \longrightarrow \operatorname{Hom}_K(K, K), \qquad \varphi \longmapsto \varphi \circ [a],
\]
描述, 其中 $\varphi \circ [a]$ 先把元素 $1 \in K$ 映到 $a$, 再映到 $\varphi(a)$.因此, 若把同构 $\iota_K \colon K \longrightarrow \operatorname{Hom}_K(K, K)$ 视为等同, 则 $\varphi \circ [a] = \varphi(a)$.这里的策略性考虑因而已经以简略形式给出了解.

\noindent\textbf{解题路径:} 映射 $\iota_V \colon V \longrightarrow \operatorname{Hom}_K(K, V)$, $a \longmapsto [a]$, 是 $K$-线性的, 因为显然对 $a, b \in V$ 和 $\alpha \in K$ 有 $[a + b] = [a] + [b]$ 以及 $[\alpha a] = \alpha [a]$.同样, $\kappa_V \colon \operatorname{Hom}_K(K, V) \longrightarrow V$, $\sigma \longmapsto \sigma(1)$, 是 $K$-线性映射, 且有 $\kappa_V \circ \iota_V = \operatorname{id}_V$ 以及 $\iota_V \circ \kappa_V = \operatorname{id}_{\operatorname{Hom}_K(K, V)}$.由此推出 $\iota_V$ 和 $\kappa_V$ 互为逆映射, 从而是同构.

现在对元素 $a \in V$, 考虑相应的线性映射 $[a] \colon K \longrightarrow V$, 以及它的对偶映射
\[
[a]^* \colon V^* \longrightarrow \operatorname{Hom}_K(K, K), \quad \varphi \longmapsto \varphi \circ [a].
\]
则对所有 $\varphi \in V^*$ 有 $(\varphi \circ [a])(1) = \varphi(a)$, 而 $\varphi \circ [a]$ 在 $\kappa_K \colon \operatorname{Hom}_K(K, K) \longrightarrow K$ 下, 即 $\iota_K$ 的逆下, 被等同为 $\varphi(a)$, 这正是要证明的.

\subsection*{习题 2.3/6}

\noindent\textbf{起点:} 要考虑 $K$-向量空间之间的 $K$-线性映射 $f \colon V \longrightarrow V'$, 它的对偶映射 $f^* \colon V'^* \longrightarrow V^*$, 以及 $f^*$ 的所谓余核, 即 coker $f^* = V^*/\operatorname{im} f^*$.

\noindent\textbf{目标:} 要构造典范同构 coker $f^* \longrightarrow (\ker f)^*$.

\noindent\textbf{策略:} 目标尤其要求考察线性子空间 $\ker f \subset V$ 以及相应的对偶空间 $(\ker f)^*$.这把我们引向考察典范正合序列
\begin{equation}
0 \longrightarrow \ker f \longrightarrow V \xrightarrow{f} V', \tag{(*)}
\end{equation}
以及它的对偶序列
\begin{equation}
V'^* \xrightarrow{f^*} V^* \longrightarrow (\ker f)^* \longrightarrow 0, \tag{(**)}
\end{equation}
后者按定理~\mtref{thm:2.3.4} 也正合.为把 $(\ker f)^*$ 与 coker $f^*$ 联系起来, 自然想到同态定理~\mtref{thm:2.2.8}.

\noindent\textbf{解题路径:} 我们从正合序列 (*) 出发, 考虑它的对偶正合序列 (**).则因 (**) 的正合性, $V^* \longrightarrow (\ker f)^*$ 是满射 $K$-线性映射, 其核与 $\operatorname{im} f^*$ 一致.同态定理(定理~\mtref{thm:2.2.8})于是给出所希望的典范同构 coker $f^* = V^*/\operatorname{im} f^* \longrightarrow (\ker f)^*$.

\section*{对应 3.1-3.5: 矩阵}

\subsection*{习题 3.1/4}

\noindent\textbf{起点:} 给定有限维 $K$-向量空间之间的一个 $K$-线性映射
$f \colon V \longrightarrow W$.设 $r = \operatorname{rg} f = \dim_K f(V)$ 为 $f$ 的秩.

\noindent\textbf{目标:} 要构造 $V$ 的基 $X$ 与 $W$ 的基 $Y$, 使得属于 $f$ 的矩阵
$A_{f, X, Y}$ 具有形状
\[
\begin{pmatrix} E_r & 0 \\ 0 & 0 \end{pmatrix},
\]
其中 $E_r$ 是 $(r \times r)$-单位矩阵.换言之, 要找到 $V$ 的基
$X = (x_1, \ldots, x_n)$ 与 $W$ 的基 $Y = (y_1, \ldots, y_m)$, 使得
\[
f(x_j) = \begin{cases} y_j, & \text{对 } 1 \leq j \leq r, \\ 0, & \text{对 } r < j \leq n\text{.} \end{cases}
\]
特别地, 这包含了不等式 $r \leq m$ 与 $r \leq n$.

\noindent\textbf{策略:} 每个 $K$-线性映射 $f \colon V \longrightarrow W$ 都确定了 $V$ 与
$W$ 中的某些线性子空间, 即 $\ker f \subset V$ 与 $\operatorname{im} f \subset W$.若存在
具有所希望性质的基 $X = (x_1, \ldots, x_n)$ 与 $Y = (y_1, \ldots, y_m)$, 则显然有
\[
\langle y_1, \ldots, y_r \rangle = \langle f(x_1), \ldots, f(x_r) \rangle = \operatorname{im} f,
\]
因为像 $f(x_{r+1}), \ldots, f(x_n)$ 都消失, 即 $x_{r+1}, \ldots, x_n \in \ker f$.由此我们
读出: $y_1, \ldots, y_r$ 构成 $\operatorname{im} f$ 的一个基.进一步我们可以证明,
$x_{r+1}, \ldots, x_n$ 构成 $\ker f \subset V$ 的一个基.这个元素组一方面线性无关,
另一方面又是 $\ker f$ 的一个生成系.事实上, 设 $x \in \ker f$, 并设
\[
x = \sum_{j = 1}^{n} \alpha_j x_j
\]
是 $x$ 关于基 $X$ 的一个线性组合表示.于是
\[
0 = f(x) = \sum_{j = 1}^{n} \alpha_j f(x_j) = \sum_{j = 1}^{r} \alpha_j f(x_j) = \sum_{j = 1}^{r} \alpha_j y_j,
\]
从而对 $j = 1, \ldots, r$ 有 $\alpha_j = 0$.因此得到
\[
x = \sum_{j = r + 1}^{n} \alpha_j x_j,
\]
并把 $x_{r+1}, \ldots, x_n$ 认作 $\ker f$ 的一个生成系, 即基.

从上述思考可以看出, 要构造一个"穿过" $\ker f$ 的 $V$ 的基 $X$, 以及一个"穿过"
$\operatorname{im} f$ 的 $W$ 的基 $Y$.此外还要实现: $X$ 中不属于 $\ker f$ 的所有元素都被
映到 $Y$ 的元素.在基的理论中可能有用的普遍构造原则, 特别有定理~\mtref{thm:1.5.6}(有限生成
$K$-向量空间中基的存在性)与定理~\mtref{thm:1.5.8}(基扩充定理).此外, 我们由
推论~\mtref{cor:1.5.14} (i) 推断, 有限生成 $K$-向量空间的每个线性子空间都具有一个(有限)基.

\noindent\textbf{解题路径:} 我们可以类似于定理~\mtref{thm:2.1.10} 的证明来处理.具体地说,
我们选取 $\operatorname{im} f$ 的一个基 $y_1, \ldots, y_r$, 并按照定理~\mtref{thm:1.5.8}
用元素 $y_{r+1}, \ldots, y_m$ 把它扩充为 $W$ 的一个基 $Y$.然后选取元素 $y_1, \ldots, y_r$
的原像 $x_1, \ldots, x_r \in V$.再设 $x_{r+1}, \ldots, x_n$ 是 $\ker f$ 的一个基.我们断言,
此时 $X = (x_1, \ldots, x_n)$ 是 $V$ 的一个基.首先想检验 $X$ 生成向量空间 $V$.对
$x \in V$, 由于 $y_1, \ldots, y_r$ 构成 $\operatorname{im} f$ 的一个生成系, 存在
$\alpha_1, \ldots, \alpha_r \in K$ 使得
\[
f(x) = \sum_{j = 1}^{r} \alpha_j y_j = \sum_{j = 1}^{r} \alpha_j f(x_j)\text{.}
\]
由此得 $x - \sum_{j = 1}^{r} \alpha_j x_j \in \ker f$, 于是存在系数
$\alpha_{r+1}, \ldots, \alpha_n \in K$ 使得
\[
x - \sum_{j = 1}^{r} \alpha_j x_j = \sum_{j = r + 1}^{n} \alpha_j x_j,
\]
即
\[
x = \sum_{j = 1}^{n} \alpha_j x_j\text{.}
\]
所以 $x_1, \ldots, x_n$ 构成 $V$ 的一个生成系.这个组线性无关也不难验证.若对某些系数
$\alpha_j \in K$ 有 $\sum_{j = 1}^{n} \alpha_j x_j = 0$, 则由方程
\[
\sum_{j = 1}^{r} \alpha_j y_j = f\left(\sum_{j = 1}^{n} \alpha_j x_j\right) = 0
\]
看出系数 $\alpha_1, \ldots, \alpha_r$ 消失.于是得到 $\sum_{j = r + 1}^{n} \alpha_j x_j = 0$,
并由于组 $x_{r+1}, \ldots, x_n$ 线性无关而推出 $\alpha_{r+1} = \cdots = \alpha_n = 0$.

综上, 我们把 $X$ 认作 $V$ 的一个基, 并确认矩阵 $A_{f, X, Y}$ 具有所要求的形状.

\noindent\textbf{补充:} 另一种解题路径(根据我们的策略性思考也许显得更自然)如下: 选取 $V$ 的
一个"穿过" $\ker f$ 的基 $X = (x_1, \ldots, x_n)$.这里我们把元素 $x_1, \ldots, x_n$ 编号,
使得 $\ker f$ 的基元素位于末尾, 即对某个指标 $s \geq 0$, 元素 $x_{s+1}, \ldots, x_n$ 构成
$\ker f$ 的一个基.然后对 $i = 1, \ldots, s$ 令 $y_i = f(x_i)$, 并断言 $y_i$ 构成
$\operatorname{im} f$ 的一个基.自然, $y_1, \ldots, y_s$ 生成线性子空间
$\operatorname{im} f \subset W$; 但它们也构成一个线性无关组, 因为由限制 $f$ 得到的 $K$-线性映射
$\langle x_1, \ldots, x_s \rangle \longrightarrow \operatorname{im} f$ 先是满射, 又由于
$\langle x_1, \ldots, x_s \rangle \cap \ker f = 0$ 也是单射, 因而是 $K$-向量空间的一个同构.
特别地有 $s = \operatorname{rg} f$.若现在把 $y_1, \ldots, y_s$ 扩充为 $W$ 的一个基, 则矩阵
$A_{f, X, Y}$ 具有所要求的形状.

\subsection*{习题 3.1/7}

\noindent\textbf{起点:} 要考虑一个 $K$-向量空间 $V \neq 0$ 及其 $K$-线性映射
$f \colon V \longrightarrow V$, 后者关于一个基 $X$ 由显式给出的矩阵
$A_{f, X, X} \in K^{n \times n}$ 描述.特别地, $X$ 是有限的, 设
$X = (x_1, \ldots, x_n)$.从矩阵 $A_{f, X, X}$ 我们读出: 对 $j = 1, \ldots, n$ 有
\[
f(x_j) = \sum_{i < j} x_i\text{.}
\]

\noindent\textbf{目标:} 要确定 $\dim_K(\ker f)$.

\noindent\textbf{策略:} 对于显式给出的线性映射 $f \colon V \longrightarrow V$, 原则上
有两种方法确定核的维数.其一, 可以直接把核确定为 $V$ 的线性子空间; 一般这意味着解一个
线性方程组.随后还必须给出核的一个基, 才能把其维数读作该基的长度.其二, 也可以利用维数公式
定理~\mtref{thm:2.1.10}.此时只需计算 $f$ 的秩, 即 $\dim(\operatorname{im} f)$.

在我们的情形下马上可以看出, 向量 $f(x_2), \ldots, f(x_n)$ 生成线性子空间
$\langle x_1, \ldots, x_{n-1} \rangle \subset V$, 因而按照大定理~\mtref{bigthm:1.5.9}
在 $V$ 中构成一个线性无关组.由于 $f(x_1) = 0$, 得 $\operatorname{rg} f = n - 1$,
并由维数公式得 $\dim_K(\ker f) = 1$.另一方面, 在这种情形下人们也"看到" $\ker f$ 由
$x_1$ 生成, 由此同样读出 $\dim_K(\ker f) = 1$.

\noindent\textbf{解题路径:} 像 $\operatorname{im} f$ 由于 $f(x_1) = 0$ 而由向量
$f(x_j)$($j = 2, \ldots, n$)生成.我们想对 $n \geq 1$ 证明
\[
\operatorname{im} f = \langle f(x_2), \ldots, f(x_n) \rangle = \langle x_1, \ldots, x_{n-1} \rangle\text{.}
\]
自然有 $\langle f(x_2), \ldots, f(x_n) \rangle \subset \langle x_1, \ldots, x_{n-1} \rangle$.
为证明反向包含, 我们利用
\[
f(x_j) = \sum_{i < j} x_i = f(x_{j-1}) + x_{j-1} \quad \text{对 } j = 2, \ldots, n,
\]
即 $x_{j-1} = f(x_j) - f(x_{j-1}) \in \langle f(x_2), \ldots, f(x_n) \rangle$, 这就给出所求的
等式.于是由大定理~\mtref{bigthm:1.5.9} 得 $\operatorname{rg} f = \dim_K(\operatorname{im} f) = n - 1$,
而维数公式定理~\mtref{thm:2.1.10} 给出 $\dim_K(\ker f) = 1$.

\noindent\textbf{补充:} 在我们的情形下, $f$ 的核也可以容易地确定.首先有 $x_1 \in \ker f$.
若再使用 $\dim_K(\ker f) = 1$, 则已经得到 $\ker f = \langle x_1 \rangle$.但我们也可以不用
维数公式来证明这一点.设 $x = \sum_{i = 1}^{n} \alpha_i x_i \in \ker f$, 即满足 $f(x) = 0$
的一个向量.由于 $f(x_1) = 0$, 由此得 $\sum_{i = 2}^{n} \alpha_i f(x_i) = 0$, 从而
$\alpha_2 = \cdots = \alpha_n = 0$, 因为 $f(x_2), \ldots, f(x_n)$ 构成一个线性无关组.
于是得 $x \in \langle x_1 \rangle$, 因而 $\ker f = \langle x_1 \rangle$.对于更复杂的矩阵
$A_{f, X, X}$, $\ker f$ 的显式确定将费力得多.

\subsection*{习题 3.1/8}

\noindent\textbf{起点:} 要考虑一个矩阵 $N = (\alpha_{ij})_{i, j = 1, \ldots, m} \in K^{m \times m}$,
$m \geq 1$, 其系数 $\alpha_{ij} \in K$ 满足: 对 $i \geq j$ 有 $\alpha_{ij} = 0$.因此 $N$
是一个上三角矩阵, 其对角元素 $\alpha_{ii}$ 全部消失, 对角线下方的所有元素
$\alpha_{ij}$(即 $i > j$)亦然.只有对角线以上的元素 $\alpha_{ij}$(即 $i < j$)才可能出现
非平凡的取值.设 $E \in K^{m \times m}$ 是单位矩阵.

\noindent\textbf{目标:} 首先要证明 $N^m = 0$.在进一步的一步中, 应以几何级数为范本证明矩阵
$B = E + N$ 可逆, 即存在矩阵 $C \in K^{m \times m}$ 使得 $BC = CB = E$.

\noindent\textbf{策略:} 先考察平方 $N^2$, 就会发现第一条上副对角线变成平凡的, 即在位置
$(i, i+1)$($i = 1, \ldots, m-1$)上 $N^2$ 的所有元素都是 $0$, 这条副对角线以下的所有元素
亦然.再与 $N$ 相乘则表明, 在 $N^3$ 中第二条上副对角线的所有元素, 即在位置
$(i, i+2)$($i = 1, \ldots, m-2$)上的元素, 也变成平凡的.上副对角线总共有 $m - 1$ 条,
其中最后一条只由一个位于位置 $(1, m)$ 的元素组成.因此可以推测, 事实上确有 $N^m = 0$.
在这些思考中重要的是如下平凡的观察: 一个行(元素为 $\alpha_1, \ldots, \alpha_m$)与一个列
(元素为 $\beta_1, \ldots, \beta_m$)的乘积 $\sum_{i = 1}^{m} \alpha_i \beta_i$ 是平凡的,
只要存在一个指标 $s$, 使得对 $i \leq s$ 有 $\alpha_i = 0$, 且对 $i > s$ 有
$\beta_i = 0$.

一旦我们认定了 $N^m = 0$, 自然就对一切指数 $\mu \geq m$ 有 $N^{\mu} = 0$.因此在几何级数的
公式
\[
(1 - q) \cdot \sum_{\mu = 0}^{\infty} q^{\mu} = 1
\]
中把参数 $q$ 换成 $N$, 把单位元 $1$ 换成单位矩阵 $E$, 则方程的两边都有定义, 于是可以检验
等式是否保持.

\noindent\textbf{解题路径:} 我们考虑满足 $\alpha_{ij} = 0$($i \geq j$)的矩阵
$N = (\alpha_{ij})_{i, j = 1, \ldots, m}$, 以及另一个矩阵
$P = (\beta_{jk})_{j, k = 1, \ldots, m} \in K^{m \times m}$, 并带有参数 $s \in \mathbb{N}$,
使得对 $j \geq k - s$ 有 $\beta_{jk} = 0$.然后形成乘积 $N \cdot P = (\gamma_{ik})_{i, k = 1, \ldots, m}$,
\[
\gamma_{ik} = \sum_{j = 1}^{m} \alpha_{ij} \beta_{jk},
\]
并断言: 对所有满足 $i \geq k - s - 1$ 的 $i, k$ 有 $\gamma_{ik} = 0$.为看出这一点,
选取 $i, k \in \{1, \ldots, m\}$, 考察 $N$ 的第 $i$ 行的元素 $\alpha_{i1}, \ldots, \alpha_{im}$
以及 $P$ 的第 $k$ 列的元素 $\beta_{1k}, \ldots, \beta_{mk}$.其中 $\alpha_{i1}, \ldots, \alpha_{ii}$
与 $\beta_{\max(1, k-s), k}, \ldots, \beta_{mk}$ 基于我们对 $N$ 与 $P$ 的假定而消失.于是只要
$i \geq k - s - 1$ 成立, 乘积矩阵 $N \cdot P$ 的元素
\[
\gamma_{ik} = \sum_{j = 1}^{m} \alpha_{ij} \beta_{jk}
\]
也消失, 因为在各个被加项中两个因子中总有一个是 $0$.这就证实了我们的断言.

现在特别令 $P = N^{\mu} = (\beta_{jk}^{(\mu)})_{j, k = 1, \ldots, m}$ 带有指数 $\mu \geq 1$,
则由归纳得对 $j \ge k - \mu + 1$ 有 $\beta_{jk}^{(\mu)} = 0$, 从而尤其有 $N^m = 0$.此外,
\[
(E - N) \cdot \sum_{\mu = 0}^{m} N^{\mu} = \sum_{\mu = 0}^{m} N^{\mu} - \sum_{\mu = 1}^{m + 1} N^{\mu} = \sum_{\mu = 0}^{m} N^{\mu} - \sum_{\mu = 1}^{m} N^{\mu} = E,
\]
而同样的计算可以对乘积 $(\sum_{\mu = 0}^{m} N^{\mu}) \cdot (E - N)$ 施行.因此 $E - N$
在 $K^{m \times m}$ 中可逆, 并且把 $N$ 换成 $-N$ 后, 我们可以类似地对 $E + N$ 论证.

\noindent\textbf{补充:} 另一种方式下, 等式 $N^m = 0$ 也可以不经过计算就看出来.为此考察由 $N$
定义的线性映射
\[
f \colon K^m \longrightarrow K^m, \qquad x \longmapsto N \cdot x,
\]
以及线性子空间的链 $0 = U_0 \subset U_1 \subset \cdots \subset U_m = K^m$, 其中
$U_i = \langle e_1, \ldots, e_i \rangle$; 这里 $e_1, \ldots, e_m$ 是 $K^m$ 的标准基.矩阵
$N = (\alpha_{ij})_{i, j = 1, \ldots, m}$ 就是 $f$ 关于这一基的描述矩阵, 而元素
$\alpha_{ij}$($i \geq j$)的消失恰好意味着对 $i = 1, \ldots, m$ 有 $f(U_i) \subset U_{i-1}$.
由此容易用归纳看出, 对 $\mu = 0, \ldots, m$ 有 $f^{\mu}(U_m) \subset U_{m - \mu}$, 特别地
$f^m(U_m) = 0$, 因而 $f^m = 0$.

最后我们还想证明, $K$-线性映射 $\operatorname{id} + f \colon K^m \longrightarrow K^m$ 是单射,
从而按照推论~\mtref{cor:2.1.11} 是一个自同构.事实上, 设 $x \in \ker(\operatorname{id} + f)$.
则 $f(x) = -x$, 应用 $f$ 得 $f^2(x) = -f(x)$, 即 $f^2(x) = x$.用完全归纳法推断
$f^m(x) = (-1)^m x$, 从而由于 $f^m = 0$ 得 $x = 0$.因此 $\operatorname{id} + f$ 是一个
同构.按照注记~\mtref{rem:2.1.2}, 相应的逆映射 $g \colon K^m \longrightarrow K^m$ 满足
$g \circ (\operatorname{id} + f) = \operatorname{id} = (\operatorname{id} + f) \circ g$.
用定理~\mtref{thm:3.1.9} 把它翻译到矩阵的语言中, 就意味着 $(E + N)$ 可逆.

\subsection*{习题 3.2/2}

\noindent\textbf{起点:} 要考虑域 $K = \mathbb{Q}$ 或 $K = \mathbb{F}_5$ 上的一个具体给出的矩阵
$A \in K^{3 \times 4}$.其中 $\mathbb{F}_5$ 是含 5 个元素的域, 即由零元素 $0$,单位元 $1$
以及元素
\[
2 \cdot 1 = 1 + 1, \qquad 3 \cdot 1 = 1 + 1 + 1, \qquad 4 \cdot 1 = 1 + 1 + 1 + 1
\]
组成的域.对 $n = 0, 1, 2, 3, 4$, 我们又写 $n$ 来代替 $n \cdot 1$.一般地有 $n \cdot 1 = r$,
其中 $r$ 是 $n$ 除以 5 的(非负)余数.

\noindent\textbf{目标:} 要计算 $A$ 的行秩, 分别对 $K = \mathbb{Q}$ 与 $K = \mathbb{F}_5$.

\noindent\textbf{策略:} 宜用大定理~\mtref{bigthm:3.2.4} 的高斯消元法.

\noindent\textbf{解题路径:} 我们把 $A$ 化为行阶梯形, 首先对域 $K = \mathbb{Q}$:
\[
A = \begin{pmatrix} 1 & 1 & 3 & 2 \\ 1 & 1 & 2 & 3 \\ 1 & 1 & 0 & 0 \end{pmatrix}
\longmapsto \begin{pmatrix} 1 & 1 & 3 & 2 \\ 0 & 0 & -1 & 1 \\ 0 & 0 & -3 & -2 \end{pmatrix}
\longmapsto \begin{pmatrix} 1 & 1 & 3 & 2 \\ 0 & 0 & -1 & 1 \\ 0 & 0 & 0 & -(3 + 2) \end{pmatrix}
\]
对 $K = \mathbb{Q}$ 由此读出 $\operatorname{rg} A = 3$.在 $K = \mathbb{F}_5$ 的情形, 我们可以
施行同样的计算, 其中负号分别表示相应的负元素, 即关于加法的逆元素.但在 $\mathbb{F}_5$ 中
有 $3 + 2 = 0$, 从而在 $K = \mathbb{F}_5$ 的情形 $\operatorname{rg} A = 2$.

\subsection*{习题 3.2/5}

\noindent\textbf{起点:} 要考虑 $\mathbb{R}^5$ 的两个线性子空间
$U = \langle a_1, a_2, a_3 \rangle$ 与 $U' = \langle a_1', a_2', a_3' \rangle$, 其中向量
$a_1, a_2, a_3$ 与 $a_1', a_2', a_3'$ 具体给出.把它们看作列向量, 即为矩阵
\[
A = (a_1, a_2, a_3) = \begin{pmatrix} 1 & 2 & 0 \\ 0 & 3 & 3 \\ 1 & 4 & 2 \\ 0 & 1 & 1 \\ 1 & 2 & 0 \end{pmatrix},
\qquad
A' = (a_1', a_2', a_3') = \begin{pmatrix} 1 & 1 & 1 \\ -1 & 3 & 2 \\ 1 & 3 & 3 \\ 0 & 1 & 1 \\ 2 & 1 & 2 \end{pmatrix}\text{.}
\]

\noindent\textbf{目标:} 要计算 $\dim_{\mathbb{R}} U$ 与 $\dim_{\mathbb{R}} U'$, 并证明
$U \subset U'$.

\noindent\textbf{策略:} 结果定理~\mtref{thm:3.2.11} 给我们提供了一种可能性: 当已知一个(有限的)
生成系时, 计算 $\mathbb{R}^5$ 中线性子空间的维数.因此我们应该把向量 $a_i$ 与 $a_i'$ 视为
列向量, 并由此构造矩阵, 例如
\[
A = (a_1, a_2, a_3), \qquad A' = (a_1', a_2', a_3') \in \mathbb{R}^{5 \times 3}\text{.}
\]
利用定理~\mtref{thm:3.2.11} 可以计算 $A$ 与 $A'$ 的秩, 从而计算 $U$ 与 $U'$ 的维数.为证明
$U \subset U'$, 我们应该思考: 此处是否也能用一个维数或秩的论证达到目的.例如我们可用
推论~\mtref{cor:1.5.14} (ii): 若线性子空间的维数一致且有限, 则由一个线性子空间的包含可以推出
它们相等.这使我们想到, $U \subset U'$ 等价于 $U' = U + U'$, 从而由于 $U' \subset U + U'$
也等价于 $\dim U' = \dim_K(U + U')$.后一关系可以用 $\operatorname{rg}(A') = \operatorname{rg}(A, A')$
来检验, 其中 $(A, A')$ 表示由列 $a_1, a_2, a_3, a_1', a_2', a_3'$ 构造的矩阵.

\noindent\textbf{解题路径:} 我们首先确定由列向量 $a_1', a_2', a_3', a_1, a_2, a_3$ 构造的矩阵
$(A', A)$ 的一个行阶梯形:
\[
\begin{pmatrix} 1 & 1 & 1 & 1 & 2 & 0 \\ -1 & 3 & 2 & 0 & 3 & 3 \\ 1 & 3 & 3 & 1 & 4 & 2 \\ 0 & 1 & 1 & 0 & 1 & 1 \\ 2 & 1 & 2 & 1 & 2 & 0 \end{pmatrix}
\longmapsto \begin{pmatrix} 1 & 1 & 1 & 1 & 2 & 0 \\ 0 & 4 & 3 & 1 & 5 & 3 \\ 0 & 2 & 2 & 0 & 2 & 2 \\ 0 & 1 & 1 & 0 & 1 & 1 \\ 0 & -1 & 0 & -1 & -2 & 0 \end{pmatrix}
\]
\[
\longmapsto \begin{pmatrix} 1 & 1 & 1 & 1 & 2 & 0 \\ 0 & 1 & 1 & 0 & 1 & 1 \\ 0 & 4 & 3 & 1 & 5 & 3 \\ 0 & 2 & 2 & 0 & 2 & 2 \\ 0 & -1 & 0 & -1 & -2 & 0 \end{pmatrix}
\longmapsto \begin{pmatrix} 1 & 1 & 1 & 1 & 2 & 0 \\ 0 & 1 & 1 & 0 & 1 & 1 \\ 0 & 0 & -1 & 1 & 1 & -1 \\ 0 & 0 & 0 & 0 & 0 & 0 \\ 0 & 0 & 1 & -1 & -1 & 1 \end{pmatrix}
\]
\[
\longmapsto \begin{pmatrix} 1 & 1 & 1 & 1 & 2 & 0 \\ 0 & 1 & 1 & 0 & 1 & 1 \\ 0 & 0 & -1 & 1 & 1 & -1 \\ 0 & 0 & 0 & 0 & 0 & 0 \\ 0 & 0 & 0 & 0 & 0 & 0 \end{pmatrix}
\]
由定理~\mtref{thm:3.2.11} 我们读出, $U'$, 即由 $A'$ 的列生成的 $\mathbb{R}^5$ 的线性子空间,
有维数 3, 由 $(A', A)$ 的列生成的线性子空间亦然, 后者与 $U + U'$ 一致.因此有
$\dim_K U' = \dim_K(U + U')$.利用 $U' \subset U + U'$, 由推论~\mtref{cor:1.5.14} (ii) 得
$U' = U + U'$, 从而 $U \subset U'$.

还剩要确定 $U$ 的维数, 即 $A$ 的秩.为此把矩阵 $A$ 化为行阶梯形:
\[
\begin{pmatrix} 1 & 2 & 0 \\ 0 & 3 & 3 \\ 1 & 4 & 2 \\ 0 & 1 & 1 \\ 1 & 2 & 0 \end{pmatrix}
\longmapsto \begin{pmatrix} 1 & 2 & 0 \\ 0 & 3 & 3 \\ 0 & 2 & 2 \\ 0 & 1 & 1 \\ 0 & 0 & 0 \end{pmatrix}
\longmapsto \begin{pmatrix} 1 & 2 & 0 \\ 0 & 1 & 1 \\ 0 & 0 & 0 \\ 0 & 0 & 0 \\ 0 & 0 & 0 \end{pmatrix}
\]
由定理~\mtref{thm:3.2.11} 得 $\operatorname{rg} A = 2$ 与 $\dim_K U = 2$.

\noindent\textbf{补充:} 注意我们并没有可能直接从矩阵 $(A', A)$ 的行阶梯形读出 $A$ 的秩.不过,
我们也可以由 $(A', A)$ 的行阶梯形的最后 3 列构造一个新矩阵, 再把它化为行阶梯形.由于所得的
矩阵归根到底是从 $A$ 经初等行变换得到的, 其秩与 $U$ 的维数一致.在我们的情形下, 这个
方法比把 $A$ 单独化为行阶梯形略微简单一些.

\subsection*{习题 3.2/9}

\noindent\textbf{起点:} 设 $\mathbb{F}_2$ 是含 2 个元素的域, 即由零元素 $0$ 与单位元 $1$ 组成,
其中 $1 + 1 = 0$.要考虑 $\mathbb{F}_2$-向量空间 $\mathbb{F}_2^{2000}$ 及其一个线性子空间
$U \subset \mathbb{F}_2^{2000}$, 后者由一族向量 $(v_i)_{i = 1, \ldots, 2000}$ 生成.向量
$v_i = (v_{ij})_{j = 1, \ldots, 2000}$(利用 $1 + 1 = 0$)定义如下:
\[
v_{ij} = \begin{cases} 0, & \text{对 } i \equiv 1, 2, 3, 5, 6, 7 \ (8) \text{ 且 } j \text{ 任意}, \\ 0, & \text{对 } i \equiv 4 \ (8) \text{ 且 } j = i, \\ 0, & \text{对 } i \equiv 0 \ (8) \text{ 且 } j \in \{i - 1, i + 1\}, \\ 1, & \text{其余情形}. \end{cases}
\]
其中对自然数 $n > 0$ 与 $r \in \{0, \ldots, n - 1\}$, $i \equiv r \ (n)$ 表示 $i$ 除以 $n$
的余数为 $r$.

\noindent\textbf{目标:} 给出 $U$ 的一个基, 并确定 $\dim U$.

\noindent\textbf{策略:} 宜把向量 $v_i$ 按定理~\mtref{thm:3.2.11} 视为一个矩阵 $A$ 的列向量.
然后我们可以尝试借助初等行变换把 $A$ 化为行阶梯形, 从而在利用定理~\mtref{thm:3.2.11} 的同时
读出 $U$ 的维数以及一个基.此时只需考察向量 $v_i \neq 0$, 即指标 $i \equiv 0 \ (4)$ 的那些.
为理解所需变换的原则, 先选取 $A$ 的左上角一个易处理的片断, 并把该片断化为行阶梯形.由此
获得这样的印象: 指标 $i \equiv 0 \ (4)$ 的向量 $v_i$ 构成一个线性无关组.因此可以首先尝试用
常规方式验证这一点, 即设一个表示零的线性组合.

\noindent\textbf{解题路径:} 我们只需考虑指标 $i \equiv 0 \ (4)$ 的向量 $v_i$, 因为其余向量都是
平凡的.为证明这些向量构成一个线性无关组, 考察一个方程
\begin{equation}
\sum_{k = 1}^{500} \alpha_k v_{4 \cdot k} = 0, \tag{$*$}
\end{equation}
其系数 $\alpha_k \in \mathbb{F}_2$.把 $v_i = (v_{ij})_{j = 1, \ldots, 2000}$ 写成分量形式,
对 $k = 1, \ldots, 500$ 有 $v_{4 \cdot k, 1} = 1$.由此从 $(*)$ 得方程
\begin{equation}
\sum_{k = 1}^{500} \alpha_k = 0\text{.} \tag{$**$}
\end{equation}
为证明系数 $\alpha_k$ 消失, 选取一个指标 $k_0 \in \{1, \ldots, 500\}$, 并先假定 $k_0$ 是奇数.
对 $j_0 = 4k_0$ 则得 $j_0 \equiv 4 \ (8)$, 且 $v_{4 \cdot k, j_0} = 0$ 恰好对 $k = k_0$ 成立.
这意味着 $\sum_{k \neq k_0} \alpha_k = 0$, 并与 $(**)$ 一起给出 $\alpha_{k_0} = 0$.另一方面,
设 $k_0$ 为偶数, $j_0 = 4k_0 + 1$.则 $v_{4 \cdot k, j_0} = 0$ 恰好对 $k = k_0$ 成立.相应地
得到 $\sum_{k \neq k_0} \alpha_k = 0$, 并与 $(**)$ 一起对偶数的 $k_0$ 也给出
$\alpha_{k_0} = 0$.总之, 对所有 $k = 1, \ldots, 500$ 得 $\alpha_k = 0$.因此指标
$i \equiv 0 \ (4)$ 的向量 $v_i$ 线性无关, 并构成 $U$ 的一个基.这个基由 500 个元素组成,
所以 $\dim_{\mathbb{F}_2} U = 500$.

\noindent\textbf{补充:} 为求解本题, 自然也可以把向量 $v_i$ 视为一个矩阵 $A$ 的列向量, 即
\[
A = (v_{ij})_{j, i = 1, \ldots, 2000},
\]
其中 $j$ 是行指标而 $i$ 是列指标.此时按行来描述元素 $v_{ij}$, 即按 $A$ 的行参数 $j$ 排序,
要稍稍多花些功夫.但随后可以毫无问题地借助初等行变换把 $A$ 化为行阶梯形, 并且这里同样得到
所希望的结果: 即 $k = 1, \ldots, 500$ 的向量 $v_{4 \cdot k}$ 构成 $U$ 的一个基.

\subsection*{习题 3.3/3}

\noindent\textbf{起点:} 给定一个有限维 $K$-向量空间 $U$ 及其自同态 $f \colon U \longrightarrow U$.

\noindent\textbf{目标:} 要证明下列断言的等价性:
\begin{enumerate}
\item[(i)] $\operatorname{rg} f < \dim_K U$;
\item[(ii)] 存在满足 $f \circ g = 0$ 的 $K$-线性映射 $g \colon U \longrightarrow U$, $g \neq 0$;
\item[(iii)] 存在满足 $g \circ f = 0$ 的 $K$-线性映射 $g \colon U \longrightarrow U$, $g \neq 0$.
\end{enumerate}

\noindent\textbf{策略:} 对自同态 $f, g \colon U \longrightarrow U$, $f \circ g = 0$ 等价于
$\operatorname{im} g \subset \ker f$, 相应地 $g \circ f = 0$ 等价于 $\operatorname{im} f \subset \ker g$.
因此解题的关键在于构造对像或核带有一定要求的 $K$-线性映射 $g \colon U \longrightarrow U$.
定理~\mtref{thm:2.1.7} (ii) 是有用的, 它使我们能够通过选取 $U$ 的一个基并指定其在 $g$ 下的希望像,
来定义线性映射 $g \colon U \longrightarrow U$.此外, 借助基扩充定理~\mtref{thm:1.5.8} 以及
定理~\mtref{thm:1.5.6} 和推论~\mtref{cor:1.5.14} (i), 我们拥有选取基的灵活工具.

\noindent\textbf{解题路径:} 先设 (i) 成立, 即 $\operatorname{rg} f < \dim_K U$.为推出断言 (ii),
选取 $\ker f$ 的一个基 $x_1, \ldots, x_m$, 并用向量 $x_{m+1}, \ldots, x_n \in U$ 把它扩充为
$U$ 的一个基.由 $\operatorname{rg} f < \dim_K U$, 利用维数公式定理~\mtref{thm:2.1.10} 得
\[
m = \dim_K(\ker f) = \dim_K U - \operatorname{rg} f > 0\text{.}
\]
然后考察由线性扩张定义的 $K$-线性映射
\[
g \colon U \longrightarrow U, \qquad x_i \longmapsto \begin{cases} x_i, & \text{对 } i = 1, \ldots, m, \\ 0, & \text{对 } i = m + 1, \ldots, n\text{.} \end{cases}
\]
由于 $\operatorname{im} g \subset \ker f$, 该映射满足 $f \circ g = 0$, 并且由于 $m > 0$ 也不是
平凡的.因此由 (i) 推出 (ii).

现设 (ii) 成立, 即存在 $K$-线性映射 $g \colon U \longrightarrow U$, 使得 $g \neq 0$ 且
$f \circ g = 0$.则 $0 \neq \operatorname{im} g \subset \ker f$, 从而由维数公式
定理~\mtref{thm:2.1.10} 得 $\operatorname{rg} f < \dim_K U$, 即 (i) 成立.

最后再设 (i) 成立.为由此推出 (iii), 我们选取 $\operatorname{im} f$ 的一个基
$x_1, \ldots, x_m$, 并用向量 $x_{m+1}, \ldots, x_n$ 把它扩充为 $U$ 的一个基.然后考察由线性
扩张定义的 $K$-线性映射
\[
g \colon U \longrightarrow U, \qquad x_i \longmapsto \begin{cases} 0, & \text{对 } i = 1, \ldots, m, \\ x_i, & \text{对 } i = m + 1, \ldots, n\text{.} \end{cases}
\]
由于 $m < n$, 该映射是非平凡的, 且满足 $\operatorname{im} f \subset \ker g$, 即
$g \circ f = 0$.于是由 (i) 推出 (iii).

反之, 设 (iii) 成立, 并设 $g \colon U \longrightarrow U$ 是满足 $g \circ f = 0$ 的非平凡
$K$-线性映射.则由 $g \neq 0$ 得 $\dim_K(\ker g) < \dim_K U$, 从而由
$\operatorname{im} f \subset \ker g$ 也有 $\operatorname{rg} f < \dim_K U$, 即 (i).

\noindent\textbf{补充:} 可以思考, 对无限维 $K$-向量空间 $U$ 的自同态 $f \colon U \longrightarrow U$,
断言 (ii) 与 (iii) 的等价性将失去其有效性.例如, 若 $\ker f \neq 0$, 但 $\operatorname{im} f = U$,
则对每个满足 $g \circ f = 0$ 的 $K$-线性映射 $g \colon U \longrightarrow U$ 都有 $g = 0$.另一方面,
由于 $\ker f \neq 0$, 存在满足 $f \circ g = 0$ 的非平凡 $K$-线性映射 $g \colon U \longrightarrow U$;
只要把上面应用的构造方法推广到任意长度基的情形即可.

反之, 若 $\ker f = 0$ 且 $\operatorname{im} f \subsetneq U$, 则每个满足 $f \circ g = 0$ 的
$K$-线性映射 $g \colon U \longrightarrow U$ 都是平凡的, 尽管存在满足 $g \circ f = 0$ 的非平凡
$K$-线性映射 $g \colon U \longrightarrow U$.这里同样使用上面的构造方法, 并推广到任意长度基的情形.

\subsection*{习题 3.3/7}

\noindent\textbf{起点:} 对 $m, r, s \in \mathbb{N} - \{0\}$, 给定矩阵
$A \in K^{m \times m}$, $B \in K^{m \times s}$, $C \in K^{r \times m}$ 与
$D \in K^{r \times s}$, 其中 $A$ 可逆.由这些矩阵组成的矩阵
$T = \begin{pmatrix} A & B \\ C & D \end{pmatrix} \in K^{(m + r) \times (m + s)}$ 有秩 $m$.

\noindent\textbf{目标:} 要证明 $D = C \cdot A^{-1}B$.

\noindent\textbf{策略:} 首先并不直接看出, 在给定的秩条件下, 即 $\operatorname{rg} A = m = \operatorname{rg} T$,
$D$ 应当由 $A, B, C$ 唯一确定.为弄清事情的实质, 宜用初等行变换来考察情形.$A$ 可逆, 因此
$A$ 可以借助初等行变换化为单位矩阵 $E_m \in K^{m \times m}$.这些行变换中的每一个都对应于从
左边乘以一个相应的 $K^{m \times m}$ 中的初等矩阵, 而所有这些初等矩阵的乘积正好是 $A^{-1}$.
现在把所有这些行变换, 即从左边乘以 $A^{-1}$, 施行于矩阵 $(A, B)$, 得到作为行阶梯形的矩阵
$(E_m, A^{-1} \cdot B)$.同样的初等行变换我们也可以施行于矩阵 $T$, 这对应于从左边乘以一个由
$A^{-1}$ 与单位矩阵 $E_r \in K^{r \times r}$ 组成的矩阵:
\[
\begin{pmatrix} A^{-1} & 0 \\ 0 & E_r \end{pmatrix} \cdot \begin{pmatrix} A & B \\ C & D \end{pmatrix} = T', \qquad \text{其中 } T' = \begin{pmatrix} E_m & A^{-1}B \\ C & D \end{pmatrix}\text{.}
\]
我们现在可以对 $T'$ 应用进一步的初等行变换, 目标是把 $C$ 化为零矩阵.为此按标准方式处理:
依次对所有 $1 \leq i \leq r$ 与 $1 \leq j \leq m$, 把 $C$ 在位置 $(i, j)$ 的系数与
$(E_m, A^{-1} \cdot B)$ 的第 $j$ 行相乘, 并从 $(C, D)$ 的第 $i$ 行减去这个倍数.一旦在 $C$ 的
位置上达到零矩阵, 就已经得到 $T$ 即 $T'$ 的一个行阶梯形, 于是 $D$ 也必然已被化为零矩阵,
因为否则会有 $\operatorname{rg} T = \operatorname{rg} T' > m$.如果我们精确追踪这次化到行阶梯形的
变换的各个步骤, 就能(以略微繁琐的方式)看出, 这里 $D$ 被 $D - C \cdot A^{-1} \cdot B$ 所取代,
从而最终 $D = C \cdot A^{-1} \cdot B$ 成立.

\noindent\textbf{解题路径:} 利用单位矩阵 $E_r \in K^{r \times r}$, 如同上面已经认识到的, 有
\[
\begin{pmatrix} A^{-1} & 0 \\ 0 & E_r \end{pmatrix} \cdot \begin{pmatrix} A & B \\ C & D \end{pmatrix} = T' = \begin{pmatrix} E_m & A^{-1}B \\ C & D \end{pmatrix}\text{.}
\]
这里可以观察到, $(2 \times 2)$-矩阵相乘的规则可以搬到行数与列数合适的块矩阵的相乘上, 只要把
系数的普通乘积换成矩阵乘积即可.现在我们用初等行变换把 $T'$ 化为行阶梯形, 具体地, 从
$(C, D)$ 的行减去 $(E_m, A^{-1} \cdot B)$ 的行的合适倍数, 目标是让 $C$ 被零矩阵取代.如上面
所说明的, 此时矩阵 $D$ 基于秩的理由同样被化为零矩阵.

现在, 所需类型的一次初等行变换, 即把某第 $j$ 行($1 \leq j \leq m$)的 $\alpha$ 倍加到某第 $i$
行($m + 1 \leq i \leq m + r$)上, 对应于从左边乘以初等矩阵 $E + \alpha E_{ij}$, 即乘以一个
带合适的 $P \in K^{r \times m}$ 的类型 $\begin{pmatrix} E_m & 0 \\ P & E_r \end{pmatrix}$ 的矩阵.
由于这个类型的矩阵的乘积仍属于这个类型, 存在矩阵 $P \in K^{r \times m}$, 使得
\[
\begin{pmatrix} E_m & 0 \\ P & E_r \end{pmatrix} \cdot \begin{pmatrix} E_m & A^{-1}B \\ C & D \end{pmatrix} = \begin{pmatrix} E_m & A^{-1}B \\ P + C & PA^{-1}B + D \end{pmatrix} = \begin{pmatrix} E_m & A^{-1}B \\ 0 & 0 \end{pmatrix}\text{.}
\]
由此读出 $P = -C$, 然后如所期望地读出 $D = C \cdot A^{-1} \cdot B$.

\noindent\textbf{补充:} 我们还想给出一个替代解法, 它把给定的矩阵解释为线性映射, 从而避免初等
行变换以及矩阵层面的计算.矩阵 $A, B, C$ 与 $D$ 给出一个由 $K$-线性映射组成的图表, 其中 $A$
是一个同构.要证明的是: 由 $D$ 与 $C \circ A^{-1} \circ B$ 给出的映射 $K^s \longrightarrow K^r$
相一致.为验证这一点, 我们把 $K^{m + s}$ 与直和 $K^m \oplus K^s$ 等同起来, 相应地也把
$K^{m + r}$ 与 $K^m \oplus K^r$ 等同起来.然后考察由 $T$ 给出的线性映射
\[
T \colon K^m \oplus K^s \longrightarrow K^m \oplus K^r, \quad x \oplus y \longmapsto (Ax \oplus Cx) + (By \oplus Dy)\text{.}
\]
设 $H \colon K^m \oplus K^r \longrightarrow K^m$ 是到第一个被加项的投影.按假定有
$\operatorname{rg} T = m$, 但同样也有 $\operatorname{rg}(H \circ T) = m$, 因为复合映射
\[
K^m \xrightarrow{T|_{K^m}} K^m \oplus K^r \xrightarrow{H} K^m
\]
与 $A$ 一致, 因而是同构.特别地, 由 $H$ 诱导的线性映射 $\operatorname{im} T \longrightarrow K^m$
是一个同构, 因为其秩为 $m$.

现在我们选取 $y \in K^s$, 并考虑元素 $x = A^{-1}By \in K^m$.$x$ 在 $T$ 下的像为
$Ax \oplus Cx = By \oplus CA^{-1}By$, 而 $y$ 被映到 $By \oplus Dy$.两个像在 $H$ 下都投影到元素
$By \in K^m$.由于 $H$ 诱导同构 $\operatorname{im} T \longrightarrow K^m$, 我们得到对所有
$y \in K^s$ 有 $Dy = CA^{-1}By$, 从而 $D = CA^{-1}B$.

\subsection*{习题 3.4/1}

\noindent\textbf{起点:} 要考虑两个系统 $X$ 与 $Y$, 它们分别由具体给出的向量
$x_1, \ldots, x_5$ 与 $y_1, \ldots, y_5 \in \mathbb{R}^5$ 组成.把 $x_j$ 与 $y_j$ 视为列向量,
即为矩阵
\[
A = (x_1, \ldots, x_5) = \begin{pmatrix} 1 & 1 & 1 & 1 & 0 \\ 1 & 0 & 1 & 1 & 1 \\ 1 & 2 & 1 & 2 & 1 \\ 2 & 1 & 2 & 1 & 2 \\ 0 & 0 & 1 & 0 & 0 \end{pmatrix},
\qquad
B = (y_1, \ldots, y_5) = \begin{pmatrix} 1 & 1 & 1 & 2 & 0 \\ 2 & 1 & 3 & 3 & 0 \\ 3 & 2 & 4 & 5 & 1 \\ 4 & 3 & 6 & 6 & 1 \\ 4 & 4 & 7 & 8 & 1 \end{pmatrix}\text{.}
\]

\noindent\textbf{目标:} 要证明 $X$ 与 $Y$ 构成 $\mathbb{R}^5$ 的基.要计算基变换矩阵
$A_{\mathrm{id}, X, Y}$ 与 $A_{\mathrm{id}, Y, X}$.

\noindent\textbf{策略:} 为验证 $X$ 与 $Y$ 构成 $\mathbb{R}^5$ 的基, 只需证明 $X$ 与 $Y$ 线性无关.
我们可以为此使用大定理~\mtref{bigthm:3.2.4} 或定理~\mtref{thm:3.2.11} 形式的高斯消元法.然后, 为确定
基变换矩阵, 例如 $A_{\mathrm{id}, X, Y}$, 可以设线性组合
\[
x_j = \sum_{i = 1}^{5} \alpha_{ij} y_i, \qquad j = 1, \ldots, 5,
\]
并尝试计算系数 $\alpha_{ij} \in \mathbb{R}$.原则上, 这是 25 个未知量 $\alpha_{ij}$ 的一个含
25 个线性方程的方程组, 乍一看并不太容易处理.不过这个方程组可以相当简单地求解, 我们稍后在
补充中将作说明.

首先我们还是应该尝试应用所学的基变换理论.把 $x_j$ 解释为列向量, 就会发现
$A = (x_1, \ldots, x_5)$ 恰好给出基变换矩阵 $A_{\mathrm{id}, X, E}$, 其中 $E$ 是 $\mathbb{R}^5$
的标准基.相应地, 若把 $y_i$ 看作矩阵 $B$ 的列, 则有 $B = (y_1, \ldots, y_5) = A_{\mathrm{id}, Y, E}$.
现在我们可以把从 $X$ 到 $Y$ 的基变换分解为从 $X$ 到 $E$,再从 $E$ 到 $Y$ 的基变换.此时由
定理~\mtref{thm:3.1.9} 与注记~\mtref{rem:3.4.2} 有
\[
\begin{array}{l}
A_{\mathrm{id}, X, Y} = A_{\mathrm{id}, E, Y} \cdot A_{\mathrm{id}, X, E} = A_{\mathrm{id}, Y, E}^{-1} \cdot A_{\mathrm{id}, X, E} = B^{-1} \cdot A, \\
A_{\mathrm{id}, Y, X} = A_{\mathrm{id}, E, X} \cdot A_{\mathrm{id}, Y, E} = A_{\mathrm{id}, X, E}^{-1} \cdot A_{\mathrm{id}, Y, E} = A^{-1} \cdot B.
\end{array}
\]
因此, 除了施行矩阵乘积外, 我们只需把矩阵 $A$ 与 $B$ 求逆, 而这可以如 \ref{sec:3.3} 节所述借助初等行变换
来完成.

\noindent\textbf{解题路径:} 为证明系统 $X$ 与 $Y$ 构成 $\mathbb{R}^5$ 的基, 只需证明这些系统线性
无关.把给定的向量视为列向量, 考察矩阵 $A = (x_1, \ldots, x_5)$ 与 $B = (y_1, \ldots, y_5)$,
则要证明 $\operatorname{rg} A = 5 = \operatorname{rg} B$.我们现在为两个矩阵确定行阶梯形.若由此
得到秩 5, 就可以通过进一步的初等行变换把单位矩阵 $E_5 \in \mathbb{R}^{5 \times 5}$ 作为行阶梯形
达到.同时我们也把这些所需变换施行于 $E_5$, 使得由此在末尾得到逆矩阵 $A^{-1}$ 即 $B^{-1}$.
我们从矩阵 $A$ 开始:
\[
A \colon \begin{pmatrix} 1 & 1 & 1 & 1 & 0 \\ 1 & 0 & 1 & 1 & 1 \\ 1 & 2 & 1 & 2 & 1 \\ 2 & 1 & 2 & 1 & 2 \\ 0 & 0 & 1 & 0 & 0 \end{pmatrix}
\longmapsto^{(1)} \begin{pmatrix} 1 & 1 & 1 & 1 & 0 \\ 0 & -1 & 0 & 0 & 1 \\ 0 & 1 & 0 & 1 & 1 \\ 0 & -1 & 0 & -1 & 2 \\ 0 & 0 & 1 & 0 & 0 \end{pmatrix}
\longmapsto^{(2)} \begin{pmatrix} 1 & 1 & 1 & 1 & 0 \\ 0 & -1 & 0 & 0 & 1 \\ 0 & 0 & 0 & 1 & 2 \\ 0 & 0 & 0 & -1 & 1 \\ 0 & 0 & 1 & 0 & 0 \end{pmatrix}
\]
\[
E_5 \colon \begin{pmatrix} 1 & 0 & 0 & 0 & 0 \\ 0 & 1 & 0 & 0 & 0 \\ 0 & 0 & 1 & 0 & 0 \\ 0 & 0 & 0 & 1 & 0 \\ 0 & 0 & 0 & 0 & 1 \end{pmatrix}
\longmapsto^{(1)} \begin{pmatrix} 1 & 0 & 0 & 0 & 0 \\ -1 & 1 & 0 & 0 & 0 \\ -1 & 0 & 1 & 0 & 0 \\ -2 & 0 & 0 & 1 & 0 \\ 0 & 0 & 0 & 0 & 1 \end{pmatrix}
\longmapsto^{(2)} \begin{pmatrix} 1 & 0 & 0 & 0 & 0 \\ -1 & 1 & 0 & 0 & 0 \\ -2 & 1 & 1 & 0 & 0 \\ -1 & -1 & 0 & 1 & 0 \\ 0 & 0 & 0 & 0 & 1 \end{pmatrix}
\]
\begin{enumerate}
\item[(1)] 从第 2,3 行减去第 1 行, 从第 4 行减去第 1 行的 2 倍;
\item[(2)] 把第 2 行加到第 3 行, 并从第 4 行减去.
\end{enumerate}
\[
\longmapsto^{(3)} \begin{pmatrix} 1 & 1 & 1 & 1 & 0 \\ 0 & 1 & 0 & 0 & -1 \\ 0 & 0 & 1 & 0 & 0 \\ 0 & 0 & 0 & 1 & 2 \\ 0 & 0 & 0 & -1 & 1 \end{pmatrix}
\longmapsto^{(4)} \begin{pmatrix} 1 & 1 & 1 & 1 & 0 \\ 0 & 1 & 0 & 0 & -1 \\ 0 & 0 & 1 & 0 & 0 \\ 0 & 0 & 0 & 1 & 2 \\ 0 & 0 & 0 & 0 & 1 \end{pmatrix}
\longmapsto^{(5)} \begin{pmatrix} 1 & 1 & 1 & 1 & 0 \\ 0 & 1 & 0 & 0 & 0 \\ 0 & 0 & 1 & 0 & 0 \\ 0 & 0 & 0 & 1 & 0 \\ 0 & 0 & 0 & 0 & 1 \end{pmatrix}
\]
\[
\longmapsto^{(3)} \begin{pmatrix} 1 & 0 & 0 & 0 & 0 \\ 1 & -1 & 0 & 0 & 0 \\ 0 & 0 & 0 & 0 & 1 \\ -2 & 1 & 1 & 0 & 0 \\ -1 & -1 & 0 & 1 & 0 \end{pmatrix}
\longmapsto^{(4)} \begin{pmatrix} 1 & 0 & 0 & 0 & 0 \\ 1 & -1 & 0 & 0 & 0 \\ 0 & 0 & 0 & 0 & 1 \\ -2 & 1 & 1 & 0 & 0 \\ -1 & 0 & \frac{1}{3} & \frac{1}{3} & 0 \end{pmatrix}
\longmapsto^{(5)} \begin{pmatrix} 1 & 0 & 0 & 0 & 0 \\ 0 & -1 & \frac{1}{3} & \frac{1}{3} & 0 \\ 0 & 0 & 0 & 0 & 1 \\ 0 & 1 & \frac{1}{3} & -\frac{2}{3} & 0 \\ -1 & 0 & \frac{1}{3} & \frac{1}{3} & 0 \end{pmatrix}
\]
\begin{enumerate}
\item[(3)] 第 2 行乘以 $-1$, 把第 5 行与第 3,4 行交换;
\item[(4)] 把第 4 行加到第 5 行, 第 5 行乘以 $\frac{1}{3}$; 达到 $A$ 的行阶梯形, 故
$\operatorname{rg} A = 5$, 即 $A$ 可逆;
\item[(5)] 把第 5 行加到第 2 行, 从第 4 行减去第 5 行的 2 倍.
\end{enumerate}
\[
\longmapsto^{(6)} \begin{pmatrix} 1 & 1 & 1 & 0 & 0 \\ 0 & 1 & 0 & 0 & 0 \\ 0 & 0 & 1 & 0 & 0 \\ 0 & 0 & 0 & 1 & 0 \\ 0 & 0 & 0 & 0 & 1 \end{pmatrix}
\longmapsto^{(7)} \begin{pmatrix} 1 & 1 & 0 & 0 & 0 \\ 0 & 1 & 0 & 0 & 0 \\ 0 & 0 & 1 & 0 & 0 \\ 0 & 0 & 0 & 1 & 0 \\ 0 & 0 & 0 & 0 & 1 \end{pmatrix}
\longmapsto^{(8)} \begin{pmatrix} 1 & 0 & 0 & 0 & 0 \\ 0 & 1 & 0 & 0 & 0 \\ 0 & 0 & 1 & 0 & 0 \\ 0 & 0 & 0 & 1 & 0 \\ 0 & 0 & 0 & 0 & 1 \end{pmatrix}
\]
\[
{\setlength{\arraycolsep}{3.5pt}
\longmapsto^{(6)} \begin{pmatrix} 1 & -1 & -\frac{1}{3} & \frac{2}{3} & 0 \\ 0 & -1 & \frac{1}{3} & \frac{1}{3} & 0 \\ 0 & 0 & 0 & 0 & 1 \\ 0 & 1 & \frac{1}{3} & -\frac{2}{3} & 0 \\ -1 & 0 & \frac{1}{3} & \frac{1}{3} & 0 \end{pmatrix}
\longmapsto^{(7)} \begin{pmatrix} 1 & -1 & -\frac{1}{3} & \frac{2}{3} & -1 \\ 0 & -1 & \frac{1}{3} & \frac{1}{3} & 0 \\ 0 & 0 & 0 & 0 & 1 \\ 0 & 1 & \frac{1}{3} & -\frac{2}{3} & 0 \\ -1 & 0 & \frac{1}{3} & \frac{1}{3} & 0 \end{pmatrix}
\longmapsto^{(8)} \begin{pmatrix} 1 & 0 & -\frac{2}{3} & \frac{1}{3} & -1 \\ 0 & -1 & \frac{1}{3} & \frac{1}{3} & 0 \\ 0 & 0 & 0 & 0 & 1 \\ 0 & 1 & \frac{1}{3} & -\frac{2}{3} & 0 \\ -1 & 0 & \frac{1}{3} & \frac{1}{3} & 0 \end{pmatrix}
}
\]
\begin{enumerate}
\item[(6)] 从第 1 行减去第 4 行;
\item[(7)] 从第 1 行减去第 3 行;
\item[(8)] 从第 1 行减去第 2 行; 总之, $A$ 借助初等行变换被化为单位矩阵 $E_5$, 同时 $E_5$
被化为逆矩阵 $A^{-1}$.
\end{enumerate}

接下来处理矩阵 $B$:
\[
B \colon \begin{pmatrix} 1 & 1 & 1 & 2 & 0 \\ 2 & 1 & 3 & 3 & 0 \\ 3 & 2 & 4 & 5 & 1 \\ 4 & 3 & 6 & 6 & 1 \\ 4 & 4 & 7 & 8 & 1 \end{pmatrix}
\longmapsto^{(1)} \begin{pmatrix} 1 & 1 & 1 & 2 & 0 \\ 0 & -1 & 1 & -1 & 0 \\ 0 & -1 & 1 & -1 & 1 \\ 0 & -1 & 2 & -2 & 1 \\ 0 & 0 & 3 & 0 & 1 \end{pmatrix}
\longmapsto^{(2)} \begin{pmatrix} 1 & 1 & 1 & 2 & 0 \\ 0 & -1 & 1 & -1 & 0 \\ 0 & 0 & 0 & 0 & 1 \\ 0 & 0 & 1 & -1 & 1 \\ 0 & 0 & 3 & 0 & 1 \end{pmatrix}
\]
\[
E_5 \colon \begin{pmatrix} 1 & 0 & 0 & 0 & 0 \\ 0 & 1 & 0 & 0 & 0 \\ 0 & 0 & 1 & 0 & 0 \\ 0 & 0 & 0 & 1 & 0 \\ 0 & 0 & 0 & 0 & 1 \end{pmatrix}
\longmapsto^{(1)} \begin{pmatrix} 1 & 0 & 0 & 0 & 0 \\ -2 & 1 & 0 & 0 & 0 \\ -3 & 0 & 1 & 0 & 0 \\ -4 & 0 & 0 & 1 & 0 \\ -4 & 0 & 0 & 0 & 1 \end{pmatrix}
\longmapsto^{(2)} \begin{pmatrix} 1 & 0 & 0 & 0 & 0 \\ -2 & 1 & 0 & 0 & 0 \\ -1 & -1 & 1 & 0 & 0 \\ -2 & -1 & 0 & 1 & 0 \\ -4 & 0 & 0 & 0 & 1 \end{pmatrix}
\]
\begin{enumerate}
\item[(1)] 从第 2 行减去第 1 行的 2 倍, 从第 3 行减去 3 倍, 从第 4 行减去 4 倍, 从第 5 行减去
4 倍;
\item[(2)] 从第 3,4 行减去第 2 行.
\end{enumerate}
\[
\longmapsto^{(3)} \begin{pmatrix} 1 & 1 & 1 & 2 & 0 \\ 0 & 1 & -1 & 1 & 0 \\ 0 & 0 & 1 & -1 & 1 \\ 0 & 0 & 3 & 0 & 1 \\ 0 & 0 & 0 & 0 & 1 \end{pmatrix}
\longmapsto^{(4)} \begin{pmatrix} 1 & 1 & 1 & 2 & 0 \\ 0 & 1 & -1 & 1 & 0 \\ 0 & 0 & 1 & -1 & 1 \\ 0 & 0 & 0 & 3 & -2 \\ 0 & 0 & 0 & 0 & 1 \end{pmatrix}
\longmapsto^{(5)} \begin{pmatrix} 1 & 1 & 1 & 2 & 0 \\ 0 & 1 & -1 & 1 & 0 \\ 0 & 0 & 1 & -1 & 0 \\ 0 & 0 & 0 & 1 & 0 \\ 0 & 0 & 0 & 0 & 1 \end{pmatrix}
\]
\[
\longmapsto^{(3)} \begin{pmatrix} 1 & 0 & 0 & 0 & 0 \\ 2 & -1 & 0 & 0 & 0 \\ -2 & -1 & 0 & 1 & 0 \\ -4 & 0 & 0 & 0 & 1 \\ -1 & -1 & 1 & 0 & 0 \end{pmatrix}
\longmapsto^{(4)} \begin{pmatrix} 1 & 0 & 0 & 0 & 0 \\ 2 & -1 & 0 & 0 & 0 \\ -2 & -1 & 0 & 1 & 0 \\ 2 & 3 & 0 & -3 & 1 \\ -1 & -1 & 1 & 0 & 0 \end{pmatrix}
\longmapsto^{(5)} \begin{pmatrix} 1 & 0 & 0 & 0 & 0 \\ 2 & -1 & 0 & 0 & 0 \\ -1 & 0 & -1 & 1 & 0 \\ 0 & \frac{1}{3} & \frac{2}{3} & -1 & \frac{1}{3} \\ -1 & -1 & 1 & 0 & 0 \end{pmatrix}
\]
\begin{enumerate}
\item[(3)] 第 2 行乘以 $-1$, 把第 3 行与第 4,5 行交换;
\item[(4)] 从第 4 行减去第 3 行的 3 倍; 达到 $B$ 的行阶梯形, 故 $\operatorname{rg} B = 5$, 即
$B$ 可逆;
\item[(5)] 从第 3 行减去第 5 行, 把第 5 行的 2 倍加到第 4 行, 第 4 行乘以 $\frac{1}{3}$.
\end{enumerate}
\[
\longmapsto^{(6)} \begin{pmatrix} 1 & 1 & 1 & 0 & 0 \\ 0 & 1 & -1 & 0 & 0 \\ 0 & 0 & 1 & 0 & 0 \\ 0 & 0 & 0 & 1 & 0 \\ 0 & 0 & 0 & 0 & 1 \end{pmatrix}
\longmapsto^{(7)} \begin{pmatrix} 1 & 1 & 0 & 0 & 0 \\ 0 & 1 & 0 & 0 & 0 \\ 0 & 0 & 1 & 0 & 0 \\ 0 & 0 & 0 & 1 & 0 \\ 0 & 0 & 0 & 0 & 1 \end{pmatrix}
\longmapsto^{(8)} \begin{pmatrix} 1 & 0 & 0 & 0 & 0 \\ 0 & 1 & 0 & 0 & 0 \\ 0 & 0 & 1 & 0 & 0 \\ 0 & 0 & 0 & 1 & 0 \\ 0 & 0 & 0 & 0 & 1 \end{pmatrix}
\]
\[
{\setlength{\arraycolsep}{3.2pt}
\longmapsto^{(6)} \begin{pmatrix} 1 & -\frac{2}{3} & -\frac{4}{3} & 2 & -\frac{2}{3} \\ 2 & -\frac{4}{3} & -\frac{2}{3} & 1 & -\frac{1}{3} \\ -1 & \frac{1}{3} & -\frac{1}{3} & 0 & \frac{1}{3} \\ 0 & \frac{1}{3} & \frac{2}{3} & -1 & \frac{1}{3} \\ -1 & -1 & 1 & 0 & 0 \end{pmatrix}
\longmapsto^{(7)} \begin{pmatrix} 2 & -1 & -1 & 2 & -1 \\ 1 & -1 & -1 & 1 & 0 \\ -1 & \frac{1}{3} & -\frac{1}{3} & 0 & \frac{1}{3} \\ 0 & \frac{1}{3} & \frac{2}{3} & -1 & \frac{1}{3} \\ -1 & -1 & 1 & 0 & 0 \end{pmatrix}
\longmapsto^{(8)} \begin{pmatrix} 1 & 0 & 0 & 1 & -1 \\ 1 & -1 & -1 & 1 & 0 \\ -1 & \frac{1}{3} & -\frac{1}{3} & 0 & \frac{1}{3} \\ 0 & \frac{1}{3} & \frac{2}{3} & -1 & \frac{1}{3} \\ -1 & -1 & 1 & 0 & 0 \end{pmatrix}
}
\]
\begin{enumerate}
\item[(6)] 从第 1 行减去第 4 行的 2 倍, 从第 2 行减去第 4 行并把它加到第 3 行;
\item[(7)] 从第 1 行减去第 3 行, 把它加到第 2 行;
\item[(8)] 从第 1 行减去第 2 行; 总之, $B$ 借助初等行变换被化为单位矩阵 $E_5$, 同时 $E_5$
被化为逆矩阵 $B^{-1}$.
\end{enumerate}

作为计算的结果我们因此得到:
\[
A^{-1} = \begin{pmatrix} 1 & 0 & -\frac{2}{3} & \frac{1}{3} & -1 \\ 0 & -1 & \frac{1}{3} & \frac{1}{3} & 0 \\ 0 & 0 & 0 & 0 & 1 \\ 0 & 1 & \frac{1}{3} & -\frac{2}{3} & 0 \\ -1 & 0 & \frac{1}{3} & \frac{1}{3} & 0 \end{pmatrix},
\qquad
B^{-1} = \begin{pmatrix} 1 & 0 & 0 & 1 & -1 \\ 1 & -1 & -1 & 1 & 0 \\ -1 & \frac{1}{3} & -\frac{1}{3} & 0 & \frac{1}{3} \\ 0 & \frac{1}{3} & \frac{2}{3} & -1 & \frac{1}{3} \\ -1 & -1 & 1 & 0 & 0 \end{pmatrix}\text{.}
\]
我们在策略性思考中已经说明, 矩阵 $A$ 与 $B$ 正好是 $\mathbb{R}^5$ 的标准基 $E$ 的基变换矩阵
$A_{\mathrm{id}, X, E}$ 与 $A_{\mathrm{id}, Y, E}$, 并且由此所求的 $X$ 与 $Y$ 之间的基变换矩阵
可以如下计算:
\[
A_{\mathrm{id}, X, Y} = B^{-1}A = \begin{pmatrix} 1 & 0 & 0 & 1 & -1 \\ 1 & -1 & -1 & 1 & 0 \\ -1 & \frac{1}{3} & -\frac{1}{3} & 0 & \frac{1}{3} \\ 0 & \frac{1}{3} & \frac{2}{3} & -1 & \frac{1}{3} \\ -1 & -1 & 1 & 0 & 0 \end{pmatrix} \cdot \begin{pmatrix} 1 & 1 & 1 & 1 & 0 \\ 1 & 0 & 1 & 1 & 1 \\ 1 & 2 & 1 & 2 & 1 \\ 2 & 1 & 2 & 1 & 2 \\ 0 & 0 & 1 & 0 & 0 \end{pmatrix} = \begin{pmatrix} 3 & 2 & 2 & 2 & 2 \\ 1 & 0 & 1 & -1 & 0 \\ -1 & -\frac{5}{3} & -\frac{2}{3} & -\frac{4}{3} & 0 \\ -1 & \frac{1}{3} & -\frac{2}{3} & \frac{2}{3} & -1 \\ -1 & 1 & -1 & 0 & 0 \end{pmatrix}
\]
\[
A_{\mathrm{id}, Y, X} = A^{-1}B = \begin{pmatrix} 1 & 0 & -\frac{2}{3} & \frac{1}{3} & -1 \\ 0 & -1 & \frac{1}{3} & \frac{1}{3} & 0 \\ 0 & 0 & 0 & 0 & 1 \\ 0 & 1 & \frac{1}{3} & -\frac{2}{3} & 0 \\ -1 & 0 & \frac{1}{3} & \frac{1}{3} & 0 \end{pmatrix} \cdot \begin{pmatrix} 1 & 1 & 1 & 2 & 0 \\ 2 & 1 & 3 & 3 & 0 \\ 3 & 2 & 4 & 5 & 1 \\ 4 & 3 & 6 & 6 & 1 \\ 4 & 4 & 7 & 8 & 1 \end{pmatrix} = \begin{pmatrix} -\frac{11}{3} & -\frac{10}{3} & -\frac{20}{3} & -\frac{22}{3} & -\frac{4}{3} \\ \frac{1}{3} & \frac{2}{3} & \frac{1}{3} & \frac{2}{3} & \frac{2}{3} \\ 4 & 4 & 7 & 8 & 1 \\ \frac{1}{3} & -\frac{1}{3} & \frac{1}{3} & \frac{2}{3} & -\frac{1}{3} \\ \frac{4}{3} & \frac{2}{3} & \frac{7}{3} & \frac{5}{3} & \frac{2}{3} \end{pmatrix}
\]

\noindent\textbf{补充:} 基变换矩阵 $A_{\mathrm{id}, X, Y} = (\alpha_{ij})_{i, j = 1, \ldots, 5}$
由方程
\[
x_j = \sum_{i = 1}^{5} \alpha_{ij} y_i, \qquad j = 1, \ldots, 5
\]
刻画.把向量 $x_j$ 与 $y_i$ 解释为 $\mathbb{R}^5$ 中的列向量, 这些关系就写成形式
\[
x_j = (y_1, \ldots, y_5) \cdot (\alpha_{1j}, \ldots, \alpha_{5j})^t, \qquad j = 1, \ldots, 5,
\]
或者在利用矩阵 $A = (x_1, \ldots, x_5)$ 与 $B = (y_1, \ldots, y_5)$ 的情况下写成形式
$A = B \cdot A_{\mathrm{id}, X, Y}$.因此 $A_{\mathrm{id}, X, Y} = B^{-1}A$, 相应地
$A_{\mathrm{id}, Y, X} = A^{-1} \cdot B$, 正如开头已经确认的.

为在计算完成后检验逆矩阵 $A^{-1}$ 与 $B^{-1}$ 是否无差错地求出, 应当计算乘积 $A \cdot A^{-1}$
与 $B \cdot B^{-1}$, 并与单位矩阵比较.同样可以利用关系 $B \cdot A_{\mathrm{id}, X, Y} = A$ 与
$A \cdot A_{\mathrm{id}, Y, X} = B$, 以确保基变换矩阵 $A_{\mathrm{id}, X, Y}$ 与
$A_{\mathrm{id}, Y, X}$ 是正确的.

\subsection*{习题 3.4/4}

\noindent\textbf{起点:} 在全体 $(m \times n)$-矩阵的集合 $K^{m \times n}$ 上给定一个关系"$\sim$",
即对两个这样的矩阵 $A \sim B$ 成立, 如果存在可逆矩阵 $S \in \mathrm{Gl}(m, K)$ 与
$T \in \mathrm{Gl}(n, K)$ 使得 $B = SAT$.

\noindent\textbf{目标:} 要证明"$\sim$"具有等价关系的性质.要确定相应的等价类, 尤其是其个数.

\noindent\textbf{策略:} 设 $A \in K^{m \times n}$.在定理~\mtref{thm:3.4.8} 中我们考察了属于 $A$ 的
$K$-线性映射 $K^n \longrightarrow K^m$, $a \longmapsto A \cdot a$, 并借助基变换证明了存在可逆矩阵
$S \in \mathrm{Gl}(m, K)$ 与 $T \in \mathrm{Gl}(n, K)$, 使得
$SAT = \begin{pmatrix} E_r & 0 \\ 0 & 0 \end{pmatrix} =: E_r^{m \times n}$.其中 $E_r$ 是
$(r \times r)$-单位矩阵, 而 $r = \operatorname{rg} A$.在给定关系"$\sim$"的术语下, 这意味着
$A \sim E_r^{m \times n}$.因此, 如果我们证明"$\sim$"是一个等价关系, 那么矩阵
$E_r^{m \times n}$($r = 0, \ldots, \min\{m, n\}$)就是相应等价类的代表, 而且是两两不同的等价类的
代表, 因为由定理~\mtref{thm:3.4.6}, 从 $A \sim B$ 尤其推出 $\operatorname{rg} A = \operatorname{rg} B$.

\noindent\textbf{解题路径:} 我们首先证明"$\sim$"是一个等价关系.设 $A \in K^{m \times n}$.则
$A = SAT$ 对单位矩阵 $S \in \mathrm{Gl}(m, K)$ 与 $T \in \mathrm{Gl}(n, K)$ 成立.这意味着
$A \sim A$, 因而关系是自反的.再设 $A, B \in K^{m \times n}$.若 $A \sim B$, 即对某些可逆矩阵
$S \in \mathrm{Gl}(m, K)$ 与 $T \in \mathrm{Gl}(n, K)$ 有 $B = SAT$, 则 $A = S^{-1}BT^{-1}$,
其中 $S^{-1} \in \mathrm{Gl}(m, K)$ 且 $T^{-1} \in \mathrm{Gl}(n, K)$, 即 $B \sim A$.因此关系也是
对称的.最后设 $A, B, C \in K^{m \times n}$ 满足 $A \sim B$ 与 $B \sim C$.则存在可逆矩阵
$S, \tilde{S} \in \mathrm{Gl}(m, K)$ 与 $T, \tilde{T} \in \mathrm{Gl}(n, K)$, 使得 $B = SAT$ 与
$C = \tilde{S}B\tilde{T}$ 成立.由此得 $C = (\tilde{S}S)A(T\tilde{T})$, 其中
$\tilde{S}S \in \mathrm{Gl}(m, K)$ 且 $T\tilde{T} \in \mathrm{Gl}(n, K)$, 即 $A \sim C$.因此关系
"$\sim$"也是传递的, 从而是一个等价关系.

如已经提到的, 我们由定理~\mtref{thm:3.4.8} 推断: 对秩为 $r$ 的矩阵 $A \in K^{m \times n}$ 有
$A \sim E_r^{m \times n}$.由于"$\sim$"按照定理~\mtref{thm:3.4.6} 保持矩阵的秩, 且
$\operatorname{rg} E_r^{m \times n} = r$, 矩阵 $E_r^{m \times n}$($r = 0, \ldots, \min\{m, n\}$)
构成 $K^{m \times n}$ 中关于"$\sim$"的等价类的一个代表系.因此其个数等于 $1 + \min\{m, n\}$,
即在给定行数 $m$ 与列数 $n$ 时类型 $E_r^{m \times n}$ 的矩阵的可能个数.

\subsection*{习题 3.5/4}

\noindent\textbf{起点:} 要考虑有限维 $\mathbb{R}$-向量空间之间的一个 $\mathbb{R}$-线性映射
$f \colon V \longrightarrow W$, 它关于合适的基 $X$(对 $V$)与 $Y$(对 $W$)由具体给出的矩阵
$A_{f, X, Y} \in \mathbb{R}^{4 \times 5}$ 描述, 即
\[
A_{f, X, Y} = \begin{pmatrix} 1 & 1 & 2 & 4 & 8 \\ 1 & 2 & 1 & 2 & 1 \\ 2 & 2 & 2 & 2 & 1 \\ 1 & 2 & 2 & 2 & 2 \end{pmatrix}\text{.}
\]

\noindent\textbf{目标:} 要确定 $\ker f$ 的一个基.

\noindent\textbf{策略:} 为在 $V$ 即 $W$ 中进行具体的计算, 我们必须把 $V$ 与 $W$ 中的向量表示成
所考察基的线性组合.换言之, 我们应该使用关于 $X$ 与 $Y$ 的相应的坐标列向量.于是宜利用
推论~\mtref{cor:3.1.8} 中的交换图
\begin{center}
\begin{tikzcd}
V \arrow[r, "f"] \arrow[d, "\kappa_X"'] & W \arrow[d, "\kappa_Y"] \\
\mathbb{R}^5 \arrow[r, "\tilde{f}"'] & \mathbb{R}^4,
\end{tikzcd}
\end{center}
其中 $\kappa_X$ 与 $\kappa_Y$ 是标准的同构, 它们把 $V$ 即 $W$ 中的向量送到相应的 $\mathbb{R}^5$
即 $\mathbb{R}^4$ 中的坐标列向量, 而 $\tilde{f}$ 是与矩阵 $A_{f, X, Y}$ 相乘, 即由
$x \longmapsto A_{f, X, Y} \cdot x$ 定义的映射.$\tilde{f}$ 的核与齐次线性方程组
$A_{f, X, Y} \cdot x = 0$ 的解空间 $M_{A_{f, X, Y}, 0}$ 一致, 后者可以用高斯消元法确定.由于
$\kappa_X$ 限制为同构 $\ker f \longrightarrow M_{A_{f, X, Y}, 0}$, 可以由 $M_{A_{f, X, Y}, 0}$
的一个基重构 $\ker f$ 的一个基.

\noindent\textbf{解题路径:} 我们使用上述交换图, 并把同构 $\kappa_X$ 与 $\kappa_Y$ 解释为等同.
然后要解齐次线性方程组 $A_{f, X, Y} \cdot x = 0$.为此使用高斯消元法, 把 $A_{f, X, Y}$ 化为
\ref{sec:3.5} 节所描述的特殊行阶梯形:
\[
A_{f, X, Y} = \begin{pmatrix} 1 & 1 & 2 & 4 & 8 \\ 1 & 2 & 1 & 2 & 1 \\ 2 & 2 & 2 & 2 & 1 \\ 1 & 2 & 2 & 2 & 2 \end{pmatrix}
\longmapsto \begin{pmatrix} 1 & 1 & 2 & 4 & 8 \\ 0 & 1 & -1 & -2 & -7 \\ 0 & 0 & -2 & -6 & -15 \\ 0 & 1 & 0 & -2 & -6 \end{pmatrix}
\longmapsto \begin{pmatrix} 1 & 1 & 2 & 4 & 8 \\ 0 & 1 & -1 & -2 & -7 \\ 0 & 0 & 2 & 6 & 15 \\ 0 & 0 & 1 & 0 & 1 \end{pmatrix}
\]
\[
\longmapsto \begin{pmatrix} 1 & 1 & 2 & 4 & 8 \\ 0 & 1 & -1 & -2 & -7 \\ 0 & 0 & 1 & 0 & 1 \\ 0 & 0 & 0 & 6 & 13 \end{pmatrix}
\longmapsto \begin{pmatrix} 1 & 1 & 2 & 4 & 8 \\ 0 & 1 & -1 & -2 & -7 \\ 0 & 0 & 1 & 0 & 1 \\ 0 & 0 & 0 & 1 & \frac{13}{6} \end{pmatrix}
\longmapsto \begin{pmatrix} 1 & 1 & 2 & 0 & -\frac{2}{3} \\ 0 & 1 & -1 & 0 & -\frac{8}{3} \\ 0 & 0 & 1 & 0 & 1 \\ 0 & 0 & 0 & 1 & \frac{13}{6} \end{pmatrix}
\]
\[
\longmapsto \begin{pmatrix} 1 & 1 & 0 & 0 & -\frac{8}{3} \\ 0 & 1 & 0 & 0 & -\frac{5}{3} \\ 0 & 0 & 1 & 0 & 1 \\ 0 & 0 & 0 & 1 & \frac{13}{6} \end{pmatrix}
\longmapsto \begin{pmatrix} 1 & 0 & 0 & 0 & -1 \\ 0 & 1 & 0 & 0 & -\frac{5}{3} \\ 0 & 0 & 1 & 0 & 1 \\ 0 & 0 & 0 & 1 & \frac{13}{6} \end{pmatrix}
\]
按照 \ref{sec:3.5} 节, 向量 $(1, \frac{5}{3}, -1, -\frac{13}{6}, 1)^t \in \mathbb{R}^5$ 因此构成解空间
$M_{A_{f, X, Y}, 0}$ 的一个基.相应地, 对 $X = \{x_1, \ldots, x_5\}$,
\[
x_1 + \frac{5}{3}x_2 - x_3 - \frac{13}{6}x_4 + x_5 \in V
\]
给出 $\ker f$ 的一个基.

\subsection*{习题 3.5/6}

\noindent\textbf{起点:} 给定一个矩阵 $A \in K^{m \times n}$, 一个列向量 $b \in K^m$ 以及一个
可逆矩阵 $S \in \mathrm{Gl}(n, K)$.

\noindent\textbf{目标:} 要对线性方程组 $Ax = b$ 与 $ASy = b$ 的解空间证明
$M_{A, b} = SM_{AS, b}$.

\noindent\textbf{策略:} 目标是有道理的.事实上, 若 $y \in K^n$ 满足方程 $ASy = b$, 则可以把
它解释成形式 $(AS)y = b$ 或 $A(Sy) = b$.于是得到 $y \in M_{AS, b}$, 但还有
$Sy \in M_{A, b}$.

\noindent\textbf{解题路径:} 先设 $x \in SM_{AS, b}$, 即 $x = Sy$ 且 $y \in M_{AS, b}$.则
$ASy = b$, 因而 $x = Sy \in M_{A, b}$.反之, 设 $x \in M_{A, b}$, 即 $Ax = b$.由于 $S$ 可逆,
存在 $y \in K^n$ 使得 $x = Sy$.则 $ASy = b$, 即 $y \in M_{AS, b}$, 从而 $x = Sy \in SM_{AS, b}$.
总之得到 $M_{A, b} = SM_{AS, b}$.

\subsection*{习题 3.5/8}

\noindent\textbf{起点:} 要考虑有限维 $K$-向量空间之间的 $K$-线性映射
$f \colon V \longrightarrow W$, 以及相应的对偶映射 $f^* \colon W^* \longrightarrow V^*$.其中
$V^*$(相应地 $W^*$)是 $V$(相应地 $W$)的对偶空间, 即全体 $K$-线性映射
$V \longrightarrow K$(相应地 $W \longrightarrow K$)构成的 $K$-向量空间, 而 $f^*$ 由
$\varphi \longmapsto \varphi \circ f$ 定义.此外还要考察一个向量 $b \in W$.

\noindent\textbf{目标:} 要证明: 恰当存在向量 $x \in V$ 使得 $f(x) = b$, 即恰当 $b \in \operatorname{im} f$,
当且仅当对所有 $\varphi \in \ker f^*$ 有 $\varphi(b) = 0$.

\noindent\textbf{策略:} 首先我们要更精确地解释对偶映射 $f^*$ 的核.对偶映射
$f^* \colon W^* \longrightarrow V^*$ 由 $f^*(\varphi) = \varphi \circ f$($\varphi \in W^*$)定义.
因此 $f^*(\varphi) = 0$ 等价于 $\varphi \circ f = 0$, 即等价于 $\operatorname{im} f \subset \ker \varphi$.
于是, 若 $b \in \operatorname{im} f$, 则对所有 $\varphi \in \ker f^*$ 有 $\varphi(b) = 0$.换言之,
\begin{equation}
\operatorname{im} f \subset \bigcap_{\varphi \in \ker f^*} \ker \varphi\text{.} \tag{$*$}
\end{equation}
为达到我们的目标, 要证明上面的包含在反方向也成立, 从而是一个等式.

\noindent\textbf{解题路径:} 先设 $b \in \operatorname{im} f$.我们已经说明, 此时对所有满足
$\operatorname{im} f \subset \ker \varphi$, 即满足 $\varphi \circ f = 0$ 的线性形式
$\varphi \in W^*$ 有 $\varphi(b) = 0$, 而后者等价于 $\varphi \in \ker f^*$.

反之, 设 $b \notin \operatorname{im} f$.我们要证明, 在此前提下存在满足 $\varphi(b) \neq 0$ 的
线性形式 $\varphi \in \ker f^*$.具体地说, 我们选取 $\operatorname{im} f$ 的一个基, 并通过加入
$b$ 把它扩大为 $W$ 中的一个线性无关组; 由于 $b \notin \operatorname{im} f$, 这是可能的, 参见
引理~\mtref{lem:1.5.5}.进一步, 我们借助定理~\mtref{thm:1.5.8} 把这个组扩充为 $W$ 的一个基 $Y$.
然后按照定理~\mtref{thm:2.1.7} (ii), 可以由
\[
\varphi(y) = \begin{cases} 1, & \text{对 } y = b, \\ 0, & \text{对 } y \in Y, y \neq b \end{cases}
\]
及线性扩张定义一个线性形式 $\varphi \colon W \longrightarrow K$.按构造有
$\varphi(\operatorname{im} f) = 0$, 即 $\varphi \in \ker f^*$, 以及 $\varphi(b) \neq 0$.
因此, 若 $b \in W$ 不属于 $\operatorname{im} f$, 则也不属于 $\bigcap_{\varphi \in \ker f^*} \ker \varphi$,
由此推出上面的包含 $(*)$ 是一个等式.








% 附录: 习题训练 -- 对应第 4 章(4.1-4.5: 行列式)的解答.
% 源文档 ap4_src.txt; 术语见 CONVENTIONS.md; 交叉引用标签见 04-行列式.tex 中的实际标签.

\section*{对应 4.1-4.5: 行列式}

\subsection*{习题 4.1/2}

\noindent\textbf{起点:} 需考察 $\mathfrak{S}_n$ 中的一个 $r$-循环, 即一个置换
$\pi \in \mathfrak{S}_n$, 对它存在两两不同的元素 $a_1, \ldots, a_r \in
\{1, \ldots, n\}$, 使得对 $i = 1, \ldots, r - 1$ 有 $\pi(a_i) = a_{i+1}$,
$\pi(a_r) = a_1$, 且对 $\{1, \ldots, n\}$ 的所有其余元素 $a$ 有
$\pi(a) = a$.我们简写为 $\pi = (a_1, \ldots, a_r)$, 其中必然有
$n \geq r \geq 2$.

\noindent\textbf{目标:} 需确定 $\sgn \pi$.

\noindent\textbf{策略:} 我们考察 $\mathfrak{S}_n$ 中一个 $r$-循环的简单例子, 即
\[
\pi = (1, \ldots, r) = \begin{pmatrix} 1 & 2 & \ldots & r - 1 & r & r + 1 & \ldots & n \\ 2 & 3 & \ldots & r & 1 & r + 1 & \ldots & n \end{pmatrix}\text{.}
\]
显然, 序列 $\pi(1), \ldots, \pi(n)$ 中逆序的个数恰好是 $r - 1$, 从而借助
定理~\mtref{thm:4.1.6} 得出 $\sgn \pi = (-1)^{r-1}$.若 $\pi = (a_1, \ldots, a_r)$
是 $\mathfrak{S}_n$ 中一个任意的 $r$-循环, 则同样应有
$\sgn \pi = (-1)^{r-1}$.事实上, 若把 $\pi$ 写成对换的乘积, 例如
$\pi = \tau_1 \circ \cdots \circ \tau_s$, 则按推论~\mtref{cor:4.1.8} 有
$\sgn \pi = (-1)^s$.不过这样的乘积表示并不依赖于 $n$ 元素集合
$\{1, \ldots, n\}$ 的特殊计数.通过选取 $\{a_1, \ldots, a_r, a_{r+1}, \ldots, a_n\}$
类型的一个计数, 其中 $a_{r+1}, \ldots, a_n \in \{1, \ldots, n\}$ 适当选取, 我们
可以像上面那样假设 $\pi = (1, \ldots, r)$.于是对任意的 $r$-循环
$\pi \in \mathfrak{S}_n$ 得出 $\sgn \pi = (-1)^{r-1}$.

\noindent\textbf{解题路径:} 刚才使用的论证可以精确化, 为此利用置换
\[
\sigma = \begin{pmatrix} 1 & \ldots & n \\ a_1 & \ldots & a_n \end{pmatrix}\text{.}
\]
具体来说有 $\sigma^{-1} \circ \pi \circ \sigma = (1, \ldots, r)$, 这是容易验证的,
因此
\[
\sgn \pi = \sgn \sigma^{-1} \cdot \sgn \pi \cdot \sgn \sigma = \sgn(1, \ldots, r) = (-1)^{r-1}\text{.}
\]

\noindent\textbf{补充:} 在策略的思考中, 我们曾用 $\pi$ 作为对换的乘积的分解来论证.
这自然引出问题: 对 $r$-循环 $\pi = (a_1, \ldots, a_r) \in \mathfrak{S}_n$, 是否存在
一般类型的对换分解, 使得 $\pi$ 的符号能借助推论~\mtref{cor:4.1.8} 来确定.这条路确实
可行, 因为
\[
\pi = (a_1, \ldots, a_r) = (a_1, a_r) \circ \cdots \circ (a_1, a_2),
\]
从而得出 $\sgn \pi = (-1)^{r-1}$.

\subsection*{习题 4.1/4}

\noindent\textbf{起点:} 给定一个正规子群 $N \subset \mathfrak{S}_n$, 它含有一个对换
$\tau$, 例如 $\tau = (i, j)$.

\noindent\textbf{目标:} 需证明 $N = \mathfrak{S}_n$.

\noindent\textbf{策略:} 正规子群 $N \subset \mathfrak{S}_n$ 是一个具有下述性质的子群:
由 $\pi \in N$ 与 $\sigma \in \mathfrak{S}_n$ 恒可推出
$\sigma^{-1} \circ \pi \circ \sigma \in N$, 或者说, 如果用 $\sigma$ 替换
$\sigma^{-1}$, 则 $\sigma \circ \pi \circ \sigma^{-1} \in N$.在我们的情形下, 我们
知道对换 $\tau = (i, j)$ 属于 $N$.因此对所有 $\sigma \in \mathfrak{S}_n$ 有
$\sigma \circ (i, j) \circ \sigma^{-1} \in N$.但 $\sigma \circ (i, j) \circ
\sigma^{-1}$ 把元素 $\sigma(i)$ 映到 $\sigma(j)$, 把元素 $\sigma(j)$ 映到
$\sigma(i)$, 并且保持 $\{1, \ldots, n\}$ 的一切其他元素不动.换言之,
\[
\sigma \circ (i, j) \circ \sigma^{-1} = (\sigma(i), \sigma(j))
\]
是另一个属于 $N$ 的对换.正如我们将看到的, 这就是证明 $N = \mathfrak{S}_n$ 的关键.

\noindent\textbf{解题路径:} 我们刚才已经说明, 除对换 $(i, j)$ 外, 对置换
$\sigma \in \mathfrak{S}_n$, 一切 $(\sigma(i), \sigma(j))$ 类型的对换也都属于
$N$.但如今对任意两个预先给定的不同数 $i', j' \in \{1, \ldots, n\}$, 总能找到一个
置换 $\sigma \in \mathfrak{S}_n$ 使得 $\sigma(i) = i'$ 且 $\sigma(j) = j'$.因此
得到
\[
(i', j') = (\sigma(i), \sigma(j)) \in N,
\]
并且我们看到 $N$ 包含 $\mathfrak{S}_n$ 的全部对换.由于按
定理~\mtref{thm:4.1.3}, 每个置换 $\pi \in \mathfrak{S}_n$ 都是对换的乘积, 又由于 $N$
具有子群的性质, 最后推出 $\pi \in N$, 从而 $N = \mathfrak{S}_n$.

\subsection*{习题 4.2/2}

\noindent\textbf{起点:} 给定一个维数为 $n$ 的 $K$-向量空间 $V$, 其上有一个非平凡的
行列式函数 $\varDelta \colon V^n \longrightarrow K$, 以及一个线性无关的向量组
$x_1, \ldots, x_{n-1} \in V$.

\noindent\textbf{目标:} 存在 $x_n \in V$ 使得 $\varDelta(x_1, \ldots, x_n) = 1$.其中
要证明, $x_n$ 在 $V/\langle x_1, \ldots, x_{n-1}\rangle$ 中的剩余类被唯一确定.

\noindent\textbf{策略:} 我们先满足于找到一个向量 $x_n' \in V$ 使
$\varDelta(x_1, \ldots, x_{n-1}, x_n') \neq 0$.为此按定理~\mtref{thm:4.2.7}, 只需把
向量组 $x_1, \ldots, x_{n-1}$ 用向量 $x_n'$ 扩充为线性无关的向量组, 从而扩充为 $V$
的一个基.于是有 $\alpha := \varDelta(x_1, \ldots, x_{n-1}, x_n') \neq 0$, 并且对
$x_n = \alpha^{-1} x_n'$, 利用 $\varDelta$ 的多重线性性, 我们如愿得到
$\varDelta(x_1, \ldots, x_n) = 1$.

按构造, $x_n$ 的剩余类 $\overline{x}_n$ 给出 $\overline{V} = V/\langle x_1, \ldots,
x_{n-1}\rangle$ 的一个基.对另一个向量 $y_n \in V$ 及其剩余类
$\overline{y}_n \in \overline{V}$, 于是存在常数 $\alpha_n \in K$ 使
$\overline{y}_n = \alpha_n \overline{x}_n$.于是有 $y_n - \alpha_n x_n \in
\langle x_1, \ldots, x_{n-1}\rangle$, 并且存在常数 $\alpha_1, \ldots, \alpha_{n-1}
\in K$ 使得
\[
y_n - \alpha_n x_n = \sum_{i=1}^{n-1} \alpha_i x_i\text{.}
\]
由计算
\[
\varDelta(x_1, \ldots, x_{n-1}, y_n) = \sum_{i=1}^{n} \alpha_i \cdot \varDelta(x_1,
\ldots, x_{n-1}, x_i) = \alpha_n \cdot \varDelta(x_1, \ldots, x_n) = \alpha_n
\]
于是得出: 若 $\varDelta(x_1, \ldots, x_{n-1}, y_n) = 1$, 则
$\overline{y}_n = \overline{x}_n$.

\noindent\textbf{解题路径:} 从策略思考中我们已能读出完整解答.特别是我们说明了如何找到满足 $\varDelta(x_1, \ldots, x_n) = 1$ 的向量 $x_n \in V$.关于 $x_n$ 模子空间 $\langle x_1, \ldots, x_{n-1}\rangle$ 唯一性的论证还可略加简化.事实上, 若对另一个向量 $y_n \in V$ 有 $\varDelta(x_1, \ldots, x_{n-1}, y_n) = 1$, 则
\[
\varDelta(x_1, \ldots, x_{n-1}, x_n - y_n) = \varDelta(x_1, \ldots, x_{n-1}, x_n)
- \varDelta(x_1, \ldots, x_{n-1}, y_n) = 0\text{.}
\]
因此我们借助定理~\mtref{thm:4.2.7} 看到, 向量 $x_1, \ldots, x_{n-1}, x_n - y_n$
线性相关.但由于 $x_1, \ldots, x_{n-1}$ 构成一个线性无关的向量组, 必然有
$x_n - y_n \in \langle x_1, \ldots, x_{n-1}\rangle$, 参见
引理~\mtref{lem:1.5.5}, 于是 $x_n$ 与 $y_n$ 在 $V/\langle x_1, \ldots,
x_{n-1}\rangle$ 中的剩余类一致.

\subsection*{习题 4.3/3}

\noindent\textbf{起点:} 需考察一个矩阵 $A = (\alpha_{ij}) \in K^{n \times n}$, 其中对
$i < j$ 有 $\alpha_{ij} = 1$, 对 $i = j$ 有 $\alpha_{ij} = 0$, 对 $i > j$ 有
$\alpha_{ij} = -1$; 为准确起见, 我们用 $A_n$ 代替 $A$.于是 $A_n$ 是一个 $n$ 阶方阵,
其系数在主对角线以上为 $1$, 在主对角线以下为 $-1$, 而在主对角线上为 $0$.

\noindent\textbf{目标:} 对偶数 $n \in \mathbb{N}$, 需证明 $\det(A_n) = 1$.

\noindent\textbf{策略:} 由于矩阵 $A_n$ 不属于 \ref{sec:4.3} 节所处理的特殊例子之列, 又由于
引理~\mtref{lem:4.2.4} 中的定义公式也不能提供概览, 因此计算 $\det(A_n)$ 宜采用
定理~\mtref{thm:4.3.6} 的初等行变换或列变换方法.此时应首先尝试对 $n$ 作归纳证明,
归纳步骤由 $n - 2$ 到 $n$, 或者说由 $n$ 到 $n + 2$, 因为只需考虑偶数的行数与列数.
此外, \ref{sec:4.3} 节的习题 1 让人推测, 对奇数 $n$, $\det(A_n)$ 的性状在这一关联中不会有任何
用处.

情形 $n = 0$ 是平凡的, 对 $n = 2$ 我们得到
\[
\det(A_2) = \det \begin{pmatrix} 0 & 1 \\ -1 & 0 \end{pmatrix} = -(-1) \cdot 1 = 1\text{.}
\]
在 $n = 4$ 的情形, 我们按标准方式用初等行变换如下变换 $A_n$:
\[
A_4 = \begin{pmatrix} 0 & 1 & 1 & 1 \\ -1 & 0 & 1 & 1 \\ -1 & -1 & 0 & 1 \\ -1 & -1 & -1 & 0 \end{pmatrix} \longmapsto^{(1)} \begin{pmatrix} -1 & 0 & 1 & 1 \\ 0 & 1 & 1 & 1 \\ -1 & -1 & 0 & 1 \\ -1 & -1 & -1 & 0 \end{pmatrix}
\]
\[
\longmapsto^{(2)} \begin{pmatrix} -1 & 0 & 1 & 1 \\ 0 & 1 & 1 & 1 \\ 0 & -1 & -1 & 0 \\ 0 & -1 & -2 & -1 \end{pmatrix} \longmapsto^{(3)} \begin{pmatrix} -1 & 0 & 1 & 1 \\ 0 & 1 & 1 & 1 \\ 0 & 0 & 0 & 1 \\ 0 & 0 & -1 & 0 \end{pmatrix} = \begin{pmatrix} \begin{matrix} -1 & 0 \\ 0 & 1 \end{matrix} & * \\ 0 & A_2 \end{pmatrix}\text{.}
\]
其中我们在步骤 (1) 中把前两行互相交换, 在步骤 (2) 中从第三行与第四行减去第一行,
最后在步骤 (3) 中把第二行加到第三行与第四行上.由于行列式在步骤 (1) 中按
定理~\mtref{thm:4.3.6} 改变因子 $-1$, 此外保持不变, 因此按 \ref{sec:4.3} 节的
例~\mtref{bsp:4.3.2} 把 $\det(A_4)$ 计算为
\[
\det(A_4) = -\det \begin{pmatrix} -1 & 0 \\ 0 & 1 \end{pmatrix} \cdot \det(A_2) = 1 \cdot 1 = 1\text{.}
\]
原则上, 对一般的偶数 $n$ 也可以进行类似的论证.

\noindent\textbf{解题路径:} 对偶数 $n \geq 4$, 可以如下用初等行变换变换 $A_n$:
\[
A_n = \begin{pmatrix} 0 & 1 & 1 & \cdots & 1 & 1 \\ -1 & 0 & 1 & \cdots & 1 & 1 \\ -1 & -1 & 0 & \cdots & 1 & 1 \\ \vdots & \vdots & \vdots & \cdots & \vdots & \vdots \\ -1 & -1 & -1 & \cdots & -1 & 0 \end{pmatrix} \longmapsto \begin{pmatrix} 0 & 1 & 1 & \cdots & 1 & 1 \\ -1 & 0 & 1 & \cdots & 1 & 1 \\ 0 & 0 & 0 & \cdots & 1 & 1 \\ \vdots & \vdots & \vdots & \cdots & \vdots & \vdots \\ 0 & 0 & -1 & \cdots & -1 & 0 \end{pmatrix} = \begin{pmatrix} A_2 & * \\ 0 & A_{n-2} \end{pmatrix}\text{.}
\]
具体来说, 在 $A_n$ 中我们把它下方 $n - 2$ 行减去第二行, 随后把第一行同样加到下方
$n - 2$ 行上.换句话说, 我们把第一行与第二行的差, 即行向量 $(1, 1, 0, \ldots, 0, 0)$,
加到下方 $n - 2$ 行上.于是按 \ref{sec:4.3} 节的例~\mtref{bsp:4.3.2} 得出
\[
\det(A_n) = \det(A_2) \cdot \det(A_{n-2}),
\]
从而利用 $\det(A_2) = 1$, 可以归纳地推出: 对偶数 $n$ 有 $\det(A_n) = 1$.

\subsection*{习题 4.3/5}

\noindent\textbf{起点:} 给定一个 $n$ 维 $K$-向量空间 $V$, 其上有一个自同态
$f \colon V \longrightarrow V$ 和一个向量 $x \in V$, 使得向量 $x, f(x), \ldots,
f^{n-1}(x)$ 构成 $V$ 的生成系.特别地, $f^n(x)$ 于是可以表示为这个向量组的线性
组合, 即存在系数 $\alpha_0, \ldots, \alpha_{n-1} \in K$ 使得
\[
f^n(x) = \sum_{i=0}^{n-1} \alpha_i f^i(x)\text{.}
\]

\noindent\textbf{目标:} 需证明 $\det(f) = (-1)^{n+1} \alpha_0$.

\noindent\textbf{策略:} 为计算 $f$ 的行列式, 按定义~\mtref{def:4.3.2} 我们可以确定一个
描述矩阵 $A_{f, X, X}$ 的行列式, 其中 $X$ 是 $V$ 的任意一个基.但在 $V$ 中我们只知道
一个生成系, 即 $X = (x, f(x), \ldots, f^{n-1}(x))$.由于它的长度为 $n = \dim_K V$,
它也是线性无关的, 从而是 $V$ 的一个基.事实上, 若 $X$ 线性相关, 则可以把 $X$ 缩减为
一个极小的生成系, 从而按定理~\mtref{thm:1.5.7} 缩减为 $V$ 的一个基.但这会推出
$\dim_K V < n$, 与 $\dim V = n$ 矛盾.因此 $X$ 是 $V$ 的一个基, 并且只需确定矩阵
\[
A_{f, X, X} = \begin{pmatrix} 0 & & & & & \alpha_0 \\ 1 & 0 & & & & \alpha_1 \\ & 1 & \cdot & & & \alpha_2 \\ & & \cdot & \cdot & & \vdots \\ & & & \cdot & \cdot & \vdots \\ & & & & 0 & \alpha_{n-2} \\ & & & & 1 & \alpha_{n-1} \end{pmatrix}
\]
的行列式.简言之, 这个矩阵是
\[
A_{f, X, X} = \begin{pmatrix} 0 & \alpha_0 \\ E_{n-1} & * \end{pmatrix}
\]
类型, 其中 $E_{n-1}$ 是 $(n-1)$ 阶单位矩阵, 并且自然的做法是像 \ref{sec:4.3} 节的
例~\mtref{bsp:4.3.2} 那样计算相应的行列式.

\noindent\textbf{解题路径:} 如我们所看到的, 对上述矩阵 $A_{f, X, X}$ 有
$\det(f) = \det(A_{f, X, X})$.为计算这个矩阵的行列式, 我们使用初等列变换.具体来说,
把 $A_{f, X, X}$ 的第 $n$ 列与前面 $n - 1$ 列交换.于是按 \ref{sec:4.3} 节的
例~\mtref{bsp:4.3.2} 得到
\[
\det(f) = \det \begin{pmatrix} 0 & \alpha_0 \\ E_{n-1} & * \end{pmatrix} = (-1)^{n-1} \det \begin{pmatrix} \alpha_0 & 0 \\ * & E_{n-1} \end{pmatrix} = (-1)^{n+1} \alpha_0\text{.}
\]

\noindent\textbf{补充:} 按定义~\mtref{def:6.5.1} 的术语, 在自同态 $f \colon V
\longrightarrow V$ 下, $V$ 是一个所谓的 $f$-循环向量空间.当人们想研究有限维向量空间
的自同态的结构时, 这类向量空间具有根本性的意义; 参见 \ref{sec:6.5} 节.

\subsection*{习题 4.4/2}

\noindent\textbf{起点:} 给定一个可逆矩阵 $A = (\alpha_{ij}) \in \mathrm{Gl}(n,
\mathbb{R})$, 其系数 $\alpha_{ij} \in \mathbb{Z}$.

\noindent\textbf{目标:} 需证明逆矩阵 $A^{-1}$ 的系数属于 $\mathbb{Q}$.此外还要证明,
$A^{-1}$ 的系数属于 $\mathbb{Z}$, 当且仅当 $\det(A) = \pm 1$.

\noindent\textbf{策略:} 我们首先考虑, 按推论~\mtref{cor:3.3.7} 的矩阵可逆性的基本刻画能
走多远.由于 $A$ 作为 $\mathbb{R}^{n \times n}$ 的元素是可逆的, 按
推论~\mtref{cor:3.3.7}, $A$ 的行(或列)在 $\mathbb{R}$ 上线性无关, 从而特别地也在
$\mathbb{Q}$ 上线性无关.若把 $A$ 视为 $\mathbb{Q}^{n \times n}$ 中的矩阵, 则再次利用
推论~\mtref{cor:3.3.7} 得到, $A$ 在 $\mathbb{Q}^{n \times n}$ 中也可逆.但如今逆元素被
唯一确定.因此, 在 $\mathbb{Q}^{n \times n}$ 与 $\mathbb{R}^{n \times n}$ 中对 $A$
形成的逆矩阵一致.所以 $A^{-1}$ 作为 $A \in \mathrm{Gl}(n, \mathbb{R})$ 的逆矩阵,
其系数属于 $\mathbb{Q}$.

在上述方法中, 行列式 $\det(A)$ 没有起作用.这让人推测, 为回答更深入的问题, 即逆矩阵
$A^{-1}$ 何时具有 $\mathbb{Z}$ 中的系数, 我们不得不求助于
推论~\mtref{cor:4.4.5} 的公式, 从而求助于克拉默法则.

\noindent\textbf{解题路径:} 伴随矩阵 $A^{\mathrm{ad}}$ 的定义是 $A$ 的代数余子式的
转置矩阵, 由此看出 $A^{\mathrm{ad}}$ 的系数属于 $\mathbb{Z}$, 因为这对 $A$ 也成立.
此外, 由引理~\mtref{lem:4.2.4} 的定义有 $\det(A) \in \mathbb{Z}$.因此我们从
推论~\mtref{cor:4.4.5} 的等式
\[
A^{-1} = \det(A)^{-1} \cdot A^{\mathrm{ad}}
\]
读出: $A^{-1}$ 的系数属于 $\mathbb{Q}$, 甚至属于 $\mathbb{Z}$, 如果 $\det(A)$ 在
$\mathbb{Z}$ 中可逆, 即 $\det(A) = \pm 1$.反过来, 若 $A$ 与 $A^{-1}$ 的系数都属于
$\mathbb{Z}$, 则
\[
\det(A) \in \mathbb{Z}, \qquad \det(A)^{-1} = \det(A^{-1}) \in \mathbb{Z},
\]
于是 $\det(A)$ 是 $\mathbb{Z}$ 中的一个单位, 而这意味着 $\det(A) = \pm 1$.

\subsection*{习题 4.4/4}

\noindent\textbf{起点:} 给定一个方程 $\sum_{i=1}^{n} \alpha_i a_i = b$, 其中系数
$\alpha_i \in K$, 列向量 $a_i, b \in K^n$.

\noindent\textbf{目标:} 对 $i = 1, \ldots, n$, 需证明关系式
\[
\det(a_1, \ldots, a_{i-1}, b, a_{i+1}, \ldots, a_n) = \alpha_i \cdot \det(a_1, \ldots, a_n),
\]
并由此推出我们在推论~\mtref{cor:4.4.6} 中借助克拉默法则得到的线性方程组的求解公式.

\noindent\textbf{策略:} $b$ 表示为 $a_1, \ldots, a_n$ 的线性组合.把这个线性组合代入
$b$ 的位置, 并利用行列式在列向量中是多重线性和交替的这一事实, 则如所愿得到
\[
\begin{array}{c}
\det(a_1, \ldots, a_{i-1}, b, a_{i+1}, \ldots, a_n) \\
= \sum_{j=1}^{n} \alpha_j \det(a_1, \ldots, a_{i-1}, a_j, a_{i+1}, \ldots, a_n) = \alpha_i \det(a_1, \ldots, a_n)\text{.}
\end{array}
\]
现在需要在推论~\mtref{cor:4.4.6} 的框架下解释这个等式.

\noindent\textbf{解题路径:} 公式的推导我们刚才已经描述了.然后, 对矩阵 $A = (a_1,
\ldots, a_n) \in K^{n \times n}$ 和 $b \in K^n$, 考虑相应的线性方程组 $Ax = b$.
方程 $\sum_{i=1}^{n} \alpha_i a_i = b$ 无非是说, 列向量 $z = (\alpha_1, \ldots,
\alpha_n)^t$ 是这个方程组的一个解.在推论~\mtref{cor:4.4.6} 的情形中 $A$ 被假定为可逆,
从而 $z = A^{-1} b$ 是方程组 $Ax = b$ 的一个解, 而且是唯一可能的解.注意到
$\det(A) \neq 0$, 由上面的公式得出解 $z$ 的第 $i$ 个分量 $\alpha_i$ 为
\[
\alpha_i = \det(A)^{-1} \cdot \det(a_1, \ldots, a_{i-1}, b, a_{i+1}, \ldots, a_n),
\]
与推论~\mtref{cor:4.4.6} 一致.

\subsection*{习题 4.5/2}

\noindent\textbf{起点:} 需考察一个 $K$-向量空间 $V$, 其中有元素 $a_1, \ldots, a_r
\in V$.

\noindent\textbf{目标:} 需证明: 向量 $a_1, \ldots, a_r$ 线性无关, 当且仅当
$a_1 \wedge \cdots \wedge a_r \neq 0$.

\noindent\textbf{策略:} 先设向量 $a_1, \ldots, a_r$ 线性相关.则考察典范的交替多重线性
映射
\[
\mathrm{kan} \colon V^r \longrightarrow \bigwedge^r V, \quad (x_1, \ldots, x_r) \longmapsto x_1 \wedge \cdots \wedge x_r,
\]
并利用注记~\mtref{rem:4.5.1} 推出 $a_1 \wedge \cdots \wedge a_r = 0$.因此只剩下对线性
无关的向量组 $a_1, \ldots, a_r$ 证明 $a_1 \wedge \cdots \wedge a_r \neq 0$.

为能在 $\bigwedge^r V$ 中计算, 我们首先假设 $V$ 是有限维的并带有基 $X$.则按
定理~\mtref{thm:4.5.4} 存在 $\bigwedge^r V$ 的一个基, 它由元素 $x_H$ 组成, 其中 $H$
遍历基 $X$ 的一切 $r$ 元素子集.此外有
\begin{equation}
a_1 \wedge \cdots \wedge a_r = \sum_H \det_{X, H}(a_1, \ldots, a_r) x_H \tag{$*$}
\end{equation}
其中 $\det_{X, H}(a_1, \ldots, a_r)$ 是从 $r$ 列矩阵 $A = (a_{1, X}, \ldots, a_{r, X})$
中删去一切其指标不与 $H$ 的元素对应的行而得到的\emph{子矩阵的行列式}(\textup{Determinante einer Untermatrix})\index{zijuzhendehanglieshi@子矩阵的行列式(Determinante einer Untermatrix)}.若 $a_1, \ldots, a_r$
线性无关, 则 $\mathrm{rg}\, A = r$, 并且存在 $A$ 的一个 $(r \times r)$ 子矩阵, 其秩
为 $r$, 从而其行列式不为零.特别地, 表示式 $(*)$ 中相应的系数
$\det_{X, H}(a_1, \ldots, a_r)$ 是非平凡的, 于是推出 $a_1 \wedge \cdots \wedge a_r
\neq 0$.我们甚至还可以更有效地论证, 方法是把向量 $a_1, \ldots, a_r$ 扩充为 $V$ 的
一个基.此时 $a_1 \wedge \cdots \wedge a_r$ 属于定理~\mtref{thm:4.5.4} 中构造的
$\bigwedge^r V$ 的基, 从而是非平凡的.原则上, 这个论证对任意的 $K$-向量空间 $V$
仍然有效, 尽管此时定理~\mtref{thm:4.5.4} 的结果首先需要推广到无限维向量空间.

\noindent\textbf{解题路径:} 不过我们想在此更直接地进行.首先像上面那样论证: 由
$a_1 \wedge \cdots \wedge a_r \neq 0$ 恒可推出 $a_1, \ldots, a_r$ 线性无关.反过来,
设 $a_1, \ldots, a_r$ 线性无关, 则对 $U = \langle a_1, \ldots, a_r \rangle$ 构造一个
直和补空间 $U'$, 例如把 $a_1, \ldots, a_r$ 扩充为 $V$ 的一个基; 其中需要把基扩充定理
~\mtref{thm:1.5.8} 借助佐恩引理~\mtref{lem:1.5.15} 推广到无限维的情形.然后考虑到第一个
直和项的投影
\[
\pi \colon V = U \oplus U' \longrightarrow U,
\]
并按定理~\mtref{thm:4.2.8} 选取一个非平凡的行列式函数 $\varDelta \colon U^r
\longrightarrow K$, 其中 $\varDelta(a_1, \ldots, a_r) \neq 0$; 参见
定理~\mtref{thm:4.2.7}.则
\[
\Phi \colon V^r \longrightarrow K, \quad (x_1, \ldots, x_r) \longmapsto \varDelta(\pi(x_1), \ldots, \pi(x_r))
\]
是一个交替多重线性映射.根据 $V$ 的 $r$ 重外积的泛性质, $\Phi$ 经由一个 $K$-线性映射
$\varphi \colon \bigwedge^r V \longrightarrow K$ 如下分解:

\begin{center}
\begin{tikzcd}[row sep=huge, column sep=huge]
V^r \arrow[rr, "\mathrm{kan}"] \arrow[dr, "\Phi"'] && \bigwedge^r V \arrow[dl, "\varphi"] \\
& K
\end{tikzcd}
\end{center}

其中 $\varphi(a_1 \wedge \cdots \wedge a_r) = \Phi(a_1, \ldots, a_r) =
\varDelta(a_1, \ldots, a_r) \neq 0$, 特别地 $a_1 \wedge \cdots \wedge a_r \neq 0$,
这正是要证明的.

\subsection*{习题 4.5/5}

\noindent\textbf{起点:} 需考察特征为 $0$ 的域 $K$ 上的向量空间 $V$ 的二重外积
$\bigwedge^2 V$.对 $z \in \bigwedge^2 V$, 秩 $\mathrm{rg}\, z$ 定义为: 使得存在分解
\[
z = \sum_{i=1}^{r} x_i \wedge y_i
\]
(其中向量 $x_i, y_i \in V$) 的一切 $r \in \mathbb{N}$ 中的最小值.

\noindent\textbf{目标:} 需证明: 数 $r = \mathrm{rg}\, z$ 由关系式 $z^r \neq 0$ 与
$z^{r+1} = 0$ 唯一刻画, 其中这些幂在外代数 $\bigwedge V$ 中取.

\noindent\textbf{策略:} 设 $z = \sum_{i=1}^{r} x_i \wedge y_i \in \bigwedge^2 V$, 其中
向量 $x_i, y_i \in V$.我们想首先获得对幂 $z^s$($s \in \mathbb{N}$) 的性状的直观
感受.事实上成立
\[
z^s = \sum_{i_1, \ldots, i_s = 1}^{r} (x_{i_1} \wedge y_{i_1}) \wedge \cdots \wedge (x_{i_s} \wedge y_{i_s}),
\]
其中右端一切其指标 $i_1, \ldots, i_s \in \{1, \ldots, r\}$ 不两两不同的乘积都消失.
由此已经推出 $z^s = 0$ 对 $s > r$ 成立.对 $s = r$, 至多乘积
\[
(x_{i_1} \wedge y_{i_1}) \wedge \cdots \wedge (x_{i_r} \wedge y_{i_r})
\]
可以非零, 其中指标 $i_1, \ldots, i_r$ 构成 $1, \ldots, r$ 的一个置换.这些保留下来的
乘积可以通过交换因子化为自然顺序
\[
(x_1 \wedge y_1) \wedge \cdots \wedge (x_r \wedge y_r)\text{.}
\]
由于此时需要偶数次交换, 符号不改变.注意到置换群 $\mathfrak{S}_r$ 恰好有 $r!$ 个元素,
幂 $z^r$ 计算为
\[
z^r = r! \cdot (x_1 \wedge y_1) \wedge \cdots \wedge (x_r \wedge y_r)\text{.}
\]
现在 $K$ 的特征为 $0$.因此 $z^r$ 消失, 当且仅当乘积 $(x_1 \wedge y_1) \wedge \cdots
\wedge (x_r \wedge y_r)$ 消失, 即利用 \ref{sec:4.5} 节的习题 2, 当且仅当向量 $x_1, y_1, \ldots,
x_r, y_r$ 线性相关.这样, 对 $z = \sum_{i=1}^{r} x_i \wedge y_i \in \bigwedge^2 V$
我们认清了下列性质:
\begin{equation}
z^{r+1} = 0, \tag{$\ast$}
\end{equation}
\begin{equation}
z^r \neq 0 \iff x_1, y_1, \ldots, x_r, y_r \text{ 线性无关.} \tag{$\ast\ast$}
\end{equation}
为把 $(*)$ 与 $(**)$ 同性质 $\mathrm{rg}\, z = r$ 联系起来, 从而求解本题, 还缺少关键
的一步.具体来说, 在 $x_1, y_1, \ldots, x_r, y_r$ 线性相关的情形, 我们必须证明: 可以
把 $z$ 分解为 $(r - 1)$ 重乘积和, 即此时 $\mathrm{rg}\, z < r$.我们尝试在特殊情形
$r = 2$ 中把相应的计算做一遍.设 $x_1, y_1, x_2, y_2$ 线性相关, 例如
\[
y_2 = \alpha_1 x_1 + \alpha_2 x_2 + \beta_1 y_1
\]
其中常数 $\alpha_1, \alpha_2, \beta_1 \in K$.则
\[
\begin{array}{r l}
z & = x_1 \wedge y_1 + x_2 \wedge y_2 = x_1 \wedge y_1 + x_2 \wedge (\alpha_1 x_1 + \alpha_2 x_2 + \beta_1 y_1) \\
& = x_1 \wedge y_1 + x_2 \wedge (\alpha_1 x_1 + \beta_1 y_1) = (x_1 + \beta_1 x_2) \wedge y_1 + \alpha_1 x_2 \wedge x_1
\end{array}
\]
但如今在最后一项 $\alpha_1 x_2 \wedge x_1$ 中, 我们可以把因子 $x_1$ 换成 $x_1 +
\beta_1 x_2$, 因为 $x_2 \wedge x_2$ 是平凡的.借助这个小小的技巧得到
\[
\begin{array}{l}
z = (x_1 + \beta_1 x_2) \wedge y_1 + \alpha_1 x_2 \wedge x_1 \\
= (x_1 + \beta_1 x_2) \wedge y_1 + \alpha_1 x_2 \wedge (x_1 + \beta_1 x_2) \\
= (x_1 + \beta_1 x_2) \wedge (y_1 - \alpha_1 x_2)
\end{array}
\]
从而 $\mathrm{rg}\, z < 2$.类似的计算在一般情况下也能进行, 正如我们马上将看到的.

\noindent\textbf{解题路径:} 我们从元素 $z = \sum_{i=1}^{r} x_i \wedge y_i \in
\bigwedge^2 V$(其中向量 $x_i, y_i \in V$) 出发, 并像上面那样证明性质 $(*)$ 与
$(**)$.然后假设 $x_1, y_1, \ldots, x_r, y_r$ 线性相关, 并考虑例如一个如下类型的线性
关系
\[
y_r = \sum_{i=1}^{r} \alpha_i x_i + \sum_{i=1}^{r-1} \beta_i y_i,
\]
其中系数 $\alpha_i, \beta_i \in K$.我们想由此推出 $\mathrm{rg}\, z < r$, 并如下计算:
\[
\begin{array}{l}
z = \sum_{i=1}^{r} x_i \wedge y_i = \sum_{i=1}^{r-1} x_i \wedge y_i + x_r \wedge \sum_{i=1}^{r-1} (\alpha_i x_i + \beta_i y_i) \\
= \sum_{i=1}^{r-1} (x_i + \beta_i x_r) \wedge y_i + x_r \wedge \sum_{i=1}^{r-1} \alpha_i x_i \\
= \sum_{i=1}^{r-1} (x_i + \beta_i x_r) \wedge y_i + x_r \wedge \sum_{i=1}^{r-1} \alpha_i (x_i + \beta_i x_r) \\
= \sum_{i=1}^{r-1} (x_i + \beta_i x_r) \wedge (y_i - \alpha_i x_r)
\end{array}
\]
于是得出 $\mathrm{rg}\, z \leq r - 1$.由这些事实, 我们现在可以容易地读出论断 (i) 与
(ii) 的等价性:
\begin{enumerate}
\item[(i)] $\mathrm{rg}\, z = r$;
\item[(ii)] $z^r \neq 0$, $z^{r+1} = 0$;
\end{enumerate}
对任意的 $z \in \bigwedge^2 V$ 与 $r \in \mathbb{N}$ 成立.先设 (i), 即
$\mathrm{rg}\, z = r$.由 $(*)$ 推出 $z^{r+1} = 0$.此外, 只要我们知道 $z$ 的一个极小
分解所需的 $2r$ 个向量线性无关, 就有 $z^r \neq 0$, 参见 $(**)$.但这些向量不可能线性
相关, 因为否则会有 $\mathrm{rg}\, z < r$, 正如我们刚才证明的.于是推出 (ii).

反过来, 设 (ii) 成立, 即 $z^r \neq 0$ 且 $z^{r+1} = 0$.设 $s = \mathrm{rg}\, z$.由
$z^r \neq 0$ 结合 $(*)$ 推出 $s = \mathrm{rg}\, z \geq r$.在 $s > r$ 的情形, 由于
$z^{r+1} = 0$ 有 $z^s = 0$, 并且我们从 $(**)$ 读出, $z$ 的一个极小分解所需的 $2s$ 个
向量线性相关.这又如我们所证明的, 会意味着 $\mathrm{rg}\, z < s$, 与 $\mathrm{rg}\, z
= s$ 矛盾.于是只剩下情形 $\mathrm{rg}\, z = r$.由此蕴含 (ii) $\Longrightarrow$ (i)
明确, 并且我们总体上得到 (i) 与 (ii) 之间的期望等价性.

% 附录: 习题训练 -- 对应第 5 章(5.1-5.3: 多项式)的解答.
% 源文档 ap5_src.txt; 术语按 CONVENTIONS.md; 引用标签按 05-多项式.tex 中的实际标签.

\section*{对应 5.1-5.3: 多项式}

\subsection*{习题 5.1/1}

\noindent\textbf{起点:} 给定一个环 $R$ 以及一个元素 $t \in R$.考虑 $R$-代数同态
$\Phi \colon R[T] \longrightarrow R$, 它把 $t$ 代入 $T$ 的位置, 即对系数
$a_i \in R$ 通过
\[
\sum_{i \in \mathbb{N}} a_i T^i \longmapsto \sum_{i \in \mathbb{N}} a_i t^i
\]
定义.我们更精确地写 $\Phi_t$ 来代替 $\Phi$.

\noindent\textbf{目标:} 要求确定 $\Phi_t$ 的核, 即证明
$\ker \Phi_t = R[T] \cdot (T - t)$.

\noindent\textbf{策略:} $\Phi_t$ 的核按定义由所有多项式
$\sum_{i \in \mathbb{N}} a_i T^i \in R[T]$ 组成, 它们在 $\Phi_t$ 下的像为零, 即满足
$\sum_{i \in \mathbb{N}} a_i t^i = 0$.对于一般的 $t$, 这个集合相当难以把握, 但对
$t = 0$ 却很容易描述, 因为
\[
\ker \Phi_0 = \left\{ \sum_{i \in \mathbb{N}} a_i T^i;\; a_0 = 0 \right\} = R[T] \cdot T,
\]
正如立即看到的那样.

在习题文本中给出提示: 考虑 $R$-代数同态 $\Psi_t \colon R[T] \longrightarrow R[T]$,
它把 $T$ 换成 $T + t$.如果随后再把 $T$ 换成 $0$, 那么作为复合
$\Phi_0 \circ \Psi_t$ 正好得到把 $T$ 换成 $t$ 的代入同态
$\Phi_t \colon R[T] \longrightarrow R$.因此, 为了确定 $\ker \Phi_t$, 我们可以利用
$\ker \Phi_0 = R[T] \cdot T$, 然后只需证明 $\Psi_t$ 关于此理想的逆像正好是
$R[T] \cdot (T - t)$.

\noindent\textbf{解题路径:} 我们首先证明 $\ker \Phi_0 = R[T] \cdot T$, 即
$\ker \Phi_0$ 与由 $T$ 在 $R[T]$ 中生成的主理想一致.设 $f \in R[T] \cdot T$, 例如
$f = g \cdot T$, 其中 $g \in R[T]$ 是一个多项式, 则由 $\Phi_0(T) = 0$ 得
\[
\Phi_0(f) = \Phi_0(g \cdot T) = \Phi_0(g) \cdot \Phi_0(T) = 0\text{.}
\]
从而 $f \in \ker \Phi_0$.反过来, 由 $f = \sum_{i \in \mathbb{N}} a_i T^i \in \ker
\Phi_0$ 得 $a_0 = \Phi_0(f) = 0$, 于是
\[
f = \sum_{i > 0} a_i T^i = \left( \sum_{i \in \mathbb{N}} a_{i + 1} T^i \right) \cdot T \in R[T] \cdot T,
\]
从而如所期望地推出 $\ker \Phi_0 = R[T] \cdot T$.

接下来我们确认, $R$-代数同态
\[
\Psi_t \colon R[T] \longrightarrow R[T], \qquad T \longmapsto T + t,
\]
是一个同构.因为作为把 $T$ 换成 $T$ 的 $R$-代数同态,
$\Psi_{-t} \circ \Psi_t \colon R[T] \longrightarrow R[T]$ 是恒等映射.对复合
$\Psi_t \circ \Psi_{-t}$ 同样成立, 因此我们把 $\Psi_t$ 识别为同构, 其逆映射由把
$T$ 换成 $T - t$ 的代入同态 $\Psi_{-t}$ 给出.现在 $\Psi_t$ 把理想
$R[T] \cdot (T - t)$ 映到理想 $R[T] \cdot T$ 中, 进而 $\Psi_{-t}$ 又把这个理想映回
理想 $R[T] \cdot (T - t)$.因此同构 $\Psi_t$ 限制为一个双射
\[
R[T] \cdot (T - t) \xrightarrow{\ \sim\ } R[T] \cdot T,
\]
并且有 $\Psi_t^{-1}(R[T] \cdot T) = R[T] \cdot (T - t)$.最后利用关系
$\Phi_t = \Phi_0 \circ \Psi_t$, 得到
\[
\ker \Phi_t = \Psi_t^{-1}(\ker \Phi_0) = \Psi_t^{-1}(R[T] \cdot T) = R[T] \cdot (T - t),
\]
这就是要证明的.

\subsection*{习题 5.1/2}

\noindent\textbf{起点:} 给定一个 $K$-向量空间 $V \neq 0$ 以及一个自同态
$\varphi \colon V \longrightarrow V$, 即 $\varphi \in \operatorname{End}_K(V)$.考虑
$K$-代数同态 $\Phi \colon K[T] \longrightarrow \operatorname{End}_K(V)$, 它把
$\varphi$ 代入 $T$ 的位置, 其中我们更精确地写 $\Phi_\varphi$ 来代替 $\Phi$.

\noindent\textbf{目标:} 要求确定在 $\varphi = \mathrm{id}$ 与 $\varphi = 0$ 两种情形下的核
$\ker \Phi_\varphi$.

\noindent\textbf{策略:} 系数为 $a_i \in K$ 的多项式
$\sum_{i \in \mathbb{N}} a_i T^i \in K[T]$ 属于 $\ker \Phi_\varphi$, 当且仅当
$\sum_{i \in \mathbb{N}} a_i \varphi^i = 0$.对 $\varphi = 0$ 特别得出
$\ker \Phi_0 = \{ \sum_{i \in \mathbb{N}} a_i T^i;\; a_0 \varphi^0 = 0 \}$.但这里
$\varphi^0 = \mathrm{id} \neq 0$, 因为 $V$ 非平凡, 因此
\[
\ker \Phi_0 = \left\{ \sum_{i \in \mathbb{N}} a_i T^i \in K[T];\; a_0 = 0 \right\} = K[T] \cdot T\text{.}
\]
由于 $\mathrm{id} \neq 0$, 相应地得出
\[
\ker \Phi_{\mathrm{id}} = \Bigl\{ \sum_{i \in \mathbb{N}} a_i T^i \in K[T];\; \sum_{i \in \mathbb{N}} a_i = 0 \Bigr\}\text{.}
\]
特别地, $T - 1 \in \ker \Phi_{\mathrm{id}}$, 从而
$K[T] \cdot (T - 1) \subset \ker \Phi_{\mathrm{id}}$.可以猜想, 后一个包含关系实际上
就是等式.为确认这一点, 我们尝试把 $\ker \Phi_{\mathrm{id}}$ 中的元素写成 $T - 1$ 的
倍数.为此考虑一个多项式 $f = \sum_{i \in \mathbb{N}} a_i T^i \in \ker
\Phi_{\mathrm{id}}$, 并建立一个方程
\[
\left( \sum_{i \in \mathbb{N}} b_i T^i \right) \cdot (T - 1) = -b_0 + \sum_{i > 0} (b_{i - 1} - b_i) T^i = \sum_{i \in \mathbb{N}} a_i T^i\text{.}
\]
对系数 $b_i \in K$ 得出下列条件:
\[
\begin{array}{l}
b_0 = -a_0, \\
b_1 = b_0 - a_1 = -a_0 - a_1, \\
b_2 = b_1 - a_2 = -a_0 - a_1 - a_2, \\
\dots
\end{array}
\]
由于由 $f \in \ker \Phi_{\mathrm{id}}$ 有 $\sum_{i \in \mathbb{N}} a_i = 0$, 我们得到
$b_i = 0$(对 $i \geq \operatorname{grad} f$).于是 $g = \sum_{i \in \mathbb{N}} b_i T^i$
确实是 $K[T]$ 中的一个多项式, 满足 $f = g \cdot (T - 1)$, 并如所猜想地得出
$\ker \Phi_{\mathrm{id}} = K[T] \cdot (T - 1)$.

\noindent\textbf{解题路径:} 我们比习题文本中的情形更一般地考虑对 $c \in K$ 的自同态
$c \cdot \mathrm{id} \colon V \longrightarrow V$.由于 $T - c \in \ker
\Phi_{c \cdot \mathrm{id}}$, 得出 $K[T] \cdot (T - c) \subset \ker
\Phi_{c \cdot \mathrm{id}}$, 并且我们断言, 这个包含关系实际上就是等式.为此把
$K[T]$ 作为 $K$-向量空间分解为和
\begin{equation}
K[T] = K + K[T] \cdot (T - c), \tag{$*$}
\end{equation}
其中这个和甚至是直和, 因为在 $K[T] \cdot (T - c)$ 中从次数理由看不可能有非平凡的
常数多项式.和关系本身可由下列方程递归地推出:
\[
T^n = T^{n - 1} \cdot T = T^{n - 1} \cdot (T - c) + cT^{n - 1}, \qquad n > 0\text{.}
\]
现设 $f \in \ker \Phi_{c \cdot \mathrm{id}}$, 且按分解 ($*$) 写 $f = g \cdot (T - c) + a$,
其中 $g \in K[T]$ 是一个多项式, $a \in K$ 是常数.于是得
\[
a = \Phi_{c \cdot \mathrm{id}}(g \cdot (T - c) + a) = \Phi_{c \cdot \mathrm{id}}(f) = 0,
\]
从而 $f \in K[T] \cdot (T - c)$.因此如所断言,
$\ker \Phi_{c \cdot \mathrm{id}}$ 确定为 $K[T] \cdot (T - c)$.习题文本中问的是
$c = 1$ 与 $c = 0$ 两种情形.

\noindent\textbf{补充:} 可以把习题文本看成习题 1 的一个特殊情形.由于 $V \neq 0$, 有
$\operatorname{End}_K(V) \neq 0$, 而把 $\operatorname{End}_K(V)$ 解释为 $K$-代数的
典范环同态 $K \longrightarrow \operatorname{End}_K(V)$ 是单射.因此所要考虑的代入同态
$\Phi_{c \cdot \mathrm{id}} \colon K[T] \longrightarrow \operatorname{End}_K(V)$ 都经由
$K \longrightarrow \operatorname{End}_K(V)$ 分解, 而对 $R = K$ 这正好是习题 1 中所考虑
的情形.

\subsection*{习题 5.1/6}

\noindent\textbf{起点:} 给定一个环 $R$ 以及一个理想 $\mathfrak{a} \subset R$.用
$\operatorname{rad} \mathfrak{a}$ 表示所有满足以下条件的元素 $a \in R$ 组成的集合:
存在 $n \in \mathbb{N}$, 使 $a^n \in \mathfrak{a}$.

\noindent\textbf{目标:} 证明 $\operatorname{rad} \mathfrak{a}$ 是 $R$ 中的理想.

\noindent\textbf{策略:} 平凡地有 $\mathfrak{a} \subset \operatorname{rad} \mathfrak{a}$,
从而特别是 $\operatorname{rad} \mathfrak{a} \neq \varnothing$.再设 $a \in \operatorname{rad}
\mathfrak{a}$, 例如 $a^n \in \mathfrak{a}$.则对 $r \in R$ 有
\[
(ra)^n = r^n \cdot a^n \in \mathfrak{a},
\]
从而 $ra \in \operatorname{rad} \mathfrak{a}$.令 $r = -1$, 由此看出, 若
$a \in \operatorname{rad} \mathfrak{a}$, 则 $-a$ 也属于 $\operatorname{rad} \mathfrak{a}$.

只剩下证明: 由 $a, b \in \operatorname{rad} \mathfrak{a}$ 也推出 $a + b \in
\operatorname{rad} \mathfrak{a}$.为了看出可以怎样着手, 我们范例性地考察
$a^2, b^2 \in \mathfrak{a}$ 的情形.即:
\[
\begin{array}{l}
(a + b)^2 = a^2 + 2ab + b^2, \\
(a + b)^3 = a^3 + 3a^2b + 3ab^2 + b^3\text{.}
\end{array}
\]
由第一个等式还不能推出 $(a + b)^2 \in \mathfrak{a}$, 因为其中有 $2ab$ 这一项.但
第三个等式确实给出 $(a + b)^3 \in \mathfrak{a}$, 因为由 $a^2, b^2 \in \mathfrak{a}$,
右边所有的被加项都属于 $\mathfrak{a}$.一般情形应当可以用二项式公式类似地处理.

\noindent\textbf{解题路径:} 我们为 $\operatorname{rad} \mathfrak{a}$ 验证
定义~\mtref{def:5.1.9} 中的条件.首先, 由 $0 \in \operatorname{rad} \mathfrak{a}$ 有
$\operatorname{rad} \mathfrak{a} \neq \varnothing$.再设 $a, b \in \operatorname{rad}
\mathfrak{a}$, 例如 $a^r, b^s \in \mathfrak{a}$.则二项式公式给出
\[
(a - b)^n = \sum_{i = 0}^{n} (-1)^i a^{n - i} b^i,
\]
显然有 $a - b \in \operatorname{rad} \mathfrak{a}$, 只要对 $i = 0, \ldots, n$ 恒有
$n - i \geq r$ 或 $i \geq s$, 从而有 $a^{n - i} \in \mathfrak{a}$ 或
$b^i \in \mathfrak{a}$.现在令 $t = \max \{ r, s \}$ 并选取 $n \geq 2t$, 则对
$i \leq t$ 有 $n - i \geq r$, 对 $i \geq t$ 有 $i \geq s$, 于是得出 $a - b \in
\operatorname{rad} \mathfrak{a}$.

此外, 在策略性思考中我们早已平凡地看出, 由 $a \in \operatorname{rad} \mathfrak{a}$
与 $r \in R$ 也推出 $ra \in \operatorname{rad} \mathfrak{a}$.

\subsection*{习题 5.2/3}

\noindent\textbf{起点:} 要考察的是有理数域 $\mathbb{Q}$.

\noindent\textbf{目标:} 要求确定 $\mathbb{Q}$ 的所有子环.

\noindent\textbf{策略:} $\mathbb{Q}$ 的每个子环都包含单位元 $1$, 从而也包含一切自然数
以及相应的负数.换言之, $\mathbb{Q}$ 的每个子环都包含整数环 $\mathbb{Z}$.这是
$\mathbb{Q}$ 的最小子环.

现设 $R \subset \mathbb{Q}$ 是一个子环, 它包含一个有理数 $x \in \mathbb{Q} - \mathbb{Z}$.
把 $x$ 写成由整数 $a, b \in \mathbb{Z}$ 组成的分数 $x = \frac{a}{b}$, 并假定它是既约的.
则 $\operatorname{ggT}(a, b) = 1$, 按定理~\mtref{thm:5.2.16} 存在整数 $r, s \in \mathbb{Z}$
满足方程 $ra + sb = 1$.于是得
\[
r \cdot \frac{a}{b} + s = \frac{1}{b},
\]
从而 $\frac{1}{b} \in R$.现在显然 $\mathbb{Z}\bigl[ \frac{1}{b} \bigr]$--定义为代入同态
\[
\mathbb{Z}[T] \longrightarrow \mathbb{Q}, \qquad T \longmapsto \frac{1}{b}
\]
的像--是 $\mathbb{Q}$ 的一个子环.它包含分数 $\frac{a}{b}$, 并且显然是具有此性质的
$\mathbb{Q}$ 的最小子环, 从而 $\mathbb{Z}\bigl[ \frac{1}{b} \bigr] \subset R$.

形如 $\mathbb{Z}\bigl[ \frac{1}{b} \bigr]$($b \in \mathbb{Z} - \{ 0 \}$)的环有一个有趣的
性质.若 $d \in \mathbb{Z}$ 是 $b$ 的因子, 则 $\frac{1}{d}$ 是 $\frac{1}{b}$ 的整数倍,
并且有 $\mathbb{Z}\bigl[ \frac{1}{d} \bigr] \subset \mathbb{Z}\bigl[ \frac{1}{b} \bigr]$.
反之, 若甚至 $b = d^n$ 其中指数 $n > 0$, 则可推出 $\frac{1}{b} = \bigl(
\frac{1}{d} \bigr)^n \in \mathbb{Z}\bigl[ \frac{1}{d} \bigr]$, 从而
$\mathbb{Z}\bigl[ \frac{1}{d} \bigr] = \mathbb{Z}\bigl[ \frac{1}{b} \bigr]$.我们从中引出一个
简单推论.若
\[
b = p_1^{n_1} \cdots p_t^{n_t}
\]
是 $b$ 的素因子分解, 指数 $n_i > 0$, 则
\[
\mathbb{Z}\left[ \frac{1}{b} \right] = \mathbb{Z}\left[ \frac{1}{p_1 \cdots p_t} \right]\text{.}
\]
事实上, $p_1 \cdots p_t \mid b$ 蕴含包含关系"$\supset$", 而 $b \mid (p_1 \cdots p_t)^e$
(对某个合适的幂 $e$)蕴含包含关系"$\subset$".因此我们在考察形如
$\mathbb{Z}\bigl[ \frac{1}{b} \bigr]$ 的环时, 总可以假定 $b$ 是两两不同的素数的乘积.
此外还有一个现象显现出来: 不同素数 $p_1, p_2, \ldots$ 的无限序列给出子环的上升链
\[
\mathbb{Z}\left[ \frac{1}{p_1} \right] \subset \mathbb{Z}\left[ \frac{1}{p_1 p_2} \right] \subset \cdots \subset \mathbb{Q},
\]
而这些环的并构成 $\mathbb{Q}$ 的一个子环, 它显然不是形如
$\mathbb{Z}\bigl[ \frac{1}{b} \bigr]$($b \in \mathbb{Z}$)的环.我们接下来要证明, 除上述
例子外, $\mathbb{Q}$ 不再有其他子环.

\noindent\textbf{解题路径:} 对素数集合 $P \subset \mathbb{N}$ 令
\[
R_P = \left\{ \frac{a}{b} \in \mathbb{Q};\; a, b \in \mathbb{Z}, \ b\ \text{是素数}\ p \in P\ \text{的乘积} \right\}\text{.}
\]
则 $R_P$ 是 $\mathbb{Q}$ 的一个子环, 并且我们断言: 当素数集合 $P$ 与 $P'$ 不同时,
$R_P$ 与 $R_{P'}$ 不同.为论证这一点, 考虑一个素数 $p$ 和一个分数
$\frac{a}{b} \in \mathbb{Q}$($a, b \in \mathbb{Z}$).若存在方程 $\frac{1}{p} =
\frac{a}{b}$, 即 $pa = b$, 则 $p \mid b$, 从而对一切不包含 $p$ 的素数集合 $P$, 有
$\frac{1}{p} \notin R_P$.因此若 $P$ 与 $P'$ 是不同的素数集合, 则 $R_P \neq R_{P'}$.

现在我们还要证明, 除类型 $R_P$ 的环外, $\mathbb{Q}$ 没有其他子环.设 $R$ 是
$\mathbb{Q}$ 的一个子环, 并设 $P$ 是所有满足 $\frac{1}{p} \in R$ 的素数 $p$ 组成的集合.
则 $R_P \subset R$, 并且我们断言这已经是等式.设 $x \in R$, 并把 $x$ 写成既约分数
$x = \frac{a}{b}$, 其中 $a, b \in \mathbb{Z}$.则如我们在策略性思考中看到的, 有
$\frac{1}{b} \in R$, 从而对一切整除 $b$ 的素数 $p$ 有 $\frac{1}{p} \in R$.特别地得出
$x = \frac{a}{b} \in R_P$, 于是看出 $R = R_P$ 成立.因此 $\mathbb{Q}$ 的子环由类型
$R_P$ 的环双射地参数化, 其中 $P \subset \mathbb{N}$ 遍历一切只由素数组成的子集.此时
$P = \varnothing$ 给出 $R_P = \mathbb{Z}$, 而 $P$ 取一切素数的集合时给出
$R_P = \mathbb{Q}$.

\noindent\textbf{补充:} 考虑一个环 $R$ 以及一个乘法系 $S \subset R$.后者意味着 $1 \in S$,
且由 $a, b \in S$ 恒推出 $ab \in S$.于是在一定针对零因子的谨慎措施下, 可以考察由一切
分数 $\frac{a}{s}$($a \in R$, $s \in S$)组成的环 $R_S$; 参见"Bosch, Algebra"
(Springer), 2.7 节.称 $R_S$ 为 $R$ 关于乘法系 $S$ 的局部化.上面考察的,对素数集合
$P \subset \mathbb{N}$ 的类型 $R_P$ 的环, 各自都是 $\mathbb{Z}$ 关于由 $P$ 生成的乘法系
$S$(即由 $P$ 中元素的乘积组成的集合)的局部化.

\subsection*{习题 5.2/6}

\noindent\textbf{起点:} 要考察的是整数环 $\mathbb{Z}$ 上的一元多项式环 $\mathbb{Z}[T]$.

\noindent\textbf{目标:} 要证明 $\mathbb{Z}[T]$ 不是主理想环.

\noindent\textbf{策略:} $\mathbb{Z}$ 中典型的素元是素数, 而域 $K$ 上的多项式环
$K[T]$ 中典型的素元是变量 $T$.现在考虑一个素数 $p$ 以及变量 $T \in \mathbb{Z}[T]$,
并假定 $\mathbb{Z}[T]$ 是主理想环.由于 $p$ 与 $T$ 在 $\mathbb{Z}[T]$ 中显然互素, 于是
按定理~\mtref{thm:5.2.16} 存在多项式 $r, s \in \mathbb{Z}[T]$ 满足方程
$r \cdot p + s \cdot T = 1$.由于 $s \cdot T$ 的常数项为零, 对 $r$ 的常数项 $r_0$ 必须
成立关系 $r_0 \cdot p = 1$.但这是不可能的, 因为 $p \in \mathbb{Z}$ 不是单位.所以
$\mathbb{Z}[T]$ 不可能是主理想环.

\noindent\textbf{解题路径:} 我们现在还要更精确地证明: 由一个素数 $p$ 和变量 $T$ 在
$\mathbb{Z}[T]$ 中生成的理想
\[
\mathfrak{a} = \mathbb{Z}[T] \cdot p + \mathbb{Z}[T] \cdot T = \left\{ \sum_{n \in \mathbb{N}} a_n T^n;\; a_n \in \mathbb{Z}, \ p \mid a_0 \right\}
\]
不是主理想.因为若 $f = \sum_{n \in \mathbb{N}} a_n T^n \in \mathfrak{a}$ 是一个生成元, 则
由于 $p \in \mathfrak{a}$, 从次数理由会推出 $a_0 \mid p$ 且 $a_n = 0$(对 $n > 0$).由于
$\mathfrak{a}$ 不是单位理想, $a_0$ 必须与 $p$ 相伴.另一方面, 在 $\mathbb{Z}[T]$ 中不存在
形如 $T = h \cdot p$ 的方程, 所以 $\mathfrak{a}$ 不可能是主理想.

\noindent\textbf{补充:} 可以借助所谓高斯定理证明, $\mathbb{Z}[T]$ 毕竟还是一个因子分解
唯一的环.这个定理断言, 因子分解唯一的环 $R$ 上的多项式环 $R[T]$ 仍是因子分解唯一的;
参见"Bosch, Algebra"(Springer), 2.7/1.

\subsection*{习题 5.3/3}

\noindent\textbf{起点:} 给定一个首一多项式 $f \in \mathbb{R}[T]$.

\noindent\textbf{目标:} 要证明 $f$ 是素的, 当且仅当 $f = T - \alpha$ 其中 $\alpha \in
\mathbb{R}$, 或 $f = (T - \alpha)(T - \bar{\alpha})$ 其中 $\alpha \in \mathbb{C} -
\mathbb{R}$.这里对复数 $\alpha \in \mathbb{C}$, $\bar{\alpha}$ 表示相应的共轭复数, 它
对于 $\alpha = u + iv$($u, v \in \mathbb{R}$)由 $\bar{\alpha} = u - iv$ 给出.

\noindent\textbf{策略:} 由于 $\mathbb{R}[T]$ 是主理想环, 我们借助定理~\mtref{thm:5.2.12}
看到, 元素 $f \in \mathbb{R}[T]$ 是素的, 当且仅当它不可约.首先, 类型
$f = T - \alpha$($\alpha \in \mathbb{R}$)的多项式在 $\mathbb{R}[T]$ 中从次数理由看
不可约.对第二类所考察的类型,带 $\alpha \in \mathbb{C}$ 的多项式 $f$,
\[
f = (T - \alpha)(T - \bar{\alpha}) = T^2 - (\alpha + \bar{\alpha})T + \alpha \bar{\alpha} \in \mathbb{R}[T],
\]
无论何种情形都是实多项式, 因为 $\alpha + \bar{\alpha} = 2\operatorname{Re}(\alpha) \in
\mathbb{R}$, 且 $\alpha \bar{\alpha} = |\alpha|^2 \in \mathbb{R}$.若现在 $\alpha$ 不实,
则 $f$ 此外在 $\mathbb{R}[T]$ 中不可约.因为否则 $f$ 从次数理由看已经会在
$\mathbb{R}[T]$ 中分解为一次因子, 从而有实零点, 这与复数零点 $\alpha, \bar{\alpha} \in
\mathbb{C} - \mathbb{R}$ 不相容.

现在我们来思考, 不可约的首一多项式 $f \in \mathbb{R}[T]$ 具有所断言的形式.习题文本中
建议考察 $\mathbb{R}$-自同构 $\mathbb{C} \longrightarrow \mathbb{C}$, $z \longmapsto
\bar{z}$, 并将其延拓为 $\mathbb{R}[T]$-自同构
\[
\tau \colon \mathbb{C}[T] \longrightarrow \mathbb{C}[T], \quad \sum_{n \in \mathbb{N}} a_n T^n \longmapsto \sum_{n \in \mathbb{N}} \bar{a}_n T^n\text{.}
\]
它保持实多项式不动.设 $f$ 是一个首一实多项式, 在 $\mathbb{R}[T]$ 中不可约, 从而次数
$\geq 1$.然后考察 $f$ 在 $\mathbb{C}[T]$ 中的素因子分解.由于域 $\mathbb{C}$ 是代数闭的
(大定理~\mtref{bigthm:5.3.3}), $f$ 完全分解为一次因子,
\[
f = (T - \alpha_1) \cdots (T - \alpha_n),
\]
其中零点 $\alpha_1, \ldots, \alpha_n \in \mathbb{C}$.但这时有
\[
(T - \alpha_1) \cdots (T - \alpha_n) = f = \tau(f) = (T - \bar{\alpha}_1) \cdots (T - \bar{\alpha}_n),
\]
而 $\mathbb{C}[T]$ 中素因子分解的唯一性表明, $f$ 的复数零点集合在复共轭下不变, 即零点
要么是实的, 要么以成对的,复共轭的非实零点出现.由此容易推出(正如我们将看到的),
$f$ 具有所断言的形式.

\noindent\textbf{解题路径:} 与上面一样, 我们考虑由复共轭诱导的 $\mathbb{R}[T]$-自同构
$\tau \colon \mathbb{C}[T] \longrightarrow \mathbb{C}[T]$.由于复数 $z \in \mathbb{C}$ 是
实的, 当且仅当 $z = \bar{z}$, 我们看到, 复多项式 $f \in \mathbb{C}[T]$ 是实的, 当且仅当
$\tau(f) = f$.特别地, 由此得出, 类型 $(T - \alpha)(T - \bar{\alpha})$($\alpha \in
\mathbb{C}$)的多项式是实的.此外我们可以像上面那样论证, 多项式 $T - \alpha$($\alpha \in
\mathbb{R}$)以及 $(T - \alpha)(T - \bar{\alpha})$($\alpha \in \mathbb{C} - \mathbb{R}$)在
$\mathbb{R}[T]$ 中不可约.

反过来, 设 $f$ 是 $\mathbb{R}[T]$ 中一个不可约(从而非平凡的)首一多项式, 并设
\[
f = (T - \alpha_1) \cdots (T - \alpha_n)
\]
是 $f$ 在 $\mathbb{C}[T]$ 中,零点为 $\alpha_1, \ldots, \alpha_n \in \mathbb{C}$ 的素因子
分解.方程 $\tau(f) = f$ 借助 $\mathbb{C}[T]$ 中素因子分解的唯一性(参见
定理~\mtref{thm:5.2.15})表明, 因子 $(T - \alpha_1), \ldots, (T - \alpha_n)$ 除顺序外与因子
$(T - \bar{\alpha}_1), \ldots, (T - \bar{\alpha}_n)$ 一致.若特别地 $\alpha_1$ 是实的, 则
$\bar{\alpha}_1 = \alpha_1$, 且因子 $(T - \alpha_2), \ldots, (T - \alpha_n)$ 与因子
$(T - \bar{\alpha}_2), \ldots, (T - \bar{\alpha}_n)$ 一致.特别是 $g = (T - \alpha_2)
\cdots (T - \alpha_n)$ 是实多项式, 从而在 $\mathbb{R}[T]$ 中有分解 $f = (T - \alpha_1)
\cdot g$.但由于 $f$ 在 $\mathbb{R}[T]$ 中首一且不可约, 得 $g = 1$, 于是如所期望地
$f = T - \alpha_1$.

另一方面, 若 $\alpha_1$ 不实, 即 $\bar{\alpha}_1 \neq \alpha_1$, 则 $\bar{\alpha}_1$ 出现在
零点 $\alpha_2, \ldots, \alpha_n$ 中, 例如 $\bar{\alpha}_1 = \alpha_2$, 且因子
$(T - \alpha_3), \ldots, (T - \alpha_n)$ 除顺序外与 $(T - \bar{\alpha}_3), \ldots,
(T - \bar{\alpha}_n)$ 一致.相应地,
\[
(T - \alpha_1)(T - \alpha_2), \qquad g = (T - \alpha_3) \cdots (T - \alpha_n)
\]
是满足 $f = (T - \alpha_1)(T - \alpha_2) \cdot g$ 的实多项式.由于 $f$ 已假定首一且不可约,
得 $g = 1$, 从而 $f = (T - \alpha_1)(T - \alpha_2)$.

% 附录: 习题训练 -- 对应第 6 章(6.1-6.5: 标准形理论)的解答.
% 源文档 ap6_src.txt; 术语按 CONVENTIONS.md; 引用标签按 06-标准形理论.tex 中的实际标签.

\section*{对应 6.1-6.5: 标准形理论}

\subsection*{习题 6.1/3}

\noindent\textbf{起点:} 需要考察的是一个具体给出的上三角矩阵 $A \in \mathbb{R}^{4 \times 4}$.

\noindent\textbf{目标:} 计算 $A$ 的全部特征值以及相应的特征空间.此外还要回答 $A$ 是否可对角化.

\noindent\textbf{策略:} 为使记号简便, 我们设 $V = \mathbb{R}^4$, 并考虑由矩阵 $A$ 给出的线性映射
$f \colon V \longrightarrow V$, $x \longmapsto Ax$.一个元素 $\lambda \in \mathbb{R}$ 是否为 $f$ 即 $A$ 的
特征值, 可以借助线性子空间
\[
V_\lambda = \ker(f - \lambda \operatorname{id}) = \left\{x \in V;\; f(x) = \lambda x\right\} \subset V
\]
来读判.确切地说, $\lambda$ 是 $f$ 即 $A$ 的特征值, 当且仅当 $V_\lambda \neq 0$; 在后一种情形下,
$V_\lambda$ 就是属于特征值 $\lambda$ 的特征空间.为了最终检验
$V_\lambda = \ker(f - \lambda \operatorname{id}) \neq 0$, 我们回顾维数公式
\[
\dim_{\mathbb{R}} V = \dim_{\mathbb{R}} \ker(f - \lambda \operatorname{id}) + \dim_{\mathbb{R}} \operatorname{im}(f - \lambda \operatorname{id})
\]
(定理~\mtref{thm:2.1.10}).它表明 $V_\lambda$ 非平凡当且仅当
\[
\operatorname{rg}(A - \lambda E) = \operatorname{rg}(f - \lambda \operatorname{id}) = \dim_{\mathbb{R}} \operatorname{im}(f - \lambda \operatorname{id}) < \dim_{\mathbb{R}} V = 4,
\]
其中 $E \in \mathbb{R}^{4 \times 4}$ 表示单位矩阵.由此我们问题的解题路径已经勾画出来:
$A - \lambda E$ 型矩阵的秩可以用高斯消元法(大定理~\mtref{bigthm:3.2.4})确定, 而这个方法还可以进一步用来求核, 即求所有满足 $Ax = \lambda x$ 的向量 $x \in V$; 参见 \ref{sec:3.5} 节.$A$ 的可对角化问题然后借助推论~\mtref{cor:6.1.12} 解决.

\noindent\textbf{解题路径:} 如我们已看到的, 元素 $\lambda \in \mathbb{R}$ 是矩阵 $A$ 的特征值, 当且仅当
\[
\operatorname{rg}(A - \lambda E) = \operatorname{rg} \begin{pmatrix}
2 - \lambda & 1 & 0 & 1 \\
0 & 2 - \lambda & 0 & 1 \\
0 & 0 & 2 - \lambda & 1 \\
0 & 0 & 0 & 2 - \lambda
\end{pmatrix} < 4\text{.}
\]
现在对 $\lambda \neq 2$, $A - \lambda E$ 已处于行阶梯形, 由此可以读出 $\operatorname{rg}(A - \lambda E) = 4$; 参见大定理~\mtref{bigthm:3.2.4}.另一方面, 对 $\lambda = 2$ 得
\[
\operatorname{rg}(A - 2E) = \operatorname{rg} \begin{pmatrix}
0 & 1 & 0 & 1 \\
0 & 0 & 0 & 1 \\
0 & 0 & 0 & 1 \\
0 & 0 & 0 & 0
\end{pmatrix} = 2,
\]
因为我们把第二行从第三行减去, 从而得到一个秩为 $2$ 的行阶梯形.于是 $A$ 只有唯一的特征值 $\lambda = 2$, 且相应的特征空间 $V_2$ 按公式(定理~\mtref{thm:2.1.10})具有维数
\[
\dim_{\mathbb{R}} V_2 = \dim_{\mathbb{R}} \mathbb{R}^4 - \operatorname{rg}(A - 2E) = 4 - 2 = 2\text{.}
\]
由于 $V_2$ 是 $A$ 唯一的特征空间并且真包含于 $\mathbb{R}^4$ 之中, 我们借助推论~\mtref{cor:6.1.12} 看出 $A$ 不可对角化.

还剩下特征空间 $V_2$ 要确定.如前所述, 我们得到矩阵
\[
\begin{pmatrix}
0 & 1 & 0 & 1 \\
0 & 0 & 0 & 1 \\
0 & 0 & 0 & 0 \\
0 & 0 & 0 & 0
\end{pmatrix}
\]
作为 $A - 2E$ 的行阶梯形.由此齐次线性方程组
$(A - 2E) \cdot (x_1, x_2, x_3, x_4)^t = 0$ 化简为
\[
x_2 + x_4 = 0, \qquad x_4 = 0,
\]
并且, 无论是用直接的方式还是用 \ref{sec:3.5} 节的方法, 都可以看出向量 $(1, 0, 0, 0)^t$ 与 $(0, 0, 1, 0)^t$ 构成相应解空间,从而构成特征空间 $V_2$ 的一个基.

\subsection*{习题 6.1/7}

\noindent\textbf{起点:} 给定 $K$-向量空间之间的线性映射的一个交换图
\[
\begin{tikzcd}[row sep=huge, column sep=huge]
V \arrow[r, "f"] \arrow[d, "h"'] & V \arrow[d, "h"] \\
W \arrow[r, "g"'] & W
\end{tikzcd}
\]
其中 $V$ 设为有限维.

\noindent\textbf{目标:} 对 $h$ 单射, 需证 $f$ 的每个特征值也是 $g$ 的特征值.反过来, 对 $h$ 满射,
需证 $g$ 的每个特征值也是 $f$ 的特征值.此外还应借助例子说明, 上述假设"单射"与"满射"不是多余的.

\noindent\textbf{策略:} 首先设 $h$ 为单射.若 $x \in V - \{0\}$ 是 $f$ 的属于特征值 $\lambda \in K$ 的特征向量, 则由图的交换性
\[
g(h(x)) = h(f(x)) = h(\lambda x) = \lambda h(x)\text{.}
\]
由于 $h$ 单射, 推出 $h(x) \neq 0$, 因而 $h(x) \in W$ 是 $g$ 的属于特征值 $\lambda$ 的特征向量.

现设 $h$ 满射, 且 $y \in W - \{0\}$ 是属于特征值 $\lambda \in K$ 的特征向量.若选取 $y$ 的一个 $h$-原像 $x \in V$, 则我们只知道
\[
h\bigl(f(x)\bigr) = g(y) = \lambda y = h(\lambda x),
\]
即 $f(x) - \lambda x \in \ker h$, 但不能由此推出 $f(x) = \lambda x$, 因为 $h$ 不一定单射.不过显然图
\[
\begin{tikzcd}[row sep=huge, column sep=huge]
V \arrow[r, "f_\lambda"] \arrow[d, "h"'] & V \arrow[d, "h"] \\
W \arrow[r, "g_\lambda"'] & W
\end{tikzcd}
\]
是交换的, 并且由此可以推出 $f_\lambda(\ker h) \subset \ker h$, 其中 $f_\lambda = f - \lambda \operatorname{id}$.若存在元素 $a \in \ker h$ 满足 $f_\lambda(x) = f_\lambda(a)$(例如当 $f_\lambda$ 在 $\ker h$ 上的限制是满射时), 则推出 $f_\lambda(x - a) = 0$, 且 $x - a \in V$ 是 $f$ 的属于特征值 $\lambda$ 的特征向量; 此处注意, 因 $h(x - a) = h(x) = y \neq 0$ 而有 $x - a \neq 0$.尽管 $f|_{\ker h}$ 在一般情况下不是满射, 我们仍将看到, 如果间接地(即反证地)进行, 这类论证就能导向目标.

\noindent\textbf{解题路径:} 对 $\lambda \in K$ 设 $f_\lambda = f - \lambda \operatorname{id}$ 及 $g_\lambda = g - \lambda \operatorname{id}$, 并考虑交换图
\[
\begin{tikzcd}[row sep=huge, column sep=huge]
V \arrow[r, "f_\lambda"] \arrow[d, "h"'] & V \arrow[d, "h"] \\
W \arrow[r, "g_\lambda"'] & W
\end{tikzcd}
\]
若 $h$ 单射, 则 $\ker f_\lambda = V \cap \ker g_\lambda$.若 $\lambda$ 是 $f$ 的特征值, 则推出 $\ker f_\lambda \neq 0$, 从而 $\ker g_\lambda \neq 0$, 于是把 $\lambda$ 识别为 $g$ 的特征值.

现设 $h$ 满射.我们反证, 假设 $\lambda$ 不是 $f$ 的特征值.则 $f_\lambda$ 单射, 且由于 $\dim_K V < \infty$, 按推论~\mtref{cor:2.1.11} 甚至双射.又因上述图的交换性有 $f_\lambda(\ker h) \subset \ker h$, 并且用同样的维数论证推出 $f_\lambda$ 在 $\ker h$ 上的限制是双射.特别得到 $f_\lambda^{-1}(\ker h) = \ker h$.图的追踪表明 $g_\lambda$ 双射, 从而 $\lambda$ 不可能是 $g$ 的特征值.故若 $\lambda$ 是 $g$ 的特征值, 则必然也是 $f$ 的.

为看出假设"单射"与"满射"不是多余的, 我们考虑两个平凡的例子.首先设 $V$ 是非平凡的 $K$-向量空间, 比如 $V = K$, 并设 $f \colon V \longrightarrow V$ 是恒等映射.再设 $W = 0$, 并考虑零映射 $h \colon V \longrightarrow W$ 以及 $g \colon W \longrightarrow W$.则 $1$ 是 $f$ 的特征值, 但 $g$ 没有特征值.另一方面, 考虑 $V = 0$, 零映射 $f \colon V \longrightarrow V$, 以及非平凡的 $K$-向量空间 $W$, 比如 $W = K$, 和恒等映射 $g \colon W \longrightarrow W$.再设 $h \colon V \longrightarrow W$ 是典范包含映射.则 $1$ 是 $g$ 的特征值, 但 $f$ 没有任何特征值.

\noindent\textbf{补充:} 背景是所谓的蛇引理, 我们想对 $h$ 满射的情形解释其应用; 参见 Bosch, Algebraic Geometry and Commutative Algebra (Springer), 1.5/2.确切地说, 对如上的 $f_\lambda$, $g_\lambda$ 以及 $f'_\lambda = f_\lambda|_{\ker h}$, 考虑如下行精确的交换图:
\[
\begin{tikzcd}[row sep=huge, column sep=huge]
0 \arrow[r] & \ker h \arrow[r] \arrow[d, "f'_\lambda"] & V \arrow[r, "h"] \arrow[d, "f_\lambda"] & W \arrow[r] \arrow[d, "g_\lambda"] & 0 \\
0 \arrow[r] & \ker h \arrow[r] & V \arrow[r, "h"] & W \arrow[r] & 0
\end{tikzcd}
\]
蛇引理于是给出典范的精确序列
\[
0 \longrightarrow \ker f'_\lambda \longrightarrow \ker f_\lambda \longrightarrow \ker g_\lambda \longrightarrow \operatorname{coker} f'_\lambda \longrightarrow \operatorname{coker} f_\lambda \longrightarrow \operatorname{coker} g_\lambda \longrightarrow 0,
\]
其中对 $K$-向量空间的线性映射 $\varphi \colon U' \longrightarrow U$, 其余核由
$\operatorname{coker} \varphi = U/\varphi(U')$ 给出.由此我们可以轻易读出我们问题的解.若 $f'_\lambda$ 满射, 则 $\operatorname{coker} f'_\lambda = 0$, 由此推出 $\ker f_\lambda \longrightarrow \ker g_\lambda$ 满射.我们曾通过假设 $f_\lambda$ 单射且 $V$ 有限维, 即假设 $\lambda$ 不是 $f$ 的特征值, 而建立这种情形.由 $\ker f_\lambda = 0$ 于是得出 $\ker g_\lambda = 0$.在一般情况下, 映射 $\ker f_\lambda \longrightarrow \ker g_\lambda$ 满射, 当且仅当映射 $\operatorname{coker} f'_\lambda \longrightarrow \operatorname{coker} f_\lambda$ 单射.

\subsection*{习题 6.2/2}

\noindent\textbf{起点:} 给定有限维 $K$-向量空间 $V$ 及其自同态 $f \colon V \longrightarrow V$.需要考察下列情形:
\begin{enumerate}
\item[(i)] $V = 0$;
\item[(ii)] $f = \operatorname{id}$;
\item[(iii)] $f = 0$;
\item[(iv)] $V = V_1 \oplus V_2$, 其中 $V_1, V_2 \subset V$ 是线性子空间, 且 $f|_{V_1} = \operatorname{id}$ 以及 $f|_{V_2} = 0$.
\end{enumerate}

\noindent\textbf{目标:} 在每种情形下确定 $f$ 的最小多项式 $p_f$.

\noindent\textbf{策略:} 所有思考的出发点是最小多项式 $p_f$ 的定义, 按定理~\mtref{thm:6.2.9} 与定义~\mtref{def:6.2.10}, 它作为代入同态 $\varphi_f \colon K[T] \longrightarrow \operatorname{End}_K(V)$, $T \longmapsto f$ 的核的生成元素.这里特征多项式 $\chi_f$ 连同凯莱--哈密顿定理(定理~\mtref{thm:6.2.11})也可能有帮助, 因为可以利用 $p_f$ 是 $\chi_f$ 的因子这一事实.其中 $V = 0$ 构成一个特殊情形, 因为这里只有一个自同态 $V \longrightarrow V$, 我们可以把它理解为恒等映射 $\operatorname{id}$, 也可以理解为零映射 $0$, 从而特别有 $\operatorname{id} = 0$.

\noindent\textbf{解题路径:} 我们考虑上述代入同态 $\varphi_f \colon K[T] \longrightarrow \operatorname{End}_K(V)$, $T \longmapsto f$.在情形 (i), 即 $V = 0$ 时, 有 $\operatorname{End}_K(V) = 0$, 从而 $\varphi_f = 0$.因此 $\ker \varphi_f$ 是单位理想, 于是 $p_f = 1$.我们在以下排除这一特殊情形, 从现在起设 $V$ 是维数 $n > 0$ 的 $K$-向量空间.

于是情形 (ii) 中的 $f = \operatorname{id}$ 关于 $V$ 的任意基由单位矩阵 $E \in K^{n \times n}$ 描述, 且 $\chi_f = (T - 1)^n$.再利用凯莱--哈密顿定理~\mtref{thm:6.2.11} 中 $p_f$ 是 $\chi_f$ 的因子这一事实, 并检验 $(T - 1)^0(f) = \operatorname{id} \neq 0$ 以及 $(T - 1)(f) = 0$, 得 $p_f = T - 1$.

现设情形 (iii) 中的 $f = 0$.与 (ii) 类似地推出 $\chi_f = T^n$ 以及 $p_f = T$.

最后在情形 (iv) 中我们可以设 $V_1 \neq 0$ 且 $V_2 \neq 0$, 因为否则我们又处于 (ii) 或 (iii) 的情形.设 $\dim V_1 = n_1$, $\dim V_2 = n_2$, 其中 $n_1, n_2 > 0$, 并且自然有 $n_1 + n_2 = n$.由 $f|_{V_1} = \operatorname{id}$ 以及 $f|_{V_2} = 0$, 检验得到 $f$ 被 $(T - 1) \cdot T$ 零化, 但不受它的任何真因子零化.因此唯一可能 $p_f = (T - 1) \cdot T$.

\subsection*{习题 6.2/5}

\noindent\textbf{起点:} 给定有限维 $K$-向量空间 $V$ 及其自同态 $f \colon V \longrightarrow V$.$f$ 设为自同构, 从而逆映射 $f^{-1} \colon V \longrightarrow V$ 也作为 $V$ 的自同构存在.

\noindent\textbf{目标:} 证明存在多项式 $q \in K[T]$ 使得 $f^{-1} = q(f)$.

\noindent\textbf{策略:} 如我们所知, 存在形如
\[
f^n + c_1 f^{n - 1} + \ldots + c_n f^0 = 0
\]
的方程, 其中常数 $c_i \in K$.用一个小技巧由此得
\[
f \circ (f^{n - 1} + c_1 f^{n - 2} + \ldots + c_{n - 1} f^0) = -c_n f^0 = -c_n \operatorname{id},
\]
并且, 只要 $c_n \neq 0$, 我们就可以在左边乘以 $f^{-1}$ 而推得
\[
f^{-1} = -c_n^{-1}(f^{n - 1} + c_1 f^{n - 2} + \ldots + c_{n - 1} f^0)\text{.}
\]
现在自然要猜测, 至少对最小多项式 $p_f \in K[T]$, 常数项不消失.为确认这一点, 我们反证, 假设 $p_f = T \cdot p$, 其中 $p \in K[T]$ 是某个多项式.则因 $\operatorname{grad} p < \operatorname{grad} p_f$ 有 $p(f) \neq 0$, 且存在元素 $x \in V$ 使 $y = p(f)(x) \neq 0$.但 $f(y) = p_f(f)(x) = 0$ 将蕴含 $y \in \ker f$, 从而 $\ker f \neq 0$, 而这是不可能的, 因为 $f$ 是自同构.因此形如 $p_f = T \cdot p$ 的分解不可能, 于是我们看到最小多项式 $p_f$ 的常数项非平凡.

\noindent\textbf{解题路径:} 最简单的是考察特征多项式 $\chi_f$.由凯莱--哈密顿定理, 对它成立 $\chi_f(f) = 0$.此外按定理~\mtref{thm:6.2.3}, $\chi_f$ 的常数项与 $\det f$ 相差一个符号, 而按定理~\mtref{thm:4.3.4} 有 $\det f \neq 0$.因此上面演示的技巧可以用来把 $f^{-1}$ 表示为系数取自 $K$ 的 $f$ 的多项式.

\subsection*{习题 6.3/3}

\noindent\textbf{起点:} 设 $R$ 是主理想环.给定元素 $a_{11}, \ldots, a_{1n} \in R$, 它们互素, 即满足 $\operatorname{ggT}(a_{11}, \ldots, a_{1n}) = 1$.

\noindent\textbf{目标:} 求其他元素 $a_{ij} \in R$, $i = 2, \ldots, n$, $j = 1, \ldots, n$, 使得矩阵 $A = (a_{ij})_{i, j = 1, \ldots, n}$ 在 $R^{n \times n}$ 中可逆.也就是说, $A$ 的第一行由 $(a_{11}, \ldots, a_{1n})$ 固定, 要证明可以添加 $n - 1$ 行, 使得得到一个在 $R^{n \times n}$ 中可逆的矩阵.

\noindent\textbf{策略:} 我们想把元素 $a_{11}, \ldots, a_{1n}$ 的互素性与初等因子理论联系起来, 因此考察 $A_1 = (a_{11}, \ldots, a_{1n})$ 的初等因子, 其中 $A_1$ 看作 $R^{1 \times n}$ 中的矩阵.引理~\mtref{lem:6.3.5} 中适用于主理想环的方法表明, $A_1$ 可以通过乘以 $R^{n \times n}$ 中可逆的矩阵(我们已把它们具体解释为列变换)从右边化为 $( \alpha, 0, \ldots, 0 )$ 的形式, 其中常数 $\alpha \in R$ 是 $A_1$ 的第一个(也是唯一的)初等因子.因此存在形如 $A_1 \cdot T = ( \alpha, 0, \ldots, 0 )$ 的方程, 其中 $T \in R^{n \times n}$ 可逆.于是 $A_1 = ( \alpha, 0, \ldots, 0 ) \cdot T^{-1}$, 我们看到 $\alpha$ 是 $a_{11}, \ldots, a_{1n}$ 的公因子, 从而必是单位.因此 $A_1$ 可以通过合适的可逆列变换化为行向量 $(1, 0, \ldots, 0) \in R^{1 \times n}$.若把这个行作为单位矩阵 $E \in R^{n \times n}$ 的第一行, 并把先前考察的列变换逆序作用于 $E$, 就应得到一个满足要求的矩阵 $A$.

\noindent\textbf{解题路径:} 设 $F = R^n$ 是 $R$ 上全部 $n$ 元组构成的自由 $R$-模, 我们按惯例把它们解释为列向量.再设 $M$ 是由列 $(a_{11}, \ldots, a_{1n})^t$ 生成的 $R$-子模.则初等因子定理~\mtref{bigthm:6.3.4} 在利用推论~\mtref{cor:6.3.7} 的情况下断言: 存在 $F$ 的基 $X = (x_1, \ldots, x_n)$, 使得对合适的常数 $\alpha \in R$ 有 $M = R \cdot \alpha x_1$.由于 $\alpha$ 必须整除所有元素 $a_{1j}$, 而这些元素互素, 故可以设 $\alpha = 1$, 从而甚至设 $x_1 = (a_{11}, \ldots, a_{1n})^t$.对 $R^n$ 的典范基 $e$, 现考虑基变换矩阵 $A_{\mathrm{id}, X, e}$.它可逆, 且以 $(a_{11}, \ldots, a_{1n})^t$ 为第一列.于是 $A_{\mathrm{id}, X, e}^t$ 是 $R^{n \times n}$ 中可逆矩阵, 且如所愿以 $(a_{11}, \ldots, a_{1n})$ 为第一行.

\subsection*{习题 6.3/6}

\noindent\textbf{起点:} 给定含 $1$ 的交换环 $R$ 上有限生成的模 $M$, 且 $M$ 拥有一个\emph{自由生成系}(\textup{freies Erzeugendensystem})\index{ziyoushengchengxi@自由生成系(freies Erzeugendensystem)}.

\noindent\textbf{目标:} 证明 $M$ 的每个自由生成系都是有限的, 并且由同样多的元素组成.

\noindent\textbf{策略:} 首先要指出 $M = 0$ 的情形是平凡的.因此可以设 $M \neq 0$, 从而特别 $R \neq 0$.若 $X = (x_i)_{i \in I}$ 是自由生成系, $Y$ 是 $M$ 的有限生成系, 则每个元素 $y \in Y$ 是有限多个生成元 $x_i$ 的线性组合, 于是总体上看到 $M$ 被 $X$ 的某个有限子系统生成.但由于 $X$ 是自由生成系, 必然推出 $X$ 本身有限.

再设 $X = (x_1, \ldots, x_m)$ 与 $Y = (y_1, \ldots, y_n)$ 是 $M$ 的两个(有限)自由生成系.给元素 $z \in M$ 配上其关于 $M$ 的基 $X$ 或 $Y$ 的相应坐标列向量 $z_X$ 或 $z_Y$, 得到 $R$-模同构 $M \longrightarrow R^m$ 以及 $M \longrightarrow R^n$.总体上由此得到 $R$-模同构 $\varphi \colon R^m \longrightarrow R^n$, 要证明这样的同构只在 $m = n$ 时存在.

例如, 若 $R$ 是域, 则由维数理论推出 $m = n$; 参见定理~\mtref{thm:2.1.8}.若 $R$ 不是域, 但仍为主理想环, 则为了证明推论~\mtref{cor:6.3.7} 可以选取素元素 $p \in R$, 并考察由 $\varphi$ 诱导的同构 $\varphi_p \colon (R/pR)^m \xrightarrow{\ \sim\ } (R/pR)^n$.由于按推论~\mtref{cor:5.2.17}, $R/pR$ 是域, 这里也可以推出 $m = n$.同样的论证对一般含 $1$ 的交换环 $R$ 也适用, 只要在 $R$ 中能找到理想 $\mathfrak{m}$, 使得剩余类环 $R/\mathfrak{m}$ 是域.

对整环 $R$ 还有一种可能的推理方式.即可以对 $R$ 构造域
\[
Q(R) = \left\{\frac{a}{b};\; a, b \in R,\ b \neq 0\right\},
\]
即按分数运算的通常规则与约定构造全部商.我们已在 \ref{sec:6.2} 节对域 $K$ 上的多项式环 $R = K[T]$ 用过这一构造.然后由同构 $\varphi \colon R^m \longrightarrow R^n$ 通过作商得到一个 $Q(R)$-向量空间的同构 $\varphi_{Q(R)} \colon Q(R)^m \longrightarrow Q(R)^n$, 同样可以推出 $m = n$.原则上这个方法在一般情形也可应用, 只要证明每个环 $R \neq 0$ 都含有所谓的\emph{素理想}(\textup{Primideal})\index{sulixiang@素理想(Primideal)} $\mathfrak{p}$, 它由剩余类环 $R/\mathfrak{p}$ 是整环来刻画.则同构 $\varphi \colon R^m \longrightarrow R^n$ 诱导 $(R/\mathfrak{p})$-模的同构 $\varphi_{\mathfrak{p}} \colon (R/\mathfrak{p})^m \longrightarrow (R/\mathfrak{p})^n$, 而现在 $R/\mathfrak{p}$ 是整环.

\noindent\textbf{解题路径:} $M = 0$ 的情形是平凡的, 因此设 $M \neq 0$, 特别 $R \neq 0$.如已说明的, 只需构造理想 $\mathfrak{m} \subset R$, 使得 $R/\mathfrak{m}$ 是域; 这样的理想尤其是素理想, 在 $R$ 中称为\emph{极大理想}(\textup{maximales Ideal})\index{jidalixiang@极大理想(maximales Ideal)}.

为了在环 $R \neq 0$ 中构造极大理想, 考察由包含关系给出偏序的集合 $\mathcal{A}$, 它由一切真理想 $\mathfrak{a} \subsetneq R$ 组成.按佐恩引理(引理~\mtref{lem:1.5.15}), $\mathcal{A}$ 有极大元素 $\mathfrak{m}$, 即对每个理想 $\mathfrak{a} \subset R$, 只要 $\mathfrak{m} \subset \mathfrak{a}$ 就已经推出 $\mathfrak{m} = \mathfrak{a}$ 或 $\mathfrak{a} = R$.我们断言剩余类环 $R/\mathfrak{m}$ 是域, 并为此证明 $R/\mathfrak{m}$ 中每个非零元素 $\bar{\alpha}$ 都是单位.事实上, 主理想 $(\bar{\alpha}) \subset R/\mathfrak{m}$ 在典范投影 $R \longrightarrow R/\mathfrak{m}$ 下的原像给出一个理想 $\mathfrak{a} \subset R$, 它满足 $\mathfrak{m} \subsetneq \mathfrak{a}$, 从而 $\mathfrak{a} = R$.于是 $(\bar{\alpha}) = R/\mathfrak{m}$, 我们看到 $\bar{\alpha}$ 是 $R/\mathfrak{m}$ 中的单位.因此 $\mathcal{A}$ 的极大元素 $\mathfrak{m}$ 导致一个剩余类环 $R/\mathfrak{m}$, 它是域.更一般地, 对理想 $\mathfrak{m} \subset R$ 容易证明: $R/\mathfrak{m}$ 是域, 当且仅当 $\mathfrak{m}$ 是 $\mathcal{A}$ 中的极大元素; 因此"极大理想"这一名称对这类理想是合理的.

在策略部分我们已经证明 $M$ 的自由生成系是有限的, 并且任意两个这样的系, 比如长度为 $m$ 与 $n$, 导致 $R$-模同构 $\varphi \colon R^m \longrightarrow R^n$.若 $\mathfrak{m} \subset R$ 是极大理想, 则 $\varphi$ 诱导 $(R/\mathfrak{m})$-向量空间的同构 $\varphi_{\mathfrak{m}} \colon (R/\mathfrak{m})^m \xrightarrow{\ \sim\ } (R/\mathfrak{m})^n$, 从而如愿借助维数论证(定理~\mtref{thm:2.1.8})推出 $m = n$.

\subsection*{习题 6.4/1}

\noindent\textbf{起点:} 给定有限阿贝尔群 $G$; $G$ 的元素个数记为 $\operatorname{ord} G$, 称为 $G$ 的阶.

\noindent\textbf{目标:} 若 $H \subset G$ 是子群, 则 $\operatorname{ord} H$ 是 $\operatorname{ord} G$ 的因子.另一方面, 对 $\operatorname{ord} G$ 的每个因子 $d$, 存在阶为 $d$ 的子群 $H \subset G$.

\noindent\textbf{策略:} 自然地想用结构定理(推论~\mtref{cor:6.4.4}).若 $P \subset \mathbb{N}$ 表示素数集, 则这个结果断言 $G$ 是 $\mathbb{Z}/p^n \mathbb{Z}$ 型群的直和, 其中 $p \in P$, $n \in \mathbb{N}$ 变化.由这样的分解可以确定 $G$ 的阶, 即各被加项之阶的乘积.因为若 $A$, $B$ 是有限阿贝尔群, 则
\begin{equation}
\operatorname{ord}(A \oplus B) = \operatorname{ord}(A) \cdot \operatorname{ord}(B), \tag{$*$}
\end{equation}
因为有限直和也可以解释为笛卡尔积.

群 $G_{p, n} = \mathbb{Z}/p^n \mathbb{Z}$ 的阶是 $p^n$, 这很容易验证.显然 $G_{p, n}$ 含有由 $n + 1$ 个子群构成的链
\[
\{0\} = p^n G_{p, n} \subset p^{n - 1} G_{p, n} \subset \ldots \subset p^0 G_{p, n} = G_{p, n},
\]
其中
$p^i G_{p, n} = p^i \mathbb{Z}/p^n \mathbb{Z} \simeq \mathbb{Z}/p^{n - i} \mathbb{Z}$,
从而对 $i = 0, \ldots, n$ 有 $\operatorname{ord} p^i G_{p, n} = p^{n - i}$.于是对 $\operatorname{ord} G_{p, n} = p^n$ 的每个因子 $d$ 都存在阶为 $d$ 的子群 $H \subset G_{p, n}$.借助推论~\mtref{cor:6.4.4} 并利用 $(*)$, 还可以对任意有限阿贝尔群 $G$ 看出: 对 $\operatorname{ord} G$ 的每个因子 $d$ 都存在阶为 $d$ 的子群 $H \subset G$.

还剩下要考虑: 对子群 $H \subset G$ 恒有 $\operatorname{ord} H$ 是 $\operatorname{ord} G$ 的因子.我们考察简单情形 $G = \mathbb{Z}/p^n \mathbb{Z}$, $p \in P$, $n \in \mathbb{N}$.则可以把 $\mathbb{Z}/p^n \mathbb{Z}$ 视为 $\mathbb{Z}$ 的剩余类环, 并考虑典范投影 $\pi \colon \mathbb{Z} \longrightarrow \mathbb{Z}/p^n \mathbb{Z}$.显然 $\mathbb{Z}/p^n \mathbb{Z}$ 的子群 $H$ 已经是 $\mathbb{Z}/p^n \mathbb{Z}$ 中的理想, 从而 $\mathfrak{a} = \pi^{-1}(H)$ 是 $\mathbb{Z}$ 中的理想, 它含有理想 $\ker \pi = (p^n)$.由于 $\mathbb{Z}$ 是主理想环, $\mathfrak{a}$ 是主理想, 比如 $\mathfrak{a} = (a)$.又由 $p^n \in \mathfrak{a}$ 得 $a \mid p^n$, 比如 $a = p^i$, 其中 $0 \leq i \leq n$, 由此可以推出 $H = p^i \mathbb{Z}/p^n \mathbb{Z} \simeq \mathbb{Z}/p^{n - i} \mathbb{Z}$.特别得 $\operatorname{ord} H = p^{n - i}$, 而它是 $\operatorname{ord} G = p^n$ 的因子.

若 $G$ 是任意有限阿贝尔群, 则按推论~\mtref{cor:6.4.4} 存在 $G$ 到 $\mathbb{Z}/p^n \mathbb{Z}$ 型子群的直和分解.但对子群 $H \subset G$, 一般不能指望这个分解限制成 $H$ 的相应分解.不过我们可以更小心地处理.例如, 若 $G = A \oplus B$, 其中 $A$, $B$ 是有限阿贝尔群, 则存在典范的短正合列
\[
0 \longrightarrow A \xrightarrow{\iota} A \oplus B \xrightarrow{\tau} B \longrightarrow 0,
\]
其中 $\iota \colon A \longrightarrow A \oplus B$ 是第一个被加项的包含, $\tau \colon A \oplus B \longrightarrow B$ 是到第二个被加项的投影.若 $H \subset A \oplus B$ 是子群, 则得到 $A$, $A \oplus B$ 与 $B$ 中子群的短正合列
\[
0 \longrightarrow \iota^{-1}(H) \longrightarrow H \longrightarrow \tau(H) \longrightarrow 0,
\]
从中读出 $\operatorname{ord} H = \operatorname{ord}(\iota^{-1}(H)) \cdot \operatorname{ord}(\tau(H))$.因此若已知 $A$ 与 $B$ 是这样的阿贝尔群, 使子\emph{群的阶}(\textup{Ordnung einer Gruppe})\index{qundejie@群的阶(Ordnung einer Gruppe)}恒为 $A$ 或 $B$ 之阶的因子, 则可以推出 $\operatorname{ord} H \mid \operatorname{ord} A \cdot \operatorname{ord} B$, 即 $\operatorname{ord} H \mid \operatorname{ord} G$.这个论证可以递归地应用, 并借助推论~\mtref{cor:6.4.4} 最终对任意有限阿贝尔群 $G$ 及其子群 $H \subset G$ 得到 $\operatorname{ord} H \mid \operatorname{ord} G$.

\noindent\textbf{解题路径:} 由于 $G$ 是有限群, 从而是挠群, 推论~\mtref{cor:6.4.4} 中的分解化简为
\[
G = \bigoplus_{p \in P} \bigoplus_{j_p = 1 \ldots s_p} \mathbb{Z}/p^{n(p, j_p)} \mathbb{Z}
\]
其中指数 $1 \leq n(p, 1) \leq \ldots \leq n(p, s_p)$.由于对 $p \in P$, $n \in \mathbb{N}$ 有 $\operatorname{ord} \mathbb{Z}/p^n \mathbb{Z} = p^n$, $G$ 的阶按上述规则 $(*)$ 计算得
\[
\operatorname{ord} G = \prod_{p \in P} p^{\sum_{j_p = 1 \ldots s_p} n(p, j_p)}\text{.}
\]
若 $d \in \mathbb{N}$ 是 $\operatorname{ord} G$ 的因子, 比如 $d = \prod_{p \in P} p^{r_p}$, 其中指数 $r_p \leq \sum_{j_p = 1 \ldots s_p} n(p, j_p)$, $p \in P$, 则选取自然数 $r(p, j_p) \leq n(p, j_p)$, $p \in P$, $j_p = 1 \ldots s_p$, 满足 $\sum_{j_p = 1 \ldots s_p} r(p, j_p) = r_p$.再分别考虑 $\mathbb{Z}/p^{n(p, j_p)} \mathbb{Z}$ 的一个阶为 $p^{r(p, j_p)}$ 的子群 $H_{p, j_p}$; 如我们在策略部分所见, 这是可能的.则
\[
H = \bigoplus_{p \in P} \bigoplus_{j_p = 1 \ldots s_p} H_{p, j_p} \subset \bigoplus_{p \in P} \bigoplus_{j_p = 1 \ldots s_p} \mathbb{Z}/p^{n(p, j_p)} \mathbb{Z}
\]
是 $G$ 中阶为 $d$ 的子群.故对 $\operatorname{ord} G$ 的每个因子 $d$ 都存在 $G$ 的阶为 $d$ 的子群 $H \subset G$.

另一方面设给定子群 $H \subset G$.为看出 $\operatorname{ord} H$ 是 $\operatorname{ord} G$ 的因子, 把 $G$ 看作 $\mathbb{Z}$-模, 并用结构定理(推论~\mtref{cor:6.4.4}).作为有限阿贝尔群, $G$ 是 $\mathbb{Z}$-\emph{挠模}(\textup{Torsionsmodul})\index{naomo@挠模(Torsionsmodul)}, 且 $G = \bigoplus_{p \in P} G_p$; 其中对 $p \in P$, $G_p \subset G$ 是 $p$-挠子模, 即
\[
G_p = \{x \in G;\; p^n x = 0 \text{ 对合适的 } n \in \mathbb{N}\}\text{.}
\]
相应地 $H = \bigoplus_{p \in P} H_p$, 其中对 $p$-挠子模自然必须有 $H_p \subset G_p$.推论~\mtref{cor:6.4.4} 中给出的分解 $G_p = \bigoplus_{j_p = 1 \ldots s_p} G_{p, j_p}$ 连同开头提到的规则 $(*)$ 表明 $G_p$ 以及相应 $H_p$ 都是 $p$-幂阶的阿贝尔群.那么显然由 $H_p \subset G_p$ 已经推出对 $p \in P$ 有 $\operatorname{ord} H_p \mid \operatorname{ord} G_p$.利用 $(*)$ 以及分解 $G = \bigoplus_{p \in P} G_p$, $H = \bigoplus_{p \in P} H_p$, 由此得到 $\operatorname{ord} H \mid \operatorname{ord} G$.

\noindent\textbf{补充:} 对有限(不一定是阿贝尔)群 $G$ 及其子群 $H \subset G$, 恒有 $\operatorname{ord} H \mid \operatorname{ord} G$; 这由所谓的拉格朗日(Lagrange)定理推出.为证明这个定理, 对元素 $a \in G$ 考虑形如 $aH = \{ah;\; h \in H\} \subset G$ 的子集, 即关于 $H$ 的\emph{左陪集}(\textup{Linksnebenklasse})\index{zuopeiji@左陪集(Linksnebenklasse)}.在简单计算中证明: $H$ 的任意两个左陪集含有同样多的元素, 并且只要不同就互不相交.于是 $G$ 的阶等于 $H$ 的阶与 $H$ 在 $G$ 中左陪集个数之积; 这就是\emph{拉格朗日定理}(\textup{Lagrange, Satz von})\index{lagelangridingli@拉格朗日定理(Lagrange, Satz von)}的陈述.特别有 $\operatorname{ord} H \mid \operatorname{ord} G$.

\subsection*{习题 6.4/6}

\noindent\textbf{起点:} 考察有限群 $G$, 它是域 $K$ 的单位群 $K^*$ 的乘法子群.

\noindent\textbf{目标:} 证明 $G$ 是循环的, 即存在 $n \in \mathbb{Z}$ 使 $G \simeq \mathbb{Z}/n \mathbb{Z}$.

\noindent\textbf{策略:} 我们选取元素 $g \in G$, 考察它的幂 $g^0, g^1, g^2, \ldots$.由于 $G$ 有限, 这些幂不能两两不同.因此存在指数 $r, s \in \mathbb{N}$, 比如 $r < s$, 使 $g^r = g^s$.对 $n = s - r$ 由此得 $g^n = 1$.故对每个元素 $g \in G$ 存在指数 $n > 0$ 使 $g^n = 1$.或者我们也可以考察把 $z \in \mathbb{Z}$ 映到幂 $g^z$ 的群同态 $\varphi \colon \mathbb{Z} \longrightarrow G$.由于 $G$ 只含有限多个元素, 推出 $\ker \varphi \neq 0$.$\ker \varphi$ 是 $\mathbb{Z}$ 的子群, 从而是 $\mathbb{Z}$-子模, 因而是 $\mathbb{Z}$ 中的理想, 尤其是主理想.故存在 $n > 0$ 使 $\ker \varphi = n \mathbb{Z}$, 其中 $n$ 是满足 $g^n = 1$ 的最小自然数 $> 0$.按同态定理, $\varphi$ 诱导单射 $\bar{\varphi} \colon \mathbb{Z}/n \mathbb{Z} \hookrightarrow G$, 我们看到像 $\langle g \rangle = \operatorname{im} \bar{\varphi}$ 恰由 $n$ 个两两不同的元素 $g^0, g^1, \ldots, g^{n - 1}$ 组成.由于 $g^n = 1$, $g$ 是多项式 $T^n - 1 \in K[T]$ 的零点, 而每个幂 $g^z$(其中 $z \in \mathbb{Z}$) 也有这个性质.因此 $\langle g \rangle$ 中的 $n$ 个元素都是 $T^n - 1$ 的零点.但 $n$ 次多项式在 $K$ 中至多有 $n$ 个零点, 所以除 $\langle g \rangle$ 的元素之外 $T^n - 1$ 在 $K$ 即 $G$ 中没有其他零点.我们从中得出一个重要推论: $G$ 含有一个 $\mathbb{Z}/n \mathbb{Z}$ 型的子群, 但不能含有 $\mathbb{Z}/n \mathbb{Z} \oplus \mathbb{Z}/n \mathbb{Z}$ 型,甚至 $\mathbb{Z}/n \mathbb{Z} \oplus \mathbb{Z}/d \mathbb{Z}$ 型($d > 1$ 是 $n$ 的因子)的子群, 因为后者的所有元素都是多项式 $T^n - 1$ 的零点.这尤其对结构定理(推论~\mtref{cor:6.4.4})框架下 $G$ 的分解有影响.

\noindent\textbf{解题路径:} 作为有限群, $G$ 是挠群.因此定理~\mtref{thm:6.4.2} 给出分解
\[
G \simeq \bigoplus_{j = 1}^{s} \mathbb{Z}/n_j \mathbb{Z}
\]
其中自然数 $n_1, \ldots, n_s > 1$ 满足对 $1 \leq j < s$ 有 $n_j \mid n_{j + 1}$.由于有限直和也可以解释为笛卡尔积, $G$ 恰有 $\prod_{j = 1}^{s} n_j$ 个元素, 且每个元素显然都是多项式 $T^n - 1 \in K[T]$ 的零点, 其中我们设 $n = n_s$.另一方面, 这个多项式在 $K$ 中至多有 $n$ 个零点, 从而推出 $s = 1$, 于是 $G \simeq \mathbb{Z}/n \mathbb{Z}$, 如所愿.

\subsection*{习题 6.5/2}

\noindent\textbf{起点:} 给定两个矩阵 $A$, $B \in \mathbb{R}^{3 \times 3}$, 分别考察其最小多项式 $p_A$ 与 $p_B$ 以及特征多项式 $\chi_A$ 与 $\chi_B$.

\noindent\textbf{目标:} 证明 $A$ 与 $B$ 相似, 当且仅当 $p_A = p_B$ 且 $\chi_A = \chi_B$.

\noindent\textbf{策略:} 首先设 $A$ 与 $B$ 相似, 比如 $B = S^{-1} \cdot A \cdot S$, 其中 $S \in \operatorname{Gl}(3, \mathbb{R})$ 可逆.则 $p_A(B) = S^{-1} \cdot p_A(A) \cdot S = 0$, 从而 $p_B \mid p_A$.相应地证明 $p_A \mid p_B$, 得 $p_A = p_B$.此外由定理~\mtref{thm:6.2.4} 读出 $\chi_A = \chi_B$.这些思考自然对任意域 $K$ 上的方阵更一般地成立.

现假设 $p_A = p_B$ 且 $\chi_A = \chi_B$, 我们首先从任意域 $K$ 上,具有一般行数与列数 $n$ 的矩阵 $A$, $B \in K^{n \times n}$ 出发.设 $p_A = p_B = p_1^{n_1} \cdot \ldots \cdot p_r^{n_r}$ 是素因子分解, 其中 $p_1, \ldots, p_r \in K[T]$ 是两两不同的素多项式, 指数 $n_i > 0$.则由大定理~\mtref{bigthm:6.5.11} 读出: 在 $A$ 的一般标准形以及 $B$ 的一般标准形中, 对每个 $i = 1, \ldots, r$, $p_i^{n_i}$ 的友矩阵至少作为对角方格出现一次.对 $n = 3$ 与 $K = \mathbb{R}$, 由此产生一定的限制, 尤其是利用特征多项式 $\chi_A = \chi_B$ 也是素多项式 $p_1, \ldots, p_r$ 的某些幂的乘积这一点.目标是证明 $A$ 与 $B$ 的一般标准形除方格的次序外一致.

\noindent\textbf{解题路径:} 我们上面已经核实: 相似的矩阵具有相同的最小多项式与特征多项式.因此还剩下要证明: 若 $p_A = p_B$ 且 $\chi_A = \chi_B$, 则 $A$, $B \in \mathbb{R}^{3 \times 3}$ 相似.

\begin{enumerate}
\item 情形: $\operatorname{grad} p_A = 1$, 比如 $p_A = T - \lambda$, 其中 $\lambda \in \mathbb{R}$ 是零点.则 $A = \operatorname{Diag}(\lambda, \lambda, \lambda)$, 即 $A$ 已处于一般标准形.
\item 情形: $\operatorname{grad} p_A = 2$.则推出 $p_A$ 可约.因为否则 $\chi_A$ 必须是 $p_A$ 的幂, 但这由次数理由被排除.设 $p_A = (T - \lambda_1) \cdot (T - \lambda_2)$, 其中 $\lambda_1, \lambda_2 \in \mathbb{R}$ 是零点.对 $\lambda_1 = \lambda_2$, 由大定理~\mtref{bigthm:6.5.11} 读出 $A$ 的一般标准形必为 $\operatorname{Diag}(\lambda_1, N)$ 的形式, 其中 $N$ 是 $(T - \lambda_1)^2$ 的友矩阵.反之若 $\lambda_1 \neq \lambda_2$, 且比如 $\chi_A = (T - \lambda_1)^2 \cdot (T - \lambda_2)$, 则 $A$ 的一般标准形具有 $\operatorname{Diag}(\lambda_1, \lambda_1, \lambda_2)$ 的形式; 参见大定理~\mtref{bigthm:6.5.11}.
\item 情形: $\operatorname{grad} p_A = 3$.这里可以考察 $p_A$ 的素因子分解的各种情形, 并特别利用基域 $\mathbb{R}$ 的性质.但更简单的是对由 $A$ 定义的自同态 $f \colon \mathbb{R}^3 \longrightarrow \mathbb{R}^3$ 应用大定理~\mtref{bigthm:6.5.18} 的 (iii).由次数理由有 $p_A = \chi_A$, 这按大定理~\mtref{bigthm:6.5.18} 的 (i) 给出: 在 $f$ 下把 $\mathbb{R}^3$ 看作 $\mathbb{R}[T]$-模, 它同构于 $\mathbb{R}[T]/p_A \mathbb{R}[T]$, 即 $f$-循环的.因此按引理~\mtref{lem:6.5.10} 存在 $\mathbb{R}^3$ 的基 $X$, 使得矩阵 $A_{f, X, X}$ 恰是 $p_A$ 的友矩阵.
\end{enumerate}

于是我们看到 $A$ 分别相似于一个结构仅依赖于 $p_A$ 与 $\chi_A$ 的矩阵.由于对 $B$ 代替 $A$ 也有同样的结论, 可以推出 $A$ 与 $B$ 相似.

\noindent\textbf{补充:} 给出的证明不需要 $\mathbb{R}$ 的任何特殊性质, 而是对任一基域 $K$ 代替 $\mathbb{R}$ 都成立.此外, 情形 3 中的论证表明, 引理~\mtref{lem:6.5.10} (i) 中 $V$ 为 $f$-循环的这一假设, 在 $\operatorname{grad} p_f = \dim_K V$ 时自动满足.

\subsection*{习题 6.5/6}

\noindent\textbf{起点:} 给定有限维 $K$-向量空间 $V$ 及其自同态 $f \colon V \longrightarrow V$.再设 $U \subset V$ 是 $f$-不变子空间.

\noindent\textbf{目标:} 证明下列断言:
\begin{enumerate}
\item[(i)] $f$ 诱导自同态 $\overline{f} \colon V/U \longrightarrow V/U$;
\item[(ii)] 对 $\overline{f}$ 与 $f$ 的最小多项式有 $p_{\overline{f}} \mid p_f$;
\item[(iii)] 对 $f$, $f|_U$ 与 $\overline{f}$ 的特征多项式有 $\chi_f = \chi_{f|_U} \cdot \chi_{\overline{f}}$.
\end{enumerate}

\noindent\textbf{策略:} 我们考虑自同态 $f$ 与典范投影 $\pi \colon V \longrightarrow V/U$ 的复合
\[
\tilde{f} \colon V \xrightarrow{f} V \xrightarrow{\pi} V/U\text{.}
\]
由于 $U$ 设为 $f$-不变, 有 $U \subset \ker \tilde{f}$, 于是按同态定理, $\tilde{f}$ 诱导线性映射 $\overline{f} \colon V/U \longrightarrow V/U$, 使图
\[
\begin{tikzcd}[row sep=huge, column sep=huge]
V \arrow[r, "f"] \arrow[d, "\pi"'] & V \arrow[d, "\pi"] \\
V/U \arrow[r, "\overline{f}"'] & V/U
\end{tikzcd}
\]
交换.此外 $p_f(f) = 0$, 从而 $p_f(\overline{f}) = 0$, 所以 $p_{\overline{f}}$ 是 $p_f$ 的因子.这解决了 (i) 与 (ii).

为处理 (iii), 必须回到定义~\mtref{def:6.2.2} 即定义~\mtref{def:6.2.6} 中特征多项式的定义.为简单起见我们一度假设 $V$ 中存在另一个 $f$-不变子空间 $U'$, 使 $V$ 分解为 $U$ 与 $U'$ 的直和, 即 $V = U \oplus U'$.若取 $U$ 与 $U'$ 的基并拼成 $V$ 的基, 则借助 \ref{sec:4.3} 节末尾的例~\mtref{bsp:4.3.2} 得到分解
\[
\chi_f = \chi_{f|_U} \cdot \chi_{f|_{U'}}\text{.}
\]
但限制 $\pi|_{U'} \colon U' \longrightarrow V/U$ 是同构, 所以图
\[
\begin{tikzcd}[row sep=huge, column sep=huge]
U' \arrow[r, "f|_{U'}"] \arrow[d, "\pi|_{U'}"'] & U' \arrow[d, "\pi|_{U'}"] \\
V/U \arrow[r, "\overline{f}"'] & V/U
\end{tikzcd}
\]
交换.推出 $\chi_{\overline{f}} = \chi_{f|_{U'}}$, 从而如所愿 $\chi_f = \chi_{f|_U} \cdot \chi_{\overline{f}}$.不过在一般情况下, $U$ 在 $V$ 中没有 $f$-不变的补.这意味着论证还需稍加细化.

\noindent\textbf{解题路径:} 断言 (i) 与 (ii) 是简单的, 上面已经论证.为证明 (iii), 选取 $U$ 在 $V$ 中的一个(不必要 $f$-不变的)补 $U'$.设 $X$ 是 $U$ 的基, $X'$ 是 $U'$ 的基.则 $X \cup X'$ 是 $V$ 的基, 且同构 $\pi|_{U'} \colon U' \longrightarrow V/U$ 表明 $\overline{X}'$ 作为 $X'$ 的像是 $V/U$ 的基.现在, 关于 $V$ 的基 $X \cup X'$ 描述 $f$ 的矩阵显然具有形式
\[
A_{f, X \cup X', X \cup X'} = \begin{pmatrix}
A_{f|_U, X, X} & * \\
0 & A_{\overline{f}, \overline{X}', \overline{X}'}
\end{pmatrix},
\]
并且借助 \ref{sec:4.3} 节末尾的例~\mtref{bsp:4.3.2}, 在合适的单位矩阵 $E$ 下可以如下计算:
\[
\begin{array}{l}
\det(TE - A_{f, X \cup X', X \cup X'}) = \det \begin{pmatrix}
TE - A_{f|_U, X, X} & * \\
0 & TE - A_{\overline{f}, \overline{X}', \overline{X}'}
\end{pmatrix} \\
= \det(TE - A_{f|_U, X, X}) \cdot \det(TE - A_{\overline{f}, \overline{X}', \overline{X}'})\text{.}
\end{array}
\]
但这如所愿意味着 $\chi_f = \chi_{f|_U} \cdot \chi_{\overline{f}}$.

% 第 7 章习题解答翻译: 欧几里得与酉向量空间
% 源文档 ap7_src.txt (对应附录习题训练第 7 章部分, 共 12 道题).

\section*{对应 7.1-7.6: 欧几里得与酉向量空间}

\subsection*{习题 7.1/3}

\noindent\textbf{起点:} 给定一个有限维 $\mathbb{R}$-向量空间 $V$ 及其上带正定的半双线性形式
$\varPhi \colon V \times V \longrightarrow \mathbb{R}$.再对向量 $x, y \in V$, $y \neq 0$,
考察多项式函数
\[
p(t) = |x + t y|^2 \quad \text{对} \ t \in \mathbb{R}\text{.}
\]

\noindent\textbf{目标:} 需确定 $p(t)$ 的零点.作为推论, 要在当前情形下推导出施瓦茨不等式,
即定理~\mtref{thm:7.1.4}.

\noindent\textbf{策略:} 成立
\[
p(t) = |x + t y|^2 = \langle x + t y, x + t y \rangle = \langle x, x \rangle + 2 t \langle x, y \rangle + t^2 \langle y, y \rangle,
\]
其中 $\langle y, y \rangle = |y|^2 \neq 0$, 因为双线性形式是正定的且 $y \neq 0$.为求 $p(t)$
的零点, 我们因此必须求解二次方程
\[
t^2 + 2 t \frac{\langle x, y \rangle}{|y|^2} + \frac{|x|^2}{|y|^2} = 0,
\]
通过配方法得到(复)解
\[
t_{1/2} = -\frac{\langle x, y \rangle}{|y|^2} \pm \sqrt{\frac{\langle x, y \rangle^2}{|y|^4} - \frac{|x|^2}{|y|^2}}\text{.}
\]
另一方面, 双线性形式是正定的.因此 $p(t)$ 恰好对满足 $x + t y = 0$ 的 $t \in \mathbb{R}$ 消失, 即 $x$ 线性依赖于 $y$.特别地, $p(t)$ 至多可以有一个实零点.为继续分析问题, 我们必须把这一信息与 $p(t)$ 对 $t \in \mathbb{C}$ 的零点公式化描述加以比较.

\noindent\textbf{解题路径:} 我们已经在上面计算了 $p(t)$ 的(复)零点.若 $x, y$ 线性无关, 则 $x + t y \neq 0$, 从而 $p(t) \neq 0$ 对一切 $t \in \mathbb{R}$ 成立, 因为 $\varPhi$ 是正定的.因此所求得的零点 $t_{1/2}$ 在此情形下不可能都是实的, 这对根号表达式的被开方数
\[
\frac{\langle x, y \rangle^2}{|y|^4} - \frac{|x|^2}{|y|^2}
\]
意味着它必须 $< 0$.于是得到
\[
\langle x, y \rangle^2 < |x|^2 \cdot |y|^2,
\]
即在 $x, y$ 线性无关的情形下的施瓦茨不等式.

现设 $x, y$ 线性相关, 即对某个参数 $t \in \mathbb{R}$ 成立 $x + t y = 0$, 其中我们利用了 $y \neq 0$.此时 $p(t)$ 恰有一个实的(因而是二重的)零点 $t_{1/2}$, 这意味着根号表达式的被开方数消失.从而得到
\[
t_{1/2} = -\frac{\langle x, y \rangle}{|y|^2}
\]
以及
\[
\langle x, y \rangle^2 = |x|^2 \cdot |y|^2\text{.}
\]
通过交换 $x$ 与 $y$ 的角色, 我们在 $V$ 上正定双线性形式的情形下总体得到施瓦茨不等式.即对任意向量 $x, y \in V$ 成立
\[
\langle x, y \rangle^2 \leq |x|^2 \cdot |y|^2,
\]
其中等号恰好当 $x$ 与 $y$ 线性相关时出现.

\subsection*{习题 7.1/5}

\noindent\textbf{起点:} 给定域 $K$ 上的一个向量空间 $V$.考察所有线性形式 $V \longrightarrow K$ 的
对偶空间 $V^*$, 以及所有 $K$-双线性映射 $V \times V \longrightarrow K$ 构成的 $K$-向量空间
$W$.$W$ 的向量空间结构如常定义, 即对 $\alpha \in K$ 与 $\varPhi \in W$ 通过
$(x, y) \longmapsto \alpha \cdot \varPhi(x, y)$ 定义 $\alpha \varPhi$, 两个元素
$\varPhi, \varPhi' \in W$ 的和通过 $(x, y) \longmapsto \varPhi(x, y) + \varPhi'(x, y)$ 定义.

\noindent\textbf{目标:} 需证明映射
\[
\varPsi \colon \operatorname{Hom}_K(V, V^*) \longrightarrow W, \qquad \varphi \longmapsto \bigl((x, y) \longmapsto \varphi(x)(y)\bigr),
\]
是 $K$-向量空间的同构.

\noindent\textbf{策略:} 我们试图推导出向量空间同构所要求的性质, 即线性性,单射性与满射性.给定两个
同态 $\varphi, \varphi' \in \operatorname{Hom}_K(V, V^*)$ 以及一个标量 $\alpha \in K$.则成立
\[
\begin{array}{c}
\alpha \varphi \longmapsto^{\varPsi} \bigl((x, y) \longmapsto (\alpha \varphi)(x)(y) = \alpha \cdot \bigl(\varphi(x)(y)\bigr)\bigr), \\[2pt]
\varphi + \varphi' \longmapsto^{\varPsi} \bigl((x, y) \longmapsto (\varphi + \varphi')(x)(y) = \varphi(x)(y) + \varphi'(x)(y)\bigr),
\end{array}
\]
而这意味着 $\varPsi(\alpha \varphi) = \alpha \varPsi(\varphi)$ 以及
$\varPsi(\varphi + \varphi') = \varPsi(\varphi) + \varPsi(\varphi')$.因此 $\varPsi$ 是
$K$-向量空间的同态.现设 $\varPsi(\varphi) = 0$, 即对一切 $x, y \in V$ 有
$\varphi(x)(y) = 0$.则对一切 $x \in V$ 在 $V^*$ 中有 $\varphi(x) = 0$, 从而在
$\operatorname{Hom}_K(V, V^*)$ 中有 $\varphi = 0$.故 $\varPsi$ 是单射.最后设
$\varPhi \in W$, 即 $\varPhi \colon V \times V \longrightarrow K$ 是一个 $K$-双线性映射.则
$\varPhi$ 对每个 $x \in V$ 定义一个线性形式
\[
\varPhi(x, \cdot) \colon V \longrightarrow K, \quad y \longmapsto \varPhi(x, y),
\]
并且
\[
\varphi \colon V \longrightarrow V^*, \quad x \longmapsto \varPhi(x, \cdot),
\]
是一个满足 $\varPsi(\varphi) = \varPhi$ 的 $K$-线性映射.因此 $\varPsi$ 是满射, 从而总体上是
$K$-向量空间的同构.我们将在下文中看到, 刚才为证明 $\varPsi$ 的满射性而进行的构造可以用于简化论证.

\noindent\textbf{解题路径:} 如上面一样验证
$\varPsi \colon \operatorname{Hom}_K(V, V^*) \longrightarrow W$ 是 $K$-向量空间的同态.对映射
\[
\varPsi' \colon W \longrightarrow \operatorname{Hom}_K(V, V^*), \quad \varPhi \longmapsto \bigl(x \longmapsto \varPhi(x, \cdot)\bigr),
\]
同样如此, 其中对 $x \in V$, $\varPhi(x, \cdot) \colon V \longrightarrow K$ 如上通过
$y \longmapsto \varPhi(x, y)$ 定义.但此时显然 $\varPsi' \circ \varPsi$ 是
$\operatorname{Hom}_K(V, V^*)$ 上的恒等映射, 而 $\varPsi \circ \varPsi'$ 是 $W$ 上的恒等映射.
因此 $\varPsi$ 与 $\varPsi'$ 互为逆映射, 特别地推出 $\varPsi$ 与 $\varPsi'$ 是 $K$-向量空间的同构.

\subsection*{习题 7.2/2}

\noindent\textbf{起点:} 对 $n \in \mathbb{N}$, 用 $\mathbb{R}[T]_n$ 表示 $\mathbb{R}[T]$ 中所有
次数 $\leq n$ 的多项式构成的 $\mathbb{R}$-向量空间.继而考察映射
\[
\varPhi \colon \mathbb{R}[T]_n \times \mathbb{R}[T]_n \longrightarrow \mathbb{R}, \quad (f, g) \longmapsto \langle f, g \rangle := \int_0^1 f(t) g(t)\, dt\text{.}
\]

\noindent\textbf{目标:} 需证明 $\varPhi$ 在 $\mathbb{R}[T]_n$ 上定义了一个数量积.再借助 E. 施密特
正交规范化方法, 由 $\mathbb{R}[T]_2$ 的基 $1, T, T^2$ 构造 $\mathbb{R}[T]_2$ 的一个标准正交基.

\noindent\textbf{策略:} 这里我们自然要动用分析中的积分概念.即由微积分基本定理, 对连续函数
$f \colon [0, 1] \longrightarrow \mathbb{R}$ 及 $f$ 的一个原函数 $F$(即一个连续函数
$F \colon [0, 1] \longrightarrow \mathbb{R}$, 它在开区间 $]0, 1[$ 内可导并且在那里以 $f$ 为导数)成立
\[
\int_0^1 f(t)\, dt = F(1) - F(0)\text{.}
\]
由导数的运算法则, 例如对函数 $F, G$ 及常数 $\alpha \in \mathbb{R}$ 有
$(F + G)' = F' + G'$, $(\alpha F)' = \alpha F'$, 得到相应的积分运算法则, 即
\[
\int_0^1 \bigl(f(t) + g(t)\bigr)\, dt = \int_0^1 f(t)\, dt + \int_0^1 g(t)\, dt, \qquad
\int_0^1 \alpha f(t)\, dt = \alpha \int_0^1 f(t)\, dt\text{.}
\]
在我们具体的情形中涉及的是多项式函数, 它们自动连续, 而且对它们来说原函数的存在不成为问题.这样就
预示了本题的解题路径.剩下的就是例行验证, 最后是执行 E. 施密特正交规范化方法.

\noindent\textbf{解题路径:} 由上述积分法则推出, $\varPhi$ 是一个对称的 $\mathbb{R}$-双线性映射.此外,
若 $f$ 不是零多项式, 则成立
\[
\langle f, f \rangle = \int_0^1 f(t)^2\, dt > 0\text{.}
\]
例如可由如下事实推出这一点: $f(t)^2$ 至多只有孤立零点, 除此之外恒 $> 0$, 因而 $f(t)^2$ 的每个
原函数都严格单调递增.特别地, 我们认识到 $\varPhi$ 是正定的半双线性形式, 从而是数量积.

为了把 $\mathbb{R}[T]_2$ 的基 $(x_1, x_2, x_3) = (1, T, T^2)$ 变换为一个标准正交基
$(e_1, e_2, e_3)$, 我们按照引理~\mtref{lem:7.2.4} 与定理~\mtref{thm:7.2.5} 所述的方法进行, 并从向量
$x_1 = 1$ 开始.因为
\[
\langle 1, 1 \rangle = \int_0^1 dt = 1,
\]
它已经归一化, 所以我们令 $e_1 = 1$.进一步成立
\[
\begin{array}{c}
\langle x_2, e_1 \rangle = \int_0^1 t\, dt = \left[ \frac{1}{2} t^2 \right]_0^1 = \frac{1}{2}, \\[2pt]
\left| x_2 - \frac{1}{2} e_1 \right|^2 = \int_0^1 \left( t - \frac{1}{2} \right)^2 dt = \left[ \frac{1}{3} t^3 - \frac{1}{2} t^2 + \frac{1}{4} t \right]_0^1 = \frac{1}{12},
\end{array}
\]
并且得到 $e_2 = \sqrt{12} \cdot \left( T - \frac{1}{2} \right)$.接下来计算
\[
\begin{array}{r l}
& \langle x_3, e_1 \rangle = \int_0^1 t^2\, dt = \left[ \frac{1}{3} t^3 \right]_0^1 = \frac{1}{3}, \\[2pt]
& \langle x_3, e_2 \rangle = \sqrt{12} \cdot \int_0^1 \left( t^3 - \frac{1}{2} t^2 \right) dt = \sqrt{12} \cdot \left[ \frac{1}{4} t^4 - \frac{1}{6} t^3 \right]_0^1 = \frac{1}{\sqrt{12}},
\end{array}
\]
并从 $x_3$ 中减去它在 $\mathbb{R}[T]_2$ 中由 $e_1, e_2$(或 $x_1, x_2$)\emph{张成的线性子空间}(\textup{aufgespannter Untervektorraum})\index{zhangchengdexianxingzikongjian@张成的线性子空间(aufgespannter Untervektorraum)}上的正交投影.
我们得到
\[
x_3' = x_3 - \frac{1}{\sqrt{12}} \cdot e_2 - \frac{1}{3} \cdot e_1 = T^2 - \left( T - \frac{1}{2} \right) - \frac{1}{3} = T^2 - T + \frac{1}{6},
\]
其中模的平方计算为
\[
\begin{array}{c}
|x_3'|^2 = \int_0^1 \left( t^2 - t + \frac{1}{6} \right)^2 dt = \int_0^1 \left( t^4 - 2 t^3 + \frac{4}{3} t^2 - \frac{1}{3} t + \frac{1}{36} \right) dt \\[2pt]
= \frac{1}{5} - \frac{1}{2} + \frac{4}{9} - \frac{1}{6} + \frac{1}{36} = \frac{1}{180}\text{.}
\end{array}
\]
于是 $\mathbb{R}[T]_2$ 的所求标准正交基由多项式
\[
e_1 = 1, \qquad e_2 = 2 \cdot \sqrt{3} \cdot \left( T - \frac{1}{2} \right), \qquad e_3 = 6 \cdot \sqrt{5} \cdot \left( T^2 - T + \frac{1}{6} \right)
\]
组成.

\subsection*{习题 7.2/4}

\noindent\textbf{起点:} 给定一个欧几里得或酉 $K$-向量空间 $V$, 一个线性子空间 $U \subset V$ 以及一个
向量 $v \in V - U$.

\noindent\textbf{目标:} 需证明存在唯一一个向量 $u_0 \in U$ 使得 $v - u_0 \in U^{\perp}$, 并且对一切
其他向量 $u \in U$ 成立不等式 $|v - u| > |v - u_0|$.

\noindent\textbf{策略:} 从根本上说, 这个问题已在引理~\mtref{lem:7.2.4} 中处理过.需给出一个向量
$u_0 \in U$, 使 $v - u_0$ 与 $U$ 中所有向量正交.后者可以通过对 $U$ 的一个基的元素
$x_1, \ldots, x_r$ 验证 $\langle v - u_0, x_i \rangle = 0$ 来检验.关于正交性或模的问题, 在取定
标准正交基的前提下最容易回答.因此很自然地选取 $U$ 的一个标准正交基 $e_1, \ldots, e_r$, 并用元素
$e_{r+1}, \ldots, e_n$ 把它扩充为 $V$ 的一个标准正交基; 参见推论~\mtref{cor:7.2.6}.则
$e_{r+1}, \ldots, e_n$ 张成 $U$ 的正交补 $U^{\perp}$, 正如在推论~\mtref{cor:7.2.8} 的证明中所证的那样.

\noindent\textbf{解题路径:} 我们选取 $U$ 的一个标准正交基 $e_1, \ldots, e_r$, 并用元素
$e_{r+1}, \ldots, e_n$ 把它扩充为 $V$ 的一个标准正交基.若此时 $v = \sum_{i = 1}^{n} \alpha_i e_i$
带有系数 $\alpha_i \in \mathbb{K}$, 则取 $u_0 = \sum_{i = 1}^{r} \alpha_i e_i$, 显然成立
\[
v - u_0 = \sum_{i = r + 1}^{n} \alpha_i e_i \in U^{\perp}\text{.}
\]
对任意 $u \in U$, 例如 $u = \sum_{i = 1}^{r} \beta_i e_i$ 带有系数 $\beta_i \in \mathbb{K}$, 此时成立
\[
|v - u|^2 = \sum_{i = 1}^{r} (\alpha_i - \beta_i)^2 + \sum_{i = r + 1}^{n} \alpha_i^2 \geq \sum_{i = r + 1}^{n} \alpha_i^2 = |v - u_0|^2,
\]
其中等号等价于项 $\alpha_i - \beta_i$($i = 1, \ldots, r$)全部消失.因此如所求地推出, 对一切满足
$u \neq u_0$ 的 $u \in U$ 有 $|v - u| > |v - u_0|$.

\noindent\textbf{补充:} 与 E. 施密特正交规范化方法完全一样, 向量 $v - u_0$ 的构造方式是: 从 $v$ 中
减去 $v$ 在线性子空间 $U$ 上的正交投影 $u_0$.于是 $v - u_0$ 垂直于 $U$, 并且正如我们所见, 它的长度
正是 $v$ 与 $U$ 之间的最小距离, 如果我们在点空间的意义下论证的话.

\subsection*{习题 7.3/4}

\noindent\textbf{起点:} 给定带基 $X$ 的有限维 $K$-向量空间 $V$.再设 $\varPhi$ 是 $V$ 上的一个半双线性
形式或埃尔米特形式.

\noindent\textbf{目标:} 对正半定的 $\varPhi$, 需证明描述矩阵 $A_{\varPhi, X}$ 的一切主子式都是实的且
$\geq 0$.反过来要考虑: 当 $A_{\varPhi, X}$ 的一切主子式都是实的且 $\geq 0$ 时, $\varPhi$ 是否已经
是正半定的.

\noindent\textbf{策略:} 我们考察一个简单的例子, 即域 $\mathbb{K} = \mathbb{R}$ 上的向量空间
$V = \mathbb{R}^2$.设 $\varPhi$ 是 $V$ 上的一个半双线性形式, 例如关于 $V$ 的典范基 $e$ 用矩阵
\[
A_{\varPhi, e} = \begin{pmatrix} \alpha & \gamma \\ \gamma & \beta \end{pmatrix} \in \mathbb{R}^{2 \times 2}
\]
来描述, 其中我们已考虑 $A_{\varPhi, e} = (A_{\varPhi, e})^t$.于是对 $a = (a_1, a_2)^t \in V$,
$\varPhi(a, a)$ 计算为
\[
\varPhi(a, a) = \alpha a_1^2 + 2 \gamma a_1 a_2 + \beta a_2^2\text{.}
\]
现设 $\varPhi$ 正半定, 即对一切 $a \in V$ 有 $\varPhi(a, a) \geq 0$.例如令 $a = (1, 0)^t$, 则得
$\alpha \geq 0$, 即 $A_{\varPhi, e}$ 的第一个主子式 $\geq 0$.此时甚至成立 $\alpha > 0$, 因为否则对
类型 $a = (a_1, 1)^t$ 的向量, $\varPhi(a, a)$ 将计算为 $\varPhi(a, a) = 2 \gamma a_1 + \beta$, 而这对
$a_1 \in \mathbb{R}$ 不会只取 $\geq 0$ 的值.

于是我们已经看到 $\alpha > 0$, 并且可以同样地推出 $\beta > 0$.随后存在向量
$a = \left( \frac{1}{\sqrt{\alpha}}, -\frac{1}{\sqrt{\beta}} \right)^t \in V$, 并且成立
\[
\varPhi(a, a) = 2 - 2 \frac{\gamma}{\sqrt{\alpha \beta}} \geq 0,
\]
由此得到 $\gamma \leq \sqrt{\alpha \beta}$, 即 $\gamma^2 \leq \alpha \beta$, 从而
$\alpha \beta - \gamma^2 \geq 0$.但这意味着矩阵 $A_{\varPhi, e}$ 的第二个主子式也是非负的.这样就在
当前特殊情形下证明了: 对正半定的 $\varPhi$, 描述矩阵 $A_{\varPhi, e}$ 的主子式是非负的.不过可以想象,
在更高维数下, 类似的通过显式计算进行的论证不再可行.因此这里需要的是一个更一般类型的论证.

此外, 我们在上述例子中还看到, 由 $\varPhi$ 正半定特别地推出条件 $\alpha > 0$,$\beta > 0$ 以及
$\alpha \beta - \gamma^2 \geq 0$.另一方面, 要求 $A_{\varPhi, e}$ 的主子式非负, 只意味着 $\alpha \geq 0$
以及 $\alpha \beta - \gamma^2 \geq 0$.这使人预料, 当某个相关矩阵的主子式非负时, 半双线性形式
$\varPhi$ 不一定会自动正半定.

\noindent\textbf{解题路径:} 我们把所考虑的 $V$ 的基记作 $X = (x_1, \ldots, x_n)$, 并对 $i = 1, \ldots, n$
令 $X_i = (x_1, \ldots, x_i)$.则 $X_i$ 分别是 $x_1, \ldots, x_i$ 生成的线性子空间 $V_i \subset V$ 的基.
现设 $\varPhi$ 正半定, 并设 $\varPhi_i$($i = 1, \ldots, n$)是 $\varPhi$ 到 $V_i$ 上的限制.推出
$\varPhi_i$ 也正半定.若 $\varPhi_i$ 甚至正定, 则由定理~\mtref{thm:7.3.8} 得
$\det(A_{\varPhi_i, X_i}) > 0$.否则存在一个向量 $x \in V_i - \{0\}$ 使得
$\varPhi_i(x, x) = 0$, 于是 $\varPhi_i$ 是退化的; 参见推论~\mtref{cor:7.1.5}.在此情形下, 按
定理~\mtref{thm:7.3.2} 成立 $\det(A_{\varPhi_i, X_i}) = 0$, 因此在两种情形下都有
$\det(A_{\varPhi_i, X_i}) \geq 0$.由于 $\det(A_{\varPhi_i, X_i})$ 恰好是 $A_{\varPhi, X}$ 的第 $i$ 个
主子式, 我们看到 $A_{\varPhi, X}$ 的一切主子式都是非负的.

为了最后回答这样的问题: 当描述矩阵 $A_{\varPhi, X}$ 的全部主子式非负时, $V$ 上的半双线性形式
$\varPhi$ 是否已经正半定, 我们考察 $\mathbb{R}$-向量空间 $V = \mathbb{R}^2$ 连同典范基 $e$, 以及 $V$
上由矩阵
\[
A_{\varPhi, e} = \begin{pmatrix} 0 & 0 \\ 0 & -1 \end{pmatrix}
\]
给出的半双线性形式 $\varPhi$.$A_{\varPhi, e}$ 的一切主子式都消失, 因此都 $\geq 0$, 但尽管如此, 对
向量 $a = (0, 1)^t \in V$ 仍成立等式 $\varPhi(a, a) = -1 < 0$, 即 $\varPhi$ 不是正半定的.

\subsection*{习题 7.3/7}

\noindent\textbf{起点:} 对矩阵 $A \in \mathbb{K}^{n \times n}$, 考察 $\mathbb{K}^n$ 上由
\[
\langle a, b \rangle = a^t \cdot A \cdot \overline{b}
\]
定义的半双线性形式 $\varPhi$.

\noindent\textbf{目标:} 需证明 $\varPhi$ 恰好当存在可逆矩阵 $S \in \mathrm{Gl}(n, \mathbb{K})$ 使得
$A = S^t \cdot \overline{S}$ 时, 在 $\mathbb{K}^n$ 上定义一个数量积.

\noindent\textbf{策略:} 要讨论的分解 $A = S^t \cdot \overline{S}$, 即带单位矩阵
$E \in \mathbb{K}^{n \times n}$ 的 $A = S^t \cdot E \cdot \overline{S}$, 使我们联想到定理~\mtref{thm:7.3.5}
中的公式
\[
A_{\varPhi, Y} = A_{Y, X}^t \cdot A_{\varPhi, X} \cdot \overline{A}_{Y, X},
\]
它刻画了从 $\mathbb{K}^n$ 的基 $X$ 过渡到新基 $Y$ 时, $\varPhi$ 的描述矩阵 $A_{\varPhi, X}$ 的变化行为.
此外, 在当前情形下 $A$ 正好是 $\varPhi$ 关于 $\mathbb{K}^n$ 的典范基 $e$ 的描述矩阵, 即
$A_{\varPhi, e} = A$.单位矩阵 $E$ 在这个关联中也有意义.它恰好当 $\varPhi$ 存在标准正交基时, 才可作为
$\varPhi$ 关于某个合适基的描述矩阵.基于这些预备考虑, 解题路径已或多或少地确定了.

\noindent\textbf{解题路径:} 我们首先假定 $\varPhi$ 在 $\mathbb{K}^n$ 上定义了一个数量积.则按
定理~\mtref{thm:7.2.5} 存在 $\varPhi$ 的一个标准正交基 $X$, 并且对 $\mathbb{K}^n$ 的典范基 $e$ 与 $X$
之间的基变换矩阵 $A_{e, X}$ 成立关系
\[
A = A_{\varPhi, e} = A_{e, X}^t \cdot A_{\varPhi, X} \cdot \overline{A}_{e, X} = A_{e, X}^t \cdot \overline{A}_{e, X},
\]
因为 $X$ 是标准正交基, 从而 $A_{\varPhi, X}$ 是单位矩阵.于是得到 $A = S^t \cdot \overline{S}$, 其中
$S = A_{e, X} \in \mathrm{Gl}(n, \mathbb{K})$; 参见注记~\mtref{rem:3.4.2}.

反之, 若对可逆矩阵 $S \in \mathrm{Gl}(n, \mathbb{K})$ 成立 $A = S^t \cdot \overline{S}$, 则可以把 $S$
解释为基变换矩阵 $A_{e, Y}$, 它描述从典范基 $e$ 到 $\mathbb{K}^n$ 的另一个基 $Y$ 的过渡, 即
$A = A_{e, Y}^t \cdot \overline{A}_{e, Y}$.现在由注记~\mtref{rem:3.4.2} 有
\[
A_{\varPhi, Y} = A_{Y, e}^t \cdot A_{\varPhi, e} \cdot \overline{A}_{Y, e} = A_{Y, e}^t \cdot A \cdot \overline{A}_{Y, e} = (A_{e, Y} \cdot A_{Y, e})^t \cdot \overline{(A_{e, Y} \cdot A_{Y, e})} = E,
\]
因此 $Y$ 是关于 $\varPhi$ 的标准正交基.由此可以容易地推出 $\varPhi$ 具有数量积的性质.设例如
$Y = (y_1, \ldots, y_n)$.对带系数 $\alpha_i, \beta_i \in \mathbb{K}$ 的向量
$a = \sum_{i = 1}^{n} \alpha_i y_i$ 与 $b = \sum_{i = 1}^{n} \beta_i y_i$, 此时得到
\[
\varPhi(a, b) = \sum_{i = 1}^{n} \alpha_i \overline{\beta}_i = \overline{\sum_{i = 1}^{n} \overline{\alpha}_i \beta_i} = \overline{\varPhi(b, a)},
\]
以及
\[
\varPhi(a, a) = \sum_{i = 1}^{n} \alpha_i \overline{\alpha}_i = \sum_{i = 1}^{n} |\alpha_i|^2 > 0 \quad \text{对} \ a \neq 0,
\]
我们看到 $\varPhi$ 在 $\mathbb{K}^n$ 上定义了一个数量积.

\subsection*{习题 7.4/1}

\noindent\textbf{起点:} 给定一个有限维欧几里得或酉 $K$-向量空间 $V$, 一个自同态
$\varphi \colon V \longrightarrow V$ 及其伴随映射 $\varphi^* \colon V \longrightarrow V$.

\noindent\textbf{目标:} 需证明 $\mathrm{Spur}(\varphi \circ \varphi^*) \geq 0$, 并证明
$\mathrm{Spur}(\varphi \circ \varphi^*) = 0$ 等价于 $\varphi = 0$.

\noindent\textbf{策略:} 自同态 $f \colon V \longrightarrow V$ 的迹由 $V$ 的任意一个基 $X$ 对应的描述矩阵
$A_{f, X, X}$ 的迹给出; 参见定义~\mtref{def:6.2.6}.此外, 方阵的迹按定理~\mtref{thm:6.2.3} 定义为其对角线
元素之和.在当前情形下, 宜选取 $X$ 为 $V$ 的标准正交基, 因为此时按注记~\mtref{rem:7.4.4} 有
$A_{\varphi^*, X, X} = (A_{\varphi, X, X})^*$, 并且只需计算乘积矩阵
$A_{\varphi, X, X} \cdot (A_{\varphi, X, X})^*$ 的迹.

\noindent\textbf{解题路径:} 我们选取 $V$ 的一个标准正交基 $X$, 并把 $\varphi$ 关于 $X$ 用矩阵表示, 例如
$A_{\varphi, X, X} = (\alpha_{ij})_{i, j = 1, \ldots, n}$.则由注记~\mtref{rem:7.4.4} 得
$A_{\varphi^*, X, X} = (A_{\varphi, X, X})^* = (\overline{\alpha}_{ij})_{j, i = 1, \ldots, n}$, 并且成立
\[
A_{\varphi, X, X} \cdot (A_{\varphi, X, X})^* = (\alpha_{ij})_{i, j = 1, \ldots, n} \cdot (\overline{\alpha}_{ij})_{j, i = 1, \ldots, n} = \left( \sum_{j = 1}^{n} \alpha_{ij} \cdot \overline{\alpha}_{kj} \right)_{i, k = 1, \ldots, n}\text{.}
\]
于是得到
\[
\operatorname{Spur}\bigl(A_{\varphi, X, X} \cdot (A_{\varphi, X, X})^*\bigr) = \sum_{i = 1}^{n} \sum_{j = 1}^{n} \alpha_{ij} \cdot \overline{\alpha}_{ij} = \sum_{i = 1}^{n} \sum_{j = 1}^{n} |\alpha_{ij}|^2 \geq 0,
\]
并且显然
\[
\operatorname{Spur}\bigl(A_{\varphi, X, X} \cdot (A_{\varphi, X, X})^*\bigr) = 0
\]
等价于对一切 $i, j$ 有 $\alpha_{ij} = 0$, 即等价于 $\varphi = 0$.

\subsection*{习题 7.4/4}

\noindent\textbf{起点:} 考察一个有限维酉 $\mathbb{C}$-向量空间 $V$ 及其自同态
$\varphi \colon V \longrightarrow V$ 和相应的伴随映射 $\varphi^* \colon V \longrightarrow V$.

\noindent\textbf{目标:} 证明 $\varphi$ 恰好当存在多项式 $p \in \mathbb{C}[T]$ 使得
$\varphi^* = p(\varphi)$ 时是正规的.

\noindent\textbf{策略:} 若存在多项式 $p \in \mathbb{C}[T]$ 使得 $\varphi^* = p(\varphi)$, 则平凡地成立
\[
\varphi \circ \varphi^* = \varphi \circ p(\varphi) = p(\varphi) \circ \varphi = \varphi^* \circ \varphi,
\]
即 $\varphi$ 是正规的.反过来, 现设 $\varphi$ 是正规的.显然没有任何一般理由说明 $\varphi$ 的伴随映射
应当是 $\varphi$ 的一个多项式.因此我们必须尽可能充分地利用这里特殊给定的情形.由于我们在域
$\mathbb{C}$ 上工作, 特征多项式 $\chi_\varphi$ 完全分裂为一次因式.于是按定理~\mtref{thm:7.4.8} 可以从
$\varphi$ 的正规性推出: 在 $V$ 中存在一个标准正交基 $X$, 它全部由 $\varphi$ 的特征向量组成, 从而按
推论~\mtref{cor:7.4.7} 也全部由 $\varphi^*$ 的特征向量组成.更精确地说,
$A_{\varphi, X, X}$ 是对角矩阵, 例如 $A_{\varphi, X, X} = \mathrm{Diag}(\lambda_1, \ldots, \lambda_n)$,
并且由注记~\mtref{rem:7.4.4} 得
$A_{\varphi^*, X, X} = (A_{\varphi, X, X})^* = \mathrm{Diag}(\overline{\lambda}_1, \ldots, \overline{\lambda}_n)$.
现在我们还需要如下观察: 对多项式 $p \in \mathbb{C}[T]$ 成立关系
\[
p(A_{\varphi, X, X}) = \mathrm{Diag}\bigl(p(\lambda_1), \ldots, p(\lambda_n)\bigr)\text{.}
\]
这样就预示了解题路径.需要找到一个多项式 $p \in \mathbb{C}[T]$, 使得对 $i = 1, \ldots, n$ 有
$p(\lambda_i) = \overline{\lambda}_i$.

\noindent\textbf{解题路径:} 如已说明的, 当存在多项式 $p \in \mathbb{C}[T]$ 使 $\varphi^* = p(\varphi)$ 时,
$\varphi$ 是正规的.反之, 若假定 $\varphi$ 正规, 则可以按定理~\mtref{thm:7.4.8} 在 $V$ 中选取一个标准正交
基 $X$, 使得 $A_{\varphi, X, X}$ 是对角矩阵, 例如 $A_{\varphi, X, X} = \mathrm{Diag}(\lambda_1, \ldots, \lambda_n)$.
再利用自同态与描述矩阵之间的双射对应, 即定理~\mtref{thm:3.3.2}, 问题按推论~\mtref{cor:7.4.7} 归结为: 找到
一个多项式 $p \in \mathbb{C}[T]$, 它满足关系
\[
\mathrm{Diag}\bigl(p(\lambda_1), \ldots, p(\lambda_n)\bigr) = p(A_{\varphi, X, X}) = A_{\varphi^*, X, X} = \mathrm{Diag}\bigl(\overline{\lambda}_1, \ldots, \overline{\lambda}_n\bigr)\text{.}
\]
但现在很容易给出 $\mathbb{C}[T]$ 中的一个多项式, 它在两两不同的位置 $\lambda_1, \ldots, \lambda_r \in \mathbb{C}$
(即去掉重复列举后的 $\lambda_1, \ldots, \lambda_n$) 上取预先指定的值, 在我们的情形下即
$\overline{\lambda}_1, \ldots, \overline{\lambda}_r$.事实上,
\[
p_i = c_i \cdot \prod_{\substack{j = 1, \ldots, r \\ j \neq i}} (T - \lambda_j), \qquad c_i \in \mathbb{C},
\]
是 $\mathbb{C}[T]$ 中的一个多项式, 它在一切位置 $\lambda_j$($j = 1, \ldots, r$)上消失, 但可能除位置
$\lambda_i$ 外, 而恰当地选取常数 $c_i$ 可以做到使 $p_i(\lambda_i) = \overline{\lambda}_i$.于是
$p = \sum_{i = 1}^{r} p_i$ 是 $\mathbb{C}[T]$ 中的一个多项式, 它在位置 $\lambda_1, \ldots, \lambda_n$ 上取
值 $\overline{\lambda}_1, \ldots, \overline{\lambda}_n$, 从而如所求地推出 $\varphi^* = p(\varphi)$.

\subsection*{习题 7.5/2}

\noindent\textbf{起点:} 给定两个旋转矩阵 $R(\vartheta_1), R(\vartheta_2) \in \mathbb{R}^{2 \times 2}$, 其
角度满足 $0 \leq \vartheta_1 < \vartheta_2 < 2 \pi$.

\noindent\textbf{目标:} 需证明 $R(\vartheta_1)$ 与 $R(\vartheta_2)$ 恰好当 $\vartheta_1 + \vartheta_2 = 2 \pi$
时相似.

\noindent\textbf{策略:} 定理~\mtref{thm:7.5.8} 的叙述及其相应证明中含有对解决本问题有用的基本要素.

\noindent\textbf{解题路径:} 首先设 $\vartheta_1 + \vartheta_2 = 2 \pi$.由于情形 $\vartheta_1 = \vartheta_2 = \pi$
已被排除, 我们可以例如假定 $\vartheta_1 < \pi$, 从而 $\pi < \vartheta_2 < 2 \pi$.考察由 $R(\vartheta_1)$ 定义的
$\mathbb{R}$-线性映射
\[
f \colon \mathbb{R}^2 \longrightarrow \mathbb{R}^2, \qquad x \longmapsto R(\vartheta_1) \cdot x = \begin{pmatrix} \cos \vartheta_1 & -\sin \vartheta_1 \\ \sin \vartheta_1 & \cos \vartheta_1 \end{pmatrix} \cdot x,
\]
我们得到 $A_{f, e, e} = R(\vartheta_1)$ 为关于 $\mathbb{R}^2$ 的典范基 $\boldsymbol{e} = (e_1, e_2)$ 的
描述矩阵.则对 $\mathbb{R}^2$ 的基 $\boldsymbol{e}' = (e_2, e_1)$ 显然成立
\[
A_{f, e', e'} = \begin{pmatrix} \cos \vartheta_1 & \sin \vartheta_1 \\ -\sin \vartheta_1 & \cos \vartheta_1 \end{pmatrix} = \begin{pmatrix} \cos \vartheta_2 & -\sin \vartheta_2 \\ \sin \vartheta_2 & \cos \vartheta_2 \end{pmatrix} = R(\vartheta_2),
\]
其中我们利用了 $\vartheta_2 = 2 \pi - \vartheta_1$, 以及 $\cos(-\vartheta_1) = \cos \vartheta_1$,
$\sin(-\vartheta_1) = -\sin \vartheta_1$.但这意味着 $R(\vartheta_1)$ 与 $R(\vartheta_2)$ 相似.

反之, 现设 $R(\vartheta_1)$ 与 $R(\vartheta_2)$ 相似.则由定理~\mtref{thm:6.2.4}, 相应的特征多项式一致, 即
\[
(T - e^{i \vartheta_1})(T - e^{-i \vartheta_1}) = \chi_{R(\vartheta_1)} = \chi_{R(\vartheta_2)} = (T - e^{i \vartheta_2})(T - e^{-i \vartheta_2})\text{.}
\]
素因子分解的唯一性可以推出 $e^{i \vartheta_1} = e^{\pm i \vartheta_2}$, 并且在限制
$0 \leq \vartheta_1 < \vartheta_2 < 2 \pi$ 下推出 $\vartheta_1 = \vartheta_2$ 或
$\vartheta_1 = 2 \pi - \vartheta_2$.但由于 $\vartheta_1 = \vartheta_2$ 已被排除, 得到
$\vartheta_1 + \vartheta_2 = 2 \pi$.

我们也可以利用这样的事实: 相似的矩阵具有相同的迹; 参见推论~\mtref{cor:6.2.5}.则得到条件
$\cos \vartheta_1 = \cos \vartheta_2$, 由它可以类似地推出 $\vartheta_1 + \vartheta_2 = 2 \pi$.

\subsection*{习题 7.5/4}

\noindent\textbf{起点:} 给定一个正交矩阵 $A = (\alpha_{ij}) \in \mathrm{O}(n)$, 即满足 $A^t = A^{-1}$
的实矩阵 $A \in \mathbb{R}^{n \times n}$.设 $A$ 具有下三角形, 即对 $i < j$ 成立
$\alpha_{ij} = 0$.

\noindent\textbf{目标:} 需证明 $A$ 甚至是对角矩阵.此外还要问: 对酉矩阵 $A \in \mathrm{U}(n)$ 是否也有
相应的断言成立.

\noindent\textbf{策略:} $n = 1$ 的情形是平凡的.对 $n = 2$, 我们已在定理~\mtref{thm:7.5.6} 中显式描述了
$\mathrm{O}(2)$ 中所有矩阵.特别地, 由此读出, $\mathrm{O}(2)$ 中每个下三角矩阵都已经是对角矩阵.现设
对一般的 $n$, $A = (\alpha_{ij}) \in \mathrm{O}(n)$.则 $A^t = A^{-1}$, 从而对单位矩阵
$E \in \mathbb{R}^n$ 成立 $A \cdot A^t = E$, 而这意味着:
\[
\begin{pmatrix}
\alpha_{11} & 0 & 0 & \cdots & 0 \\
\alpha_{21} & \alpha_{22} & 0 & \cdots & 0 \\
\vdots & \vdots & \vdots & \ddots & \vdots \\
\alpha_{n1} & \alpha_{n2} & \alpha_{n3} & \cdots & \alpha_{nn}
\end{pmatrix}
\cdot
\begin{pmatrix}
\alpha_{11} & \alpha_{21} & \cdots & \cdots & \alpha_{n1} \\
0 & \alpha_{22} & \cdots & \cdots & \alpha_{n2} \\
\vdots & \vdots & \ddots & \ddots & \vdots \\
0 & 0 & \cdots & 0 & \alpha_{nn}
\end{pmatrix}
=
\begin{pmatrix}
1 & 0 & \cdots & 0 \\
0 & 1 & \cdots & 0 \\
\vdots & \vdots & \ddots & \vdots \\
0 & 0 & \cdots & 1
\end{pmatrix}\text{.}
\]
现在依次把第一个因子的各行与第二个因子的各列相乘, 得到如下关系:
\[
\begin{array}{c c}
|\alpha_{11}|^2 = 1, & \text{特别地} \ \alpha_{11} \neq 0 \\[2pt]
\alpha_{11} \cdot \alpha_{i1} = 0, & \text{从而对} \ i = 2, \ldots, n \ \text{有} \ \alpha_{i1} = 0 \\[2pt]
|\alpha_{22}|^2 = 1, & \text{特别地} \ \alpha_{22} \neq 0 \\[2pt]
\alpha_{22} \cdot \alpha_{i2} = 0, & \text{从而对} \ i = 3, \ldots, n \ \text{有} \ \alpha_{i2} = 0 \\[2pt]
\text{等等.}
\end{array}
\]
由此我们看到 $A$ 确实是对角矩阵.相应的计算对酉矩阵 $A \in \mathrm{U}(n)$ 也成立.此时只需把 $A^t$
换成伴随矩阵 $A^*$.

\noindent\textbf{解题路径:} 上面提示的大部分计算可以避免, 如果我们更根本地进行的话.因此我们把
$V = \mathbb{R}^n$ 在\emph{典范数量积}(\textup{kanonisches Skalarprodukt})\index{dianfanshuliangji@典范数量积(kanonisches Skalarprodukt)}下看作欧几里得向量空间, 其中典范基 $e = (e_1, \ldots, e_n)$ 构成一个
标准正交基.现设 $A \in \mathrm{O}(n)$ 是以 $a_1, \ldots, a_n \in V$ 为列的正交矩阵, 即
$A = (a_1, \ldots, a_n)$.由于 $A \cdot A^t$ 给出单位矩阵, 我们认识到 $(a_1, \ldots, a_n)$ 也是 $V$ 的一个
标准正交基.我们再假定 $A$ 具有下三角形.则对 $i = 1, \ldots, n$ 有
$a_i \in \langle e_i, \ldots, e_n \rangle$, 并且我们要对 $n$ 作归纳, 证明必然推出对 $i = 1, \ldots, n$
有 $\langle a_i \rangle = \langle e_i \rangle$.

首先当然成立 $\langle a_n \rangle = \langle e_n \rangle$.这特别地完成了 $n = 1$ 时的归纳起点.进一步推出
\[
a_1, \ldots, a_{n-1} \in \langle a_n \rangle^{\perp} = \langle e_n \rangle^{\perp} = \langle e_1, \ldots, e_{n-1} \rangle,
\]
从而
\[
a_i \in \langle e_i, \ldots, e_n \rangle \cap \langle e_1, \ldots, e_{n-1} \rangle = \langle e_i, \ldots, e_{n-1} \rangle, \qquad i = 1, \ldots, n - 1\text{.}
\]
归纳假设因此给出对 $i = 1, \ldots, n - 1$ 有 $\langle a_i \rangle = \langle e_i \rangle$, 我们总体看到
$A = (a_1, \ldots, a_n)$ 是对角矩阵.

同样的论证方式对酉矩阵也适用.此时把 $V = \mathbb{C}^n$ 在典范数量积下看作酉向量空间, 以典范基
$\boldsymbol{e} = (e_1, \ldots, e_n)$ 为标准正交基.由于对 $A \in \mathrm{U}(n)$, $A \cdot A^*$ 给出
单位矩阵, $A$ 的列构成 $V$ 的另一个标准正交基.与欧几里得情形一样, 我们看到: 若 $A$ 具有下三角形,
则 $A$ 已经是对角矩阵.

\subsection*{习题 7.6/1}

\noindent\textbf{起点:} 给定一个有限维欧几里得或酉 $K$-向量空间及一个自伴随自同态
$\varphi \colon V \longrightarrow V$.

\noindent\textbf{目标:} 证明 $\varphi$ 恰好当对一切 $x \in V - \{0\}$ 成立
$\langle \varphi(x), x \rangle > 0$ 时, 只有正实特征值.

\noindent\textbf{策略:} 我们寻找一个简单途径来计算对 $x \in V$ 的 $\langle \varphi(x), x \rangle$.由于
$\varphi$ 自伴随, 推论~\mtref{cor:7.6.4} 的结果正适合.于是存在 $V$ 的一个标准正交基
$X = (x_1, \ldots, x_n)$, $\varphi$ 关于它由一个实对角矩阵描述, 例如
$A_{\varphi, X, X} = \mathrm{Diag}(\lambda_1, \ldots, \lambda_n)$.对带系数
$\alpha_1, \ldots, \alpha_n \in \mathbb{K}$ 的 $x = \sum_{j = 1}^{n} \alpha_j x_j \in V$, 我们可以如下计算:
\[
\langle \varphi(x), x \rangle = \left\langle \sum_{j = 1}^{n} \lambda_j \alpha_j x_j, \sum_{j = 1}^{n} \alpha_j x_j \right\rangle = \sum_{j = 1}^{n} \lambda_j |\alpha_j|^2\text{.}
\]
从这个关系中可以直接读出答案.

\noindent\textbf{解题路径:} 我们利用上面的等式.若 $\varphi$ 只有正特征值, 即对 $j = 1, \ldots, n$ 有
$\lambda_j > 0$, 则只要不是所有系数 $\alpha_i$ 都消失, 即 $x \neq 0$, 就推出
\[
\langle \varphi(x), x \rangle = \sum_{j = 1}^{n} \lambda_j |\alpha_j|^2 > 0\text{.}
\]
反之, 设对一切 $x \neq 0$, 即对不是全部消失的系数 $\alpha_j$, 成立
\[
\langle \varphi(x), x \rangle = \sum_{j = 1}^{n} \lambda_j |\alpha_j|^2 > 0\text{.}
\]
则对 $j = 1, \ldots, n$ 必然推出 $\lambda_j > 0$, 即 $\varphi$ 的一切特征值都是正的.

\subsection*{习题 7.6/3}

\noindent\textbf{起点:} 给定一个矩阵 $A \in \mathbb{R}^{n \times n}$.

\noindent\textbf{目标:} 需证明 $A$ 恰好当存在复矩阵 $S \in \mathbb{C}^{n \times n}$ 使得 $A = S^t \cdot S$
时是对称的.

\noindent\textbf{策略:} 若 $A$ 具有形式 $S^t \cdot S$(其中 $S \in \mathbb{C}^{n \times n}$), 则得
$A^t = (S^t \cdot S)^t = S^t \cdot S = A$, 于是 $A$ 对称.反之, 现设 $A$ 对称.则由关于主轴变换的
大定理~\mtref{bigthm:7.6.6}, 存在正交矩阵 $S \in \mathrm{O}(n)$, 使得 $D = S^{-1} \cdot A \cdot S$ 是一个
实对角矩阵.于是由 $S^{-1} = S^t$ 得 $A = S \cdot D \cdot S^t$, 其中如果我们把 $S$ 换成 $S^t$, 也可以
写成 $A = S^t \cdot D \cdot S$.若此时 $D$ 例如是单位矩阵, 则 $A$ 已通过 $A = S^t \cdot S$ 具有所求的形式.
否则我们不妨说从 $D$ 中提取平方根, 这在 $\mathbb{C}$ 上是可能的, 并把该平方根以适当方式与 $S$ 相乘.

\noindent\textbf{解题路径:} 等式 $A^t = (S^t \cdot S)^t = S^t \cdot S = A$ 表明, 只要 $A$ 具有形式
$A = S^t \cdot S$(其中 $S \in \mathbb{C}^{n \times n}$), $A$ 就是对称的.因此现设 $A$ 对称.由主轴变换
大定理~\mtref{bigthm:7.6.6}, 此时存在正交矩阵 $S \in \mathrm{O}(n)$, 使得 $S^t \cdot A \cdot S$ 给出一个实
对角矩阵 $D(\lambda_1, \ldots, \lambda_n)$.若再对每个 $\lambda_j$ 选取一个平方根 $\sqrt{\lambda_j}$, 这在
$\mathbb{C}$ 中无论如何都是可能的, 则得到
\[
\begin{array}{c}
A = S \cdot D(\lambda_1, \ldots, \lambda_n) \cdot S^t = S \cdot D\bigl(\sqrt{\lambda_1}, \ldots, \sqrt{\lambda_n}\bigr)^2 \cdot S^t \\[2pt]
= \Bigl(D\bigl(\sqrt{\lambda_1}, \ldots, \sqrt{\lambda_n}\bigr) \cdot S^t\Bigr)^t \cdot \Bigl(D\bigl(\sqrt{\lambda_1}, \ldots, \sqrt{\lambda_n}\bigr) \cdot S^t\Bigr),
\end{array}
\]
从而得到 $A$ 的所求类型的分解.
